% \documentclass[12pt]{article}
\documentclass[12pt]{article}
\usepackage{setspace}
\usepackage[labelfont=bf]{caption}
\usepackage[margin=1in]{geometry}
\usepackage{amsmath,amssymb,amsthm,bm,mathtools}
\usepackage{algorithm}
\usepackage{algpseudocode}
\usepackage{enumitem}
\usepackage{hyperref}

\usepackage{graphicx} % Required for inserting images
\usepackage{hyperref}
\usepackage{listings}
\usepackage{booktabs}

\usepackage{xcolor} % Required for colored text

\definecolor{impGreen}{rgb}{0.0, 0.5, 0.0}
\definecolor{worsRed}{rgb}{0.7, 0.0, 0.0}

\usepackage[style=numeric, sorting=nyt]{biblatex}
\addbibresource{references.bib}

\usepackage{amsmath}
\usepackage{amssymb}

\usepackage{graphicx}
\usepackage{subcaption}

\usepackage{tcolorbox}
\usepackage{xcolor} % Required for defining custom colors
\definecolor{lightgray}{gray}{0.9} % Define a light gray color

\usepackage{color,soul}

\setlength{\marginparwidth}{2cm}
\usepackage[textsize=tiny]{todonotes}
\usepackage{dsfont}

\usepackage{lineno}
\renewcommand{\linenumberfont}{\normalfont\bfseries\tiny\color{gray}}
\linenumbers

\usepackage{multicol}

\usepackage{makecell}
\usepackage{xcolor}

\newtheorem{lemma}{Lemma}
\newtheorem{corollary}{Corollary}

\hypersetup{colorlinks=true,linkcolor=blue,citecolor=blue,urlcolor=blue}

% \title{\vspace{-5mm}Reducing Regressivity in ML-Based Mass Appraisal:\\ Fairness-Aware Learning for Valuation Workflows\vspace{0mm}}
% \author{Nicolás Acevedo Villena\\
% {\small Operations Research Center, MIT, Cambridge, MA 02139. \texttt{nacevedo@mit.edu}}\\
% }
% \date{
% \vspace{-8mm}
% }

\title{\vspace{-5mm}Reducing Regressivity in ML-Based Mass Appraisal:\\ A Covariance-Penalty Layer for Assessor Workflows\vspace{0mm}}

\author{
    Nicolas Acevedo\\
    {\small Operations Research Center, MIT, Cambridge, MA 02139. \texttt{nacevedo@mit.edu}}
    \and
    Haihao Lu\\
    {\small Sloan School of Management, MIT, Cambridge, MA 02139. \texttt{haihao@mit.edu}}
}

\date{
\vspace{-8mm}
}
\theoremstyle{plain}
\newtheorem{theorem}{Theorem}
\newtheorem{proposition}{Proposition}
\theoremstyle{definition}
\newtheorem{definition}{Definition}
\newtheorem{assumption}{Assumption}
\newtheorem{remark}{Remark}

\newcommand{\R}{\mathbb{R}}
\newcommand{\E}{\mathbb{E}}
\newcommand{\prox}{\mathrm{prox}}
\newcommand{\argmin}{\mathop{\mathrm{arg\,min}}}
\newcommand{\norm}[1]{\left\lVert #1 \right\rVert}
\newcommand{\ip}[2]{\left\langle #1,#2 \right\rangle}
\newcommand{\1}{\mathbf{1}}
\newcommand{\Cov}{\text{Cov}}
\newcommand{\Corr}{\text{Corr}}
\newcommand{\Var}{\text{Var}}

\newcommand{\Y}{\mathcal{Y}}
\newcommand{\X}{\mathcal{X}}
\newcommand{\V}{\mathcal{V}}
\newcommand{\U}{\mathcal{U}}
\newcommand{\G}{\mathcal{G}}
\newcommand{\D}{\mathcal{D}}
\newcommand{\N}{\mathcal{N}}
\newcommand{\F}{\mathcal{F}}
\newcommand{\A}{\mathcal{A}}
\newcommand{\J}{\mathcal{J}}
\newcommand{\I}{\mathcal{I}}
\newcommand{\M}{\mathcal{M}}
\newcommand{\C}{\mathcal{C}}
\renewcommand{\L}{\mathcal{L}}
\renewcommand{\H}{\mathcal{H}}
\renewcommand{\P}{\mathcal{P}}
\renewcommand{\R}{\mathbb{R}}

\newcommand{\diag}{\mathbf{\operatorname{diag}}}

\newcommand{\muBold}{\boldsymbol{\mu}}
\renewcommand{\t}{^{\top}}
\newcommand{\qt}[1]{``#1''}

\newcommand{\Expected}[1]{\mathbb{E}[#1]}
\newcommand{\Prob}[1]{\mathbb{P}(#1)}

\newcommand{\red}[1]{{\color{red} #1}}

% \doublespacing
\begin{document}
\maketitle

\begin{abstract}

Property taxes fund a large share of U.S.\ local government services, and the values that determine them are increasingly produced by machine learning (ML) models. These models are usually optimized for predictive accuracy, which can leave \emph{regressivity} in place or even amplify it: lower-value properties are over-assessed while higher-value properties are under-assessed. We introduce a covariance-based regressivity-control layer for standard valuation learners. Starting from the observation that regressivity is, in assessor practice, a value-dependent pattern in assessment ratios, we add a soft penalty on the empirical covariance between log-residuals and log-prices. We develop two forms of this penalty: a direct squared-covariance penalty and an implementation-friendly, sample-additive surrogate that plugs into standard gradient-boosting interfaces through the usual custom gradient and Hessian routines. We apply the method inside a translated Cook County Assessor's Office (CCAO) residential pipeline with 349{,}661 sales and 95 features. A single penalty hyperparameter recovers most of the vertical-equity gap: for a representative surrogate configuration, PRD moves from $1.085$ to $1.038$, PRB from $-0.118$ to $-0.011$, and VEI from $-36.8\%$ to $-8.5\%$, while $R^2$ stays at $0.882$ and COD and COV improve slightly relative to the unconstrained LightGBM baseline. We also cast final model selection as a constrained calibration problem, selecting the model that maximizes predictive performance subject to explicit ratio-study tolerances. The method adds one hyperparameter to existing assessor workflows, keeps the same preprocessing, feature set, and diagnostics, and is designed as a targeted correction layer rather than a replacement for current valuation systems.

\end{abstract}

\section{Introduction}

Property taxes are one of the largest sources of local public revenue in the United States, providing roughly 72\% of local tax revenue and 47\% of local own-source general revenue, or about \$500 billion per year \cite{Berry2021}. The values that generate those taxes are produced by \emph{mass appraisal}: the estimation, as of a given date, of a large set of real estate values using common data, standardized methods, and statistical testing \cite{IAAO2025Mass}. Assessor offices have increasingly turned to machine learning (ML) models such as gradient-boosted trees to carry out this task, motivated by their predictive performance on tabular data and their scalability to large parcel inventories \cite{IAAOAITaskForce2022,BidansetRakow2022,Rootetal2023,Sevgen2024100111,Zilli2024}. In this setting, however, accuracy alone is insufficient. Assessments must also satisfy the operational constraints that assessors evaluate through standardized \emph{ratio studies}: equity across similar properties (horizontal equity) and equity across value strata (vertical equity) \cite{Quintos2014,10578d49-78f8-3ef9-8277-9b6add2bd071,IAAO2013Ratio,IAAO2025ExposureRatio}.

Vertical equity is typically the hardest constraint to satisfy in deployed systems. Assessments are \emph{regressive} when lower-value properties are over-assessed and higher-value properties are under-assessed; such patterns are widespread across U.S.\ counties and carry substantial distributional consequences for taxpayers \cite{McMillen2020,Berry2021,IAAO2023VerticalReview,McMillenSingh2023}. \textcite{amornsiripanitch2020why} estimates that owners of lower-valued properties face effective tax rates up to 50\% higher than owners of higher-valued properties. Adopting a more flexible learner does not, by itself, resolve the problem. Squared-error objectives target conditional means, so when some price variation remains unexplained by the available features, the fit induces a systematic non-positive dependence between residuals and log-price; we make this precise in Section~\ref{subsec:squared_loss_regressivity}. The empirical counterpart is sharp in our CCAO study: moving from linear regression to a LightGBM baseline improves $R^2$, MAE, COD, and COV, yet worsens PRD, PRB, and VEI, pushing the model out of the assessor compliance region (Table~\ref{tab:baseline_models}).

\paragraph{This paper.}
We develop a regressivity-control layer that sits on top of the standard valuation learners already used by assessors. Starting from the sample covariance $\widehat{\Cov}(e,y)$ between log-residuals $e=\hat y-y$ and log-prices $y=\log P$, we add a soft penalty that discourages value-dependent trends in the fitted assessment ratios. We study two forms of this penalty:
\begin{itemize}
    \item a \emph{direct} squared-covariance penalty, which is the closest instantiation of the dependence-control idea;
    \item an \emph{implementation-friendly} surrogate obtained as a separable upper bound on the squared covariance, which is sample-additive and therefore compatible with standard gradient-boosting interfaces such as LightGBM and XGBoost.
\end{itemize}
The direct penalty introduces dense curvature and is handled either through a diagonal Hessian approximation or through a custom leaf solver. The surrogate is separable by construction and is the variant we recommend for routine deployment. Both penalties add one scalar hyperparameter $\rho\ge 0$ and recover the unpenalized learner at $\rho=0$.

We evaluate the method in a translated version of the Cook County Assessor's Office (CCAO) residential pipeline with 349{,}661 verified sales over 2016--2023, 95 features, and an expanding-window rolling-origin validation protocol. Relative to the LightGBM baseline, a single choice of $\rho$ moves PRD from $1.085$ to $1.038$, PRB from $-0.118$ to $-0.011$, and VEI from $-36.8\%$ to $-8.5\%$, with $R^2$ essentially unchanged at $0.882$ and small improvements in COD and COV (Table~\ref{tab:penalized_models}). Because the penalty does not directly optimize ratio-study statistics, we also cast final model selection as a constrained calibration problem: choose the candidate that maximizes predictive performance subject to explicit ratio-study tolerances. Figure~\ref{fig:regressivity_motivation} shows the ratio--price pattern that the method is designed to flatten.

\begin{figure}[!ht]
    \centering
    \includegraphics[width=0.7\linewidth]{Fairness Price/CCAO meeting/img/Baseline LightGBM_0_motivation.pdf}
    \caption{Regressivity in the baseline LightGBM fit on the 2023 Cook County test set. The assessment-to-sale ratio $r=\widehat P/P$ declines with property value: lower-value properties tend to be over-assessed ($r>1$) and higher-value properties tend to be under-assessed ($r<1$).}
    \label{fig:regressivity_motivation}
\end{figure}

\paragraph{Contributions.}
\begin{itemize}
    \item \textit{A covariance-based regressivity-control layer.} We formalize regressivity as a value-dependent pattern in assessment ratios and build a soft covariance penalty that can be added to standard regression losses. We develop both a direct squared-covariance penalty and an implementation-friendly separable surrogate. The surrogate is motivated by an upper bound on the squared covariance; dropping residual centering produces a sample-additive approximation that preserves the same value-weighted intuition and plugs into standard boosting APIs.
    \item \textit{Implementation in standard learners.} We instantiate the penalties in regularized linear models and in second-order gradient-boosted trees. For the surrogate, we provide the exact gradient and Hessian that plug into LightGBM and XGBoost custom-objective interfaces. For the direct penalty, we provide both a diagonal approximation compatible with those interfaces and an exact Newton-style leaf solver (Appendix~\ref{app:boosting_algorithms}).
    \item \textit{Evidence in a real CCAO pipeline with constrained calibration.} Using the CCAO residential workflow, we show that a single penalty hyperparameter recovers most of the vertical-equity gap while preserving accuracy, and we cast final model selection as a constrained calibration problem that ranks candidates by predictive performance subject to explicit ratio-study tolerances.
\end{itemize}

The paper is organized as follows. Section~\ref{sec:problem_setting} fixes the assessment setting and the ratio-study diagnostics used throughout. Section~\ref{sec:covariance_penalty_learning_in_learners} develops the covariance-based penalties and their implementation in linear models and boosted trees. Section~\ref{sec:empirical} reports the CCAO evidence and the constrained calibration results. Section~\ref{sec:discussion} concludes.

\subsection{Related Work}

The paper sits at the intersection of three literatures: machine learning for mass appraisal, equity and regressivity in assessor practice, and fair regression with continuous attributes. We cover each briefly and in the depth required to position the contribution.

\paragraph{ML in mass appraisal and assessor workflows.}
Recent reviews document the move from regularized linear models to tree ensembles and, more recently, neural architectures, with gradient-boosted trees now among the most common production learners \cite{WangLi2019,Rootetal2023,Sevgen2024100111,Zilli2024,purdueAssessingTechniques}. Evaluation in this literature is still dominated by predictive criteria rather than equity as an explicit design target. Adoption in assessor offices is partial and uneven: \textcite{BidansetRakow2022} identify communication, technical capacity, and budget as the primary barriers, and the \textcite{IAAOAITaskForce2022} lists explainability, defensibility, administrative feasibility, and workflow integration as the conditions for real-world use. Any practical equity intervention therefore has to be compatible with these learners and layered onto the existing training pipeline.

\paragraph{Equity and regressivity in mass appraisal.}
Assessor standards decompose equity into horizontal equity, measured by the coefficient of dispersion (COD), and vertical equity, measured by the price-related differential (PRD), the coefficient of price-related bias (PRB), and the recently introduced vertical equity indicator (VEI) \cite{IAAO2013Ratio,IAAO2025Mass,IAAO2025ExposureRatio,10578d49-78f8-3ef9-8277-9b6add2bd071,cornia2006property,McMillenSingh2023}. \textcite{IAAO2023VerticalReview} recommend reporting a suite of these tests rather than relying on any single statistic, and \textcite{McMillenSingh2023} show that some regression-based regressivity detectors can be severely biased when applied to ML-generated assessments. \textcite{Berry2021} documents that regressivity is widespread across U.S.\ counties and more consistent with assessors' data and model limitations than with pure measurement error, and \textcite{AvenancioLeonHoward2022} document related tax-incidence effects. Corrections inside the valuation pipeline have been proposed via two-step spatial procedures \cite{Quintos2014}, geographically weighted regression \cite{BidansetEtAl2019,BidansetEtAl2025}, validation-time ratio studies \cite{HermansEtAl2022,HermansEtAl2025}, and fairness-aware valuation models \cite{CandoganHanLu2023}.

\paragraph{Fair regression and dependence-based regularization.}
Most early fair-ML work addresses classification, whereas property assessment is a continuous-prediction problem \cite{CatonHaas2024,BarocasHardtNarayanan2023}. Fair-regression methods treat fairness as an explicit objective through constrained empirical risk minimization \cite{AgarwalDudikWu2019,FitzsimonsEtAl2019,WeiEtAl2023}, post-processing via optimal transport \cite{ChzhenEtAl2020Wasserstein,ChzhenEtAl2020Recalibration}, or explicit dependence constraints \cite{KomiyamaEtAl2018,ScutariPaneroProissl2022}. The closest strand for our purposes is dependence-based regularization \cite{MaryEtAl2019,PerezSuayEtAl2017,LiEtAl2022Fairness,lee2022maximal,BuylDefranceDeBie2024FAIRRET}, which penalizes statistical dependence between model outputs and a continuous sensitive quantity. These regularizers require small changes to the base learner and are natural when the relevant inequity runs along a continuous variable rather than a demographic attribute, as is the case for assessment ratios and sale prices. For multi-objective selection we build on \textcite{martinez2020minimax} and \textcite{little2022fairness}.

The assessor literature defines regressivity as a value-dependent pattern in assessment ratios and addresses it mainly through diagnostics or post hoc corrections; the fair-ML literature provides training-time machinery to control such dependence but typically for demographic attributes rather than assessor-style vertical inequity. Our contribution is to recast the assessor-defined pattern as a covariance target on a standard valuation learner and to instantiate it inside the second-order boosting interfaces that assessor offices already use.


\section{The Assessment Setting and the CCAO Workflow}
\label{sec:problem_setting}

This section fixes the assessor-facing ratio objects, the predictive notation, and the ratio-study diagnostics used throughout the paper. It also recasts regressivity as a covariance-based dependence pattern that links classical ratio-study diagnostics to the covariance penalties introduced in Section~\ref{sec:covariance_penalty_learning_in_learners}. Throughout, we anchor the exposition in the operational workflow of the Cook County Assessor's Office (CCAO), which also provides the empirical setting of Section~\ref{sec:empirical}.

\paragraph{The CCAO residential workflow.}
The CCAO produces residential valuations for roughly 1.5 million parcels on a triennial cycle using a LightGBM-based automated valuation model (AVM), and publishes the pipeline, feature set, and ratio-study diagnostics as open code. Sales are verified upstream and merged with 95 property, location, and neighborhood features; the learner is trained on verified arm's-length sales and evaluated on held-out years through ratio studies reported on the original price scale. Final candidate models are promoted only if they meet the office's internal ratio-study tolerances. Throughout the paper, we take this operational loop---train on historical verified sales, evaluate with price-scale ratios, screen with IAAO-style tolerances---as the deployment context that any regressivity correction must fit into. Section~\ref{sec:empirical} describes our translation of this workflow into a Python-based implementation, the 2016--2022 training window, the 2023 held-out test year, and the 8-fold rolling-origin validation protocol.

\subsection{Assessment Ratios and Equity Measures}
\label{subsec:equity_metrics}

For a fitted model and an evaluation sample $S$, let $\widehat P_i$ denote the predicted value of property $i$ and $P_i$ its observed sale price. The \emph{assessment-to-sale ratio} is
\begin{equation}
\label{eq:ratio_def}
r_i := \frac{\widehat P_i}{P_i},
\qquad \forall i\in S,
\end{equation}
where dependence on the fitted model is suppressed when clear from context. We refer to $r_i$ as the \emph{assessment ratio}.

\emph{Ratio studies} summarize assessment quality through statistics of $\{r_i\}_{i\in S}$ \cite{IAAO2025ExposureRatio}. We use the standard IAAO equity diagnostics and the central statistics that measure the overall \emph{assessment level}.

\paragraph{Remark.} All equity diagnostics in this paper are computed on the price scale of $\widehat P_i$ and $P_i$, even though the underlying predictive models specified in later sections are trained on the log-prices scale.

\paragraph{Assessment level.}
The \emph{assessment level} is summarized by central statistics of $\{r_i\}_{i\in S}$: the median, the mean, the value-weighted mean, and the standard deviation \cite{IAAO2025ExposureRatio}. Let
\[
m_S := \operatorname{med}\{r_i: i\in S\}
\]
denote the median ratio,
\[
\bar r (S) := \frac{1}{|S|}\sum_{i\in S} r_i
\]
the mean ratio,
\[
\overline{r_W} (S) := \frac{1}{\sum_{i\in S} P_i}\sum_{i\in S} r_i P_i
\]
the value-weighted mean ratio, and
\[
\sigma_S := \sqrt{\frac{1}{|S|-1}\sum_{i\in S}(r_i-\bar r)^2}
\]
the sample standard deviation of the assessment ratios. We omit the dependence on $S$ when the context is clear.

\medskip
\noindent\textbf{Horizontal equity measures.}
We summarize \textit{horizontal equity} primarily through the coefficient of dispersion (COD), which the IAAO describes as the average percentage deviation from the median ratio $m_S$ \cite{IAAO2013Ratio,IAAO2025ExposureRatio}:
\begin{equation}
\label{eq:cod}
\operatorname{COD}(S)
:=
100\cdot\frac{1}{|S|}
\sum_{i\in S}\frac{|r_i-m_S|}{m_S}.
\end{equation}
Lower COD values indicate greater within-group uniformity. We also report the coefficient of variation (COV), defined as
\begin{equation}
\label{eq:cov_metric}
\operatorname{COV}(S)
:=
100\cdot\frac{\sigma_S}{\bar r},
\end{equation}
as a supplementary dispersion diagnostic. Both indicators summarize how dispersed the ratios are around a central tendency.

\medskip
\noindent\textbf{Vertical equity measures.}
We summarize vertical equity using standard \textit{price-related diagnostics}. First, the price-related differential (PRD) is
\begin{equation}
\label{eq:prd}
\operatorname{PRD}(S)
:=
\frac{\bar r}{\overline{r_W}}.
\end{equation}
Under standard assessor interpretation, PRD values above $1.03$ tend to indicate regressivity, while values below $0.98$ tend to indicate progressivity \cite{IAAO2013Ratio,IAAO2025ExposureRatio}.

We also report the coefficient of price-related bias (PRB), defined in terms of a quantity referred to as the \textit{market-value proxy}; \cite{IAAO2025ExposureRatio}
$
\operatorname{MV}_i
:=
\log_2\!\left(\frac{1}{2}\left(P_i+\frac{\widehat P_i}{m_S}\right)\right),
$
designed to reduce measurement-error bias \cite{IAAO2025ExposureRatio}. The PRB is then defined as
\begin{equation}
\label{eq:prb}
\operatorname{PRB}(S)
:=
\operatorname{slope}\!\left(\frac{r_i-m_S}{m_S}\sim \operatorname{MV}_i\right).
\end{equation}
That is, PRB is the slope of relative ratio deviations from the median against the market-value proxy. PRB values below $-0.05$ suggest regressivity, while values above $0.05$ suggest progressivity \cite{IAAO2013Ratio,IAAO2025ExposureRatio}.

Finally, we report the vertical equity indicator (VEI), recently introduced in the IAAO 2025 exposure draft \cite{IAAO2025ExposureRatio}, defined as
\begin{equation}
\label{eq:vei}
\operatorname{VEI}(S)
:=
100\cdot\frac{m_{S_{q_{10}}}-m_{S_{q_1}}}{m_S},
\end{equation}
where $m_{S_{q_{10}}}$ is the median ratio in the highest value decile of $S$ (ranked by the proxy $\operatorname{MV}_i$) and $m_{S_{q_1}}$ is the median ratio in the lowest value decile. VEI values below $-10\%$ suggest regressivity, while values above $10\%$ suggest progressivity \cite{IAAO2025ExposureRatio}.

Table \ref{tab:assessment_metrics_summary} summarizes the main assessor-facing diagnostics used later for empirical evaluation and constrained model selection.

\begin{table}[!ht]
\centering
\caption{Summary of core assessment ratio statistics, desired targets, and interpretation.}
\label{tab:assessment_metrics_summary}
\begin{tabular}{lcp{8.5cm}}
\toprule
\textbf{Measure} & \textbf{Target Value} & \textbf{Interpretation} \\
\midrule 
\multicolumn{3}{l}{\textit{Central Statistics (Assessment Level)}} \vspace{3mm} \\
Median ($m_S$) & \(\approx 1\) & Outlier-robust measure of overall assessment level. \\
Mean ($\bar r$) & \(\approx 1\) & Arithmetic average. \\\\
Weighted Mean ($\overline{r_W}$) & \(\approx 1\) & Value-weighted average; denominator in the PRD. \\
\midrule
\multicolumn{3}{l}{\textit{Horizontal Equity (Uniformity)}} \vspace{3mm} \\
COD & \([5.0,15.0]\textsuperscript{*}\) & Average \% deviation from the median. Lower values indicate better uniformity. \\
COV & $-$ & Ratio of standard deviation to the mean. Lower values indicate less relative variability. \\
\midrule
\multicolumn{3}{l}{\textit{Vertical Equity (Regressivity)}} \vspace{3mm} \\
PRD & \([0.98,1.03]\) & Values \(>1.03\) indicate regressivity. Values \(<.098\) indicate progressivity. \\
PRB & \([-0.05,0.05]\) & Values \(<-0.05\) indicate regressivity. Values \(>0.05\) indicate progressivity. \\
VEI & \([-10\%,10\%]\) & Values \(<0\) indicate regressivity. Values \(>0\) indicate progressivity.\\
\bottomrule
\multicolumn{3}{l}{\footnotesize \textsuperscript{*}Assessments with COD values below 5.0 are considered highly suspect. Acceptable upper limits} \\
\multicolumn{3}{l}{\footnotesize {\color{white}\textsuperscript{*}}vary by property type (up to 20.0 or 25.0 for rural/vacant) \cite{IAAO2025ExposureRatio}.}
\end{tabular}
\end{table}

\subsection{Regressivity as a Dependence Problem}
\label{subsec:dependence_problem}

The assessor literature describes regressivity as a price-related pattern in ratios: assessment ratios tend to fall as property value rises. Figure~\ref{fig:regressivity_motivation_boxed} shows this pattern as a negative trend between the ratios $\{r_i\}_{i\in S}$ and log-prices $\{y_i\}_{i\in S}$ with $y_i=\log P_i$. We plot against log-price because sale prices are highly right-skewed; log-price makes the trend comparable across value levels.

\begin{figure}[!ht]
    \centering
    \includegraphics[width=0.7\linewidth]{Fairness Price/GE Paper/img/motivation/Baseline LightGBM_0_motivation_squared_only.pdf}
    \caption{Illustration of regressive behavior in mass appraisal. Assessment-to-sale ratios $r$ tend to decrease with property value $P$: lower-value properties are more likely to be over-assessed ($r > 1$), while higher-value properties are more likely to be under-assessed ($r < 1$).}
    \label{fig:regressivity_motivation_boxed}
\end{figure}

Because ratios and prices are continuous quantities, a natural first-order summary of their dependence is the empirical covariance:
\begin{equation}
\label{eq:sample_cov}
\operatorname{Cov}(r,P)
:=
\frac{1}{|S|}
\sum_{i\in S}
(r_i-\bar r)(P_i-\bar P),
\end{equation}
where $\bar r$ and $\bar P$ are the sample means of $\{r_i\}_{i\in S}$ and $\{P_i\}_{i\in S}$. When $\Var(r)>0$ and $\Var(P)>0$, the covariance, correlation, and simple-regression slope share a sign:
\begin{equation}
\label{eq:cov_corr_slope_sign}
\operatorname{sign}\!\left(\Cov(r,P)\right)
=
\operatorname{sign}\!\left(\Corr(r,P)\right)
=
\operatorname{sign}\!\left(\operatorname{slope}(r \sim P)\right),
\end{equation}
where
$
\Corr(r,P):=\frac{\Cov(r,P)}{\sqrt{\Var(r)\Var(P)}}$, $
\operatorname{slope}(r \sim P):=\frac{\Cov(r,P)}{\Var(P)}$, and $\operatorname{sign}(\cdot) \in \{-1,0,1\}$ is the signum function.
Thus $\Cov(r,P)<0$ captures a first-order price-related negative pattern, which we associate with regressivity. The covariance is also directionally aligned with the standard price-related differential (PRD) used in assessor practice \cite{IAAO2025ExposureRatio}, as formalized in Proposition~\ref{prop:cov_prd}.

\begin{proposition}
\label{prop:cov_prd}
Let $P>0$ and $r=\widehat P/P>0$ almost surely, with $\Expected{r}<\infty$, $\Expected{P}<\infty$, and $\Expected{rP}<\infty$. Then the population covariance  and the population
$
\mathrm{PRD}=\frac{\Expected{r}\Expected{P}}{\Expected{rP}}
$
satisfy
\begin{equation}
\label{eq:cov_prd_sign}
\operatorname{sign}\!\left(\Cov(r,P)\right)
=
\operatorname{sign}\!\left(1-\mathrm{PRD}\right).
\end{equation}
\end{proposition}

\begin{proof}
Since
$
\Cov(r,P)=\Expected{rP}-\Expected{r}\Expected{P},
$
we have
\[
1-\mathrm{PRD}
=
1-\frac{\Expected{r}\Expected{P}}{\Expected{rP}}
=
\frac{\Expected{rP}-\Expected{r}\Expected{P}}{\Expected{rP}}
=
\frac{\Cov(r,P)}{\Expected{rP}}.
\]
Because $r>0$ and $P>0$ almost surely, $\Expected{rP}>0$. Therefore, $1-\mathrm{PRD}$ and $\Cov(r,P)$ have the same sign.
\end{proof}

Proposition~\ref{prop:cov_prd} shows that the covariance $\Cov(r,P)$ is directionally consistent with the PRD, which links the ratio-study view of regressivity to the dependence-based formulation developed in the next sections.

\subsection{Mass-Appraisal Prediction Setting}
\label{subsec:prediction_setting}

% Mass appraisal estimates the value of a \emph{group} of properties as of a given date using common data, standardized methods, and statistical testing \cite{IAAO2025Mass}. 

We observe a sample of $n$ sold properties
$
\D_n := \{(x_i,P_i,t_i)\}_{i=1}^n,
$
where $x_i\in\mathcal X\subseteq\R^p$ is a feature vector (e.g., square footage, bedroom count, distance to transit), $P_i\in\R_{++}$ is the verified sale price, and $t_i$ is the sale date. Because prices are positive and right-skewed, we work on the log-price scale, $y_i := \log P_i$.

Let $\Theta$ be a parameter space and $f_\theta:\mathcal X\to\R$, $\theta\in\Theta$, a candidate predictor of log-price (e.g., $f_\theta(x_i)=\theta^\top x_i$ in linear models). Each $i\in[n]$ receives a log-scale prediction $\hat y_i := f_\theta(x_i)$, with corresponding price-scale prediction $\widehat P_i := \exp(\hat y_i)$.

% \paragraph{Remark.} Assessor-facing diagnostics in Section \ref{subsec:equity_metrics} are computed on the original sale-price scale; using $P_i$ and $\widehat P_i=\exp(\hat y_i)$. This is deliberate: the workflow is evaluated through the price-scale assessment ratios, not through recovery of a bias-corrected conditional mean on the original price scale. Therefore, models are trained on the log-price scale and evaluated on the price scale.

\paragraph{The baseline tension.}
The empirical problem this paper addresses can be stated already at this level. The CCAO workflow moved from linear regression to a LightGBM AVM to exploit the predictive gains of tree ensembles on tabular real-estate data. On the translated 2023 CCAO test set (Section~\ref{sec:empirical}, Table~\ref{tab:baseline_models}), LightGBM lifts $R^2$ from $0.762$ to $0.883$ and cuts MAE from \$86{,}257 to \$72{,}064 relative to linear regression, and it also improves horizontal-equity diagnostics (COD from $25.0\%$ to $23.1\%$, COV from $36.8\%$ to $35.4\%$). Over the same comparison, the vertical-equity diagnostics move in the opposite direction: PRD rises from $1.039$ to $1.085$, PRB falls from $-0.011$ to $-0.118$, and VEI drops from $-13.5\%$ to $-36.8\%$. The median ratio slides from $1.001$ to $0.944$. In operational terms, the stronger predictor exits the assessor compliance region on exactly the dimension the office most cares about. The covariance penalty introduced next is designed to recover vertical-equity behavior without giving up this predictive ground, and the remainder of the paper is organized around that single tension.

\paragraph{Log-scale residuals as a ratio proxy.} Models are trained on the log-price scale, so we define the \emph{log-scale residual}
$
e_i := \hat y_i - y_i, \quad i\in[n].
$
By construction, $e_i$ is the log-assessment ratio:
\begin{equation}
\label{eq:log_ratio_identity}
e_i
=
\hat y_i-y_i
=
\log \widehat P_i-\log P_i
=
\log\!\left(\frac{\widehat P_i}{P_i}\right)
=
\log r_i,
\qquad \forall i\in S.
\end{equation}
Because $\operatorname{sign}(e_i)=\operatorname{sign}(r_i-1)$, the sign of $e$ preserves the interpretation of over- and under-assessment. Identity \eqref{eq:log_ratio_identity} also yields the log-scale counterpart of the price-related dependence of Section~\ref{subsec:dependence_problem}: $\Cov(e,y)=\Cov(\log r,\log P)$. This is the quantity the method of Section~\ref{sec:covariance_penalty_learning_in_learners} penalizes during training.


% Thus, the log-scale residuals play a key role in the methodological sections below. 

% This identity links price-scale ratios to log-scale residuals and will be used in the methodological sections below.

\section{Covariance-Based Penalization in Standard Learners}
\label{sec:covariance_penalty_learning_in_learners}

% \section{General Regressivity-Aware Learning Setting}


% This section introduces the general structure of the proposed solution. We first state a baseline empirical learning problem that covers the main predictive model classes used in assessors pipelines. We then characterize why general squared-loss regression problems can induce negative linear dependence between fitted assessments residuals and log-prices. Fi

% This motivates a general constrained and penalized correction class, which serves as the abstract template for the covariance-based penalizations introduced in the next section.

% This section formalizes the in-processing penalization approach required for the regressivity-aware learning problem. We formalize the baseline learner in a general form that summarizes the common machine learning approaches used by assessors. We briefly discussed why the general learning problem can induce some regressive dependencies between ratios and sale prices. Finally, we introduce the abstract constrained and penalized correction class for the underlying learning problem, which serves as the formal introduction of the exact regressivity dependence penalization terms introduced in the next section. 



\subsection{General Baseline Learning Setting}

Let
$
\H := \{f_\theta : \theta \in \Theta\}
$
denote the hypothesis class induced by a parameter space $\Theta$, and let $\D_n$ denote a training sample of size $n$ with $y_i=\log P_i$. The empirical loss is
\[
L(f, \D_n)
:=
\frac{1}{n}\sum_{i=1}^n \ell(f(x_i),y_i),
\]
where $\ell$ is a prediction-error loss on the log-price scale. The canonical choice in assessor workflows is the squared residual loss, $\ell(f(x_i),y_i)=(f(x_i)-y_i)^2=e_i^2$, which we use throughout; the general notation is kept so that the approach extends to alternative losses.

\paragraph{Baseline learner.} The \emph{baseline} learner solves the empirical risk minimization problem
% \vspace{-4mm}
\begin{align}
\label{opt:base_model}
f_{\mathrm{base}}
\in
\arg\min_{f \in \H}
\left\{
L(f, \D_n)+\Omega(f)
\right\},
\end{align}
where $\Omega:\H\to\mathbb R$ is an optional complexity regularizer on $\Theta$ (e.g., $\ell_2$ regularization in ridge regression). We refer to $f_{\mathrm{base}}$ as the baseline predictor and to \eqref{opt:base_model} as the baseline learning problem.

% Formulation \eqref{opt:base_model} is intentionally broad: it covers the main learners studied later in the paper, including regularized linear regression and boosted-trees models.

% \hrule

\subsubsection{Why regressive patterns are expected in squared-loss learning}
\label{subsec:squared_loss_regressivity}

Squared-loss learning connects directly to the dependence view of regressivity developed in Section~\ref{subsec:dependence_problem}. In that setting the baseline predictor $f_{\mathrm{base}}$ targets the conditional mean $\E[y\mid x]$, and any residual conditional variation in log-price that cannot be explained from features generates a systematic non-positive linear dependence between log-residuals and log-prices. Proposition~\ref{prop:population_covariance_identity} makes this precise.


% We next explain why squared-loss learning can create regressive pressure even before any explicit fairness criterion is imposed. Because the squared-loss Bayes predictor targets conditional mean behavior, unresolved variation around that conditional mean can translate into systematic dependence between residuals and value level. The following proposition makes this precise.

% We further motivate the necessity of regressivity correction in general learning problems, such as $\eqref{opt:base_model}$; now from a theoretical perspective. Error-based learning problems naturally induce regressive patterns when some conditional price variation remains unexplained given the available features (as expected in practice). Intuitively, machine learning models aim to learn central tendency behaviors, hence pushing extreme values to some central measure of tendency. This means that extreme lower values are pushed towards the over-prediction, and the opposite for extreme higher values. We formalize this intuition for one of the most common risk measure in ML-based mass appraisal; the squared-error loss.  


\begin{proposition}[Population covariance non-positivity under squared-error prediction]
\label{prop:population_covariance_identity}

Let $Y:=\log P$ be the log-sale price, let $X$ denote the observed feature vector, and assume $\E[Y^2]<\infty$ and $\E[P^2]<\infty$. Let
$
f^\star(X)=\E[Y\mid X]
$
be the Bayes predictor under squared-error loss. Define the associated population residual
$
e^\star:=f^\star(X)-Y.
$

Let $g:\mathbb{R}\to\mathbb{R}$ be a measurable nondecreasing function such that $\E[g(Y)^2]<\infty$. Then
\begin{equation}
\label{eq:covariance_population_identity}
\Cov(e^\star,g(Y))
=
-\E\!\big[\Cov(Y,g(Y)\mid X)\big]
\le 0.
\end{equation}
\end{proposition}
\begin{proof}
By the law of total covariance,
\[
\Cov(e^\star,g(Y))
=
\E\!\big[\Cov(e^\star,g(Y)\mid X)\big]
+
\Cov\!\big(\E[e^\star\mid X],\,\E[g(Y)\mid X]\big).
\]
Since
\[
\E[e^\star\mid X]=\E[\E[Y\mid X]-Y\mid X]=0,
\]
the second term is zero. Moreover, conditional on $X$, the quantity $\E[Y\mid X]$ is constant, so
\[
\Cov(e^\star,g(Y)\mid X)
=
\Cov(\E[Y\mid X]-Y,\;g(Y)\mid X)
=
-\Cov(Y,g(Y)\mid X).
\]
Hence
\[
\Cov(e^\star,g(Y))
=
-\E\!\big[\Cov(Y,g(Y)\mid X)\big].
\]
Finally, since $g$ is nondecreasing, Chebyshev's covariance inequality applied under the conditional law of $Y\mid X$ yields
\[
\Cov(Y,g(Y)\mid X)\ge 0
\qquad \text{a.s.}
\]
Therefore,
\[
\Cov(e^\star,g(Y))\le 0.
\]
\end{proof}
% % LINEAR VERSION

% \begin{proposition}[Population covariance non-positivity under squared-error prediction]
% \label{prop:population_covariance_identity}


% Let $Y:=\log P$ be the log-sale price, let $X$ denote the observed feature vector, and let
% $
% f^\star(X)=\E[Y\mid X]
% $
% be the Bayes predictor\todo{ref?} under squared-error loss. Define the associated population residual
% $
% e^\star:=f^\star(X)-Y.
% $

% Then
% \begin{equation}
% \label{eq:covariance_population_identity}
% \Cov(e^\star,Y)
% =
% -\E[\Var(Y\mid X)]
% \le 0.
% \end{equation}
% \end{proposition}

% \begin{proof}
% Using $f^\star(X)=\E[Y\mid X]$, we obtain
% $
% \Cov(e^\star,Y)
% =
% \Cov(\E[Y\mid X]-Y,\;Y)
% =
% \Cov(\E[Y\mid X],Y)-\Var(Y).
% $
% By the tower property,
% \[
% \Cov(\E[Y\mid X],Y)
% =
% \E\!\big[\E[Y\mid X]\,Y\big]-\E[\E[Y\mid X]]\,\E[Y]
% =
% \E\!\big[\E[Y\mid X]^2\big]-\E[Y]^2
% =
% \Var(\E[Y\mid X]).
% \]
% Therefore,
% $
% \Cov(e^\star,Y)
% =
% \Var(\E[Y\mid X])-\Var(Y).
% $
% Applying the law of total variance,
% $
% \Var(Y)=\Var(\E[Y\mid X])+\E[\Var(Y\mid X)],
% $
% we get
% $
% \Cov(e^\star,Y)
% =
% -\E[\Var(Y\mid X)]
% \le 0.
% $
% \end{proof}

Taking $g(Y)=Y$ or $g(Y)=\exp(Y)$ in \eqref{eq:covariance_population_identity} shows that the non-positive dependence holds both between log-residuals and log-prices and between log-residuals and price-scale values. This is a population statement and does not imply that every finite-sample squared-loss fit will be regressive under the assessor-side diagnostics of Section~\ref{subsec:equity_metrics}. It does motivate augmenting \eqref{opt:base_model} with a dependence-aware correction when mitigating this behavior is a goal.

% In the present log-price formulation, $y_i=\log P_i$, the relevant object is the dependence between log-residuals and log-prices. By \eqref{eq:log_ratio_identity}, $e_i=\log r_i$, so dependence between $e$ and $y$ is equivalently dependence between log-assessment ratios and log-sale prices. This makes log-space dependence a natural target for correction: it is defined on the same scale used for estimation, while still preserving an interpretation in ratio-study terms. The next subsection introduces a general constrained and penalized formulation for such dependence-aware corrections; the following section then specializes that formulation to covariance-based choices.


% Recall from \eqref{eq:log_ratio_identity} that $e_i=\log r_i$. Therefore,
% $
% \Cov(e,y)=\Cov(\log r,\log P),
% $
% so controlling dependence between log-residuals and log-prices is equivalent to controlling first-order dependence between log-assessment ratios and log-sale prices. This log-space formulation is especially convenient in settings such as \eqref{opt:base_model}, since the predictive models are already estimated on the log-scale and multiplicative price-scale distortions become additive residual patterns.

% This population argument does not imply that every finite-sample squared-loss fit will be regressive under the assessor-side diagnostics of Section \ref{subsec:equity_metrics}. It does, however, motivate augmenting the baseline learning problem with a dependence-aware correction term. In the next subsection, we introduce a general constrained and penalized formulation for such corrections, leaving the specific choice of dependence functional open until the covariance-based constructions of the following section.

% Equation \eqref{eq:covariance_population_identity} shows that whenever some conditional log-price variation remains unexplained, squared-loss prediction induces a non-positive dependence between log-residuals and log-sale prices at the population level. This makes covariance-based dependence control a natural candidate for regressivity-aware correction, while also motivating the more general constrained and penalized formulation introduced next.

% Equation \eqref{eq:covariance_population_identity} shows that whenever some conditional log-price variation remains unexplained, squared-error prediction induces a non-positive dependence between log-residuals and log-sale prices. This identifies a population mechanism through which accuracy-oriented regression can generate regressive tendencies, and it explains why penalizing $\Cov(e,y)$ is a natural way of correcting them in settings like \eqref{opt:base_model}. 


% Recall from identity \eqref{eq:log_ratio_identity} that the log-price residuals and assessment ratios satisfy $e_i = \log r_i$. Thus, penalizing $
% \Cov(e,y)
% $ is equivalent to penalizing systematic linear dependence between log-assessment ratios and log-sale prices; $
% \Cov(e,y)=\Cov(\log r,\log P).
% $
% Further, this log-space reformulation of dependence is methodologically convenient for learning problems such as \eqref{opt:base_model}: (i) it places the dependence control on the same log-scale on which the predictive models are estimated, and (ii) its translates multiplicative deviations in the price space into additive residuals, which are usually easier patterns to learn and optimize.

% % \vspace{-4mm}
% \paragraph{Remark.} Proposition \ref{prop:population_covariance_identity} is a population statement; it does not imply that every squared-loss fit will be regressive under the assessors diagnostics defined in Section \ref{subsec:equity_metrics}.


\subsection{General Constrained and Penalized Setting}
\label{subsec:general_penalized_setting}

Sections~\ref{subsec:dependence_problem} and~\ref{subsec:squared_loss_regressivity} show that ML-based mass appraisal can induce value-dependent distortions in assessment ratios, both empirically and at the population level. We now introduce a general correction class that constrains or penalizes a sample summary of that dependence.

Let $\Psi(f,\D_n)$ denote a sample functional that quantifies value-dependent distortions in the fitted log-residuals or, equivalently, in the fitted log-assessment ratios. We can reformulate \eqref{opt:base_model} as the constrained problem
\begin{align}
\label{opt:general_constrained_model}
\min_{f\in\H}
\quad
& L(f, \D_n)+\Omega(f)
\\
\text{s.t.}\quad
& \Psi(f, \D_n)\le \tau,
\nonumber
\end{align}
for a tolerance $\tau\ge 0$. The hard constraint is conceptually useful because $\tau$ directly parameterizes the maximum dependence allowed at training. For estimation and tuning across model classes, we use its soft-constrained Lagrangian counterpart, the \emph{penalized learning problem}
\begin{align}
\label{opt:general_penalized_model}
f_\rho
\in
\arg\min_{f\in\H}
\left\{
L(f, \D_n)+\Omega(f)+\rho\,\Psi(f, \D_n)
\right\},
\end{align}
where $\rho\ge 0$ controls the tradeoff between predictive fit and regressivity correction, and $\rho=0$ recovers the baseline \eqref{opt:base_model}. Formulation \eqref{opt:general_penalized_model} captures the structure of our approach: augment a standard predictive learner with a term that discourages systematic value-dependent distortions in the fitted assessments. Sections~\ref{sec:covariance_penalty_learning} and~\ref{sec:model_specific_implementations} specialize $\Psi$ to covariance-based choices and then implement those choices in linear models and second-order tree boosting.


% \section{Covariance-Based Penalization in Standard Learners}
% \label{sec:covariance_penalty_learning_in_learners}

\subsection{Covariance-Based Penalization}
\label{sec:covariance_penalty_learning}

We now instantiate \eqref{opt:general_penalized_model} with concrete covariance-based choices of $\Psi$. By \eqref{eq:log_ratio_identity}, $e_i(f)=\log r_i(f)$, so penalizing dependence between the fitted log-residuals $e(f):=(e_1(f),\dots,e_n(f))$ and the log-prices $y:=(y_1,\dots,y_n)$ is equivalent to penalizing dependence between fitted log-assessment ratios and log-sale prices.

% Since regressivity corresponds to a systematic decrease of assessment ratios with property values, covariance-based penalties provide a natural way to instantiate the abstract correction class introduced in the previous section.

% \subsubsection{Empirical covariance target}


\subsubsection{Squared covariance penalty}

From Section~\ref{subsec:dependence_problem}, the empirical covariance $\widehat\Cov(e(f),y)$ is proportional to the slope of the least-squares regression of log-residuals on log-prices, $\operatorname{slope}(e(f)\sim y)=\widehat\Cov(e(f),y)/\widehat\Var(y)$, and $\widehat\Var(y)$ does not depend on $f$. Controlling this covariance therefore controls the first-order value-dependent trend in fitted assessment ratios. The most symmetric choice penalizes the squared covariance:
\begin{align}
    \label{eq:cov2_penalty}
    \Psi^{\mathrm{cov}^2}(f)
    :=
    \frac{1}{2}\,
    \widehat{\Cov}(e(f),y)^2.
\end{align}
Substituting \eqref{eq:cov2_penalty} into \eqref{opt:general_penalized_model} yields
\begin{align}
    \label{opt:cov2_penalized_model}
    f_{\rho}^{\mathrm{cov}^2}
    \in
    \arg\min_{f\in\H}
    \left\{
        L(f)+\Omega(f)
        +
        \frac{\rho}{2}\,\widehat{\Cov}(e(f),y)^2
    \right\}.
\end{align}
This is the most direct instantiation of the dependence-control idea developed above: it suppresses regressive and progressive first-order trends symmetrically, and is appropriate when the target is to drive the covariance toward zero.

\subsubsection{A separable surrogate penalty}

The squared covariance penalty has two limitations. As a first-order summary, covariance can miss nonmonotone dependence patterns in which local positive and negative trends cancel. Moreover, $\widehat\Cov(e(f),y)^2$ is non-separable across observations, which complicates implementation in optimizers that rely on per-observation gradients and Hessians, most notably the second-order tree-boosting interfaces that dominate mass-appraisal practice. We therefore upper-bound the squared covariance by a quantity that is additive over observations:
\begin{align}
    \label{eq:covariance_upper_bound}
    \widehat{\Cov}_n(e(f),y)^2
    &=
    \left(
        \frac{1}{n}\sum_{i=1}^n
        \bigl(e_i(f)-\bar e(f)\bigr)\bigl(y_i-\bar y\bigr)
    \right)^2
    \le
    \frac{1}{n}\sum_{i=1}^n
    \bigl(e_i(f)-\bar e(f)\bigr)^2
    \bigl(y_i-\bar y\bigr)^2,
\end{align}
where the inequality follows from Jensen's inequality. The bound still couples observations through $\bar e(f)$, so we drop the residual centering and use the separable surrogate
\begin{align}
    \label{eq:surrogate_penalty}
    \Psi^{\mathrm{surr}}(f)
    :=
    \frac{1}{n}\sum_{i=1}^n
    e_i(f)^2\bigl(y_i-\bar y\bigr)^2.
\end{align}
Substituting \eqref{eq:surrogate_penalty} into \eqref{opt:general_penalized_model} gives the surrogate-penalized problem
\begin{align}
    \label{opt:surrogate_penalized_model}
    f_{\rho}^{\mathrm{surr}}
    \in
    \arg\min_{f\in\H}
    \left\{
        L(f)+\Omega(f)
        +
        \rho\,\Psi^{\mathrm{surr}}(f)
    \right\}.
\end{align}
$\Psi^{\mathrm{surr}}$ is best understood as a covariance-motivated, observation-additive regularizer rather than as exact control of $\widehat\Cov_n(e(f),y)^2$. Two features make it attractive for assessor deployment. First, it is fully additive across observations, so it plugs into standard gradient-boosting custom-objective interfaces without modifying the leaf-splitting logic. Second, it weights each squared residual by the squared price deviation $(y_i-\bar y)^2$, so residual distortions far from the center of the value distribution are penalized most heavily. This is the property that drives the strongest empirical behavior in Section~\ref{sec:empirical}.

\paragraph{Remark.} By suppressing the residual average $\bar e(f)$, the surrogate penalizes the distance between $e_i(f)$ and zero instead of the distance to the average log-residual. On the original price scale, the analogous quantity is the squared deviation $(r_i-1)^2$. Because the surrogate weights by $(y_i-\bar y)^2$ in log space, it treats equal multiplicative deviations around the geometric center of the price distribution symmetrically, which is the scale on which the ratio--price trend is naturally read.

\paragraph{Methodological roles of the two penalties.}
The two penalties are not interchangeable. The direct penalty in \eqref{eq:cov2_penalty} is the closest realization of the dependence-control target: it suppresses first-order linear dependence between log-ratios and log-prices symmetrically and in either direction. The surrogate in \eqref{eq:surrogate_penalty} sacrifices exact covariance control in exchange for three properties that matter for deployment: it is sample-additive, so it composes cleanly with standard gradient-boosting custom-objective interfaces; it cannot be zeroed out by local cancellation between positive and negative residuals, so it responds to pointwise distortions rather than only to aggregate covariance; and its per-observation weight $(y_i-\bar y)^2$ is fixed before training, so tuning the penalty amounts to a one-dimensional line search on $\rho$. In the empirical study of Section~\ref{sec:empirical}, the surrogate delivers the strongest combined behavior on the CCAO workflow, and we recommend it as the default for routine deployment. The direct penalty is retained in the main text because it exposes the object that the surrogate approximates and because it admits the exact leaf-wise Newton implementation of Appendix~\ref{app:boosting_algorithms}.


% The surrogate \eqref{eq:surrogate_penalty} has two useful features. First, it remains anchored to the same underlying idea as the covariance-based penalties: residual distortions are weighted more heavily when they occur far from the center of the value distribution. Second, unlike \eqref{eq:cov2_penalty}, it is separable across observations once the weights $(y_i-\bar y)^2$ are fixed. This makes \eqref{opt:surrogate_penalized_model} particularly attractive for model classes and optimization procedures in which observation-wise objective contributions are computationally convenient.

% \subsubsection{Penalty comparison and methodological roles}

% The two formulations above play different roles in the proposed framework. The squared covariance penalty \eqref{eq:cov2_penalty} is the closest direct instantiation of the paper's main dependence-control idea: it suppresses first-order linear dependence between log-ratios and log-values in either direction. The surrogate penalty \eqref{eq:surrogate_penalty} sacrifices exact covariance control in exchange for observation-wise separability, making it especially attractive in learning algorithms whose optimization interfaces are naturally sample-additive.

% In this sense, the two penalties occupy different points along a trade-off between conceptual fidelity to covariance control and computational convenience. The remainder of this section shows how these covariance-based objectives can be instantiated in model classes used in practice by assessors, namely regularized linear models and boosted trees.

\subsection{Model-Specific Implementations}
\label{sec:model_specific_implementations}

We now instantiate the covariance-based penalties of Section~\ref{sec:covariance_penalty_learning} in two model classes already used in mass-appraisal practice: linear models and second-order tree boosting. Linear models offer a transparent analytical benchmark that exposes the algebraic effect of each penalty. Second-order tree boosting covers the high-performing nonlinear learners (XGBoost, LightGBM) that dominate recent assessor pipelines. The implementation choices highlighted here are what make the method drop-in compatible with standard training infrastructure: the surrogate maps directly onto per-observation custom objectives, and the squared covariance reduces to a manageable rank-one correction inside the per-round tree optimization.

% These implemented learners will later define the candidate model configurations compared in the constrained calibration procedure of Section~\ref{sec:model_selection}.

\subsubsection{Linear Covariance-Penalized Models}
\label{subsec:linear_models}

Consider the linear hypothesis class
\[
\H_{\mathrm{lin}} := \{f_\theta : x \mapsto x^\top \theta \;:\; \theta \in \Theta \subseteq \R^p\},
\]
and let $X \in \R^{n\times p}$ denote the feature matrix. For a given coefficient vector $\theta \in \Theta$, we have
$
f_\theta(x_i) = x_i^\top \theta,
$ and $
e := X\theta - y.
$ 
For notation simplicity, let us define
\[
H := I - \frac{1}{n}\mathbf{1}\mathbf{1}^\top,
\qquad
c := Hy = y - \bar y\,\mathbf{1},
\qquad
D := \diag(c_1^2,\dots,c_n^2),
\]
so that
$
\widehat{\Cov}(e,y)
=
\frac{1}{n}e^\top c.
$ 

Throughout this subsection, we take the baseline loss to be
\[
L(f_\theta) = \frac{1}{2n}\|X\theta-y\|_2^2.
\]

\paragraph{Squared covariance penalty in linear models.}
For the symmetric penalty
\[
\Psi^{\mathrm{cov}^2}(f_\theta)
=
\frac{1}{2}\,\widehat{\Cov}_n(e,y)^2
=
\frac{1}{2n^2}\bigl(e^\top c\bigr)^2,
\]
the penalized estimator solves
\begin{align}
    \label{opt:linear_cov2}
    \theta_{\rho}^{\mathrm{cov}^2}
    \in
    \arg\min_{\theta\in\Theta}
    \left\{
        \frac{1}{2n}\|X\theta-y\|_2^2
        +
        \frac{\rho}{2n^2}\bigl((X\theta-y)^\top c\bigr)^2
    \right\}.
\end{align}
Define
\[
a := X^\top c,
\qquad
b := y^\top c = c^\top c.
\]
Then the first-order optimality condition for \eqref{opt:linear_cov2} is
\[
\left(
\frac{1}{n}X^\top X + \frac{\rho}{n^2}aa^\top
\right)\theta
=
\frac{1}{n}X^\top y + \frac{\rho}{n^2}ba.
\]
Hence, whenever the displayed inverse exists,
\begin{align}
    \label{eq:linear_cov2_closed_form}
    \theta_{\rho}^{\mathrm{cov}^2}
    =
    \left(
        \frac{1}{n}X^\top X
        +
        \frac{\rho}{n^2}aa^\top
    \right)^{-1}
    \left(
        \frac{1}{n}X^\top y
        +
        \frac{\rho}{n^2}ba
    \right).
\end{align}
This shows that the squared covariance penalty induces a rank-one modification of the usual normal equations.

\paragraph{Surrogate penalty in linear models.}
For the separable surrogate
\[
\Psi^{\mathrm{surr}}(f_\theta)
=
\frac{1}{n}\sum_{i=1}^n e_i^2 c_i^2
=
\frac{1}{n}e^\top D\,e,
\]
the penalized problem is
\begin{align}
    \label{opt:linear_surr}
    \theta_{\rho}^{\mathrm{surr}}
    \in
    \arg\min_{\theta\in\Theta}
    \left\{
        \frac{1}{2n}\|X\theta-y\|_2^2
        +
        \frac{\rho}{n}e^\top D\,e
    \right\}.
\end{align}
The corresponding optimality condition is
\[
\left(
\frac{1}{n}X^\top X
+
\frac{2\rho}{n}X^\top D X
\right)\theta
=
\frac{1}{n}X^\top y
+
\frac{2\rho}{n}X^\top D y,
\]
and therefore, whenever the displayed inverse exists,
\begin{align}
    \label{eq:linear_surr_closed_form}
    \theta_{\rho}^{\mathrm{surr}}
    =
    \left(
        \frac{1}{n}X^\top X
        +
        \frac{2\rho}{n}X^\top D X
    \right)^{-1}
    \left(
        \frac{1}{n}X^\top y
        +
        \frac{2\rho}{n}X^\top D y
    \right).
\end{align}

\paragraph{Extension to regularized linear models.}
The same constructions extend to penalized linear estimators
\[
\min_{\theta\in\Theta}
\left\{
    L(f_\theta) + \Omega(\theta) + \rho\,\Psi(f_\theta)
\right\}
\]
whenever $\Omega$ is a convex, differentiable regularizer. Non-differentiable cases such as the LASSO remain convex but lose the closed form and must be solved numerically. In both regimes, the covariance penalty integrates into standard regularized linear regression as a small modification of the normal equations.

\subsubsection{Covariance-Penalized Second-Order Tree Boosting}
\label{subsec:tree_boosting}

We next consider second-order additive tree boosting under squared loss, as used in gradient-boosted tree systems such as XGBoost and LightGBM. Let \(f^{(m-1)}\) denote the current prediction function at boosting iteration \(m\), and write
\[
f_i := f^{(m-1)}(x_i),
\qquad
e_i := f_i - y_i,
\qquad
c_i := y_i - \bar y,
\qquad
z := \widehat{\Cov}_n(e,y) = \frac{1}{n}\sum_{i=1}^n e_i c_i.
\]
At step \(m\), one adds a regression tree \(t\) to the current predictor, so that
\[
f^{(m)}(x) = f^{(m-1)}(x) + \nu\,t(x),
\]
where \(\nu\in(0,1]\) is the learning rate.

\paragraph{Adding dependence penalties into the boosting objective.}
For standard second-order boosting, the objective is approximated locally by
\begin{align}
    \label{eq:boosting_local_objective}
    \widetilde{\mathcal J}(t)
    =
    \sum_{i=1}^n
    \left(
        g_i\,t_i + \frac{1}{2}h_i\,t_i^2
    \right)
    +
    \Omega(t),
\end{align}
where \(t_i := t(x_i)\), \(\Omega(t)\) is the tree-complexity penalty, and
\[
g_i := \frac{\partial}{\partial f_i}
\left(
L(f) + \rho\,\Psi(f)
\right),
\qquad
h_i := \frac{\partial^2}{\partial f_i^2}
\left(
L(f) + \rho\,\Psi(f)
\right)
\]
are evaluated at the current predictions. Under squared loss,
\[
\frac{\partial L(f)}{\partial f_i} = \frac{e_i}{n},
\qquad
\frac{\partial^2 L(f)}{\partial f_i^2} = \frac{1}{n}.
\]
The covariance-based penalties modify these quantities as follows.

\paragraph{Squared covariance penalty in boosting objective.}
For \(\Psi^{\mathrm{cov}^2}(f)=\tfrac12 z^2\), the gradient is
\begin{align}
    \label{eq:boosting_cov2_gradient}
    g_i^{\mathrm{cov}^2}
    =
    \frac{e_i}{n}
    +
    \frac{\rho}{n}z\,c_i.
\end{align}
The exact Hessian of the penalized objective with respect to the prediction vector \(f=(f_1,\dots,f_n)\) is
\begin{align}
    \label{eq:boosting_cov2_exact_hessian}
    \nabla_f^2
    \bigl(L(f)+\rho\,\Psi^{\mathrm{cov}^2}(f)\bigr)
    =
    \frac{1}{n}I
    +
    \frac{\rho}{n^2}cc^\top.
\end{align}
Unlike the separable boosting objectives used in standard libraries, \eqref{eq:boosting_cov2_exact_hessian} is dense: each observation is coupled to every other observation through the rank-one term \(cc^\top\). Consequently, there are two natural implementation routes.

The first is a practical approximation that keeps only the diagonal terms,
\begin{align}
    \label{eq:boosting_cov2_diag_hessian}
    h_i^{\mathrm{cov}^2,\mathrm{diag}}
    =
    \frac{1}{n}
    +
    \frac{\rho}{n^2}c_i^2.
\end{align}
This produces a custom objective compatible with standard second-order boosting APIs, at the price of ignoring the cross-observation curvature induced by the covariance penalty.

The second is an exact leaf-wise implementation for a fixed tree structure. Suppose a candidate tree \(t\) has leaves \(I_1,\dots,I_J\) and leaf weights \(w=(w_1,\dots,w_J)\), so that \(t_i=w_j\) for \(i\in I_j\). Let
\[
n_j := |I_j|,
\qquad
b_j := \sum_{i\in I_j} e_i,
\qquad
s_j := \sum_{i\in I_j} c_i,
\]
and let \(s=(s_1,\dots,s_J)^\top\). If \(\Omega(t)=\gamma J + \frac{\lambda}{2}\|w\|_2^2\), then the exact local objective for the squared covariance penalty can be written as
\begin{align}
    \label{eq:boosting_cov2_exact_leaf_objective}
    \widetilde{\mathcal J}_{\mathrm{cov}^2}(w)
    =
    g_T^\top w
    +
    \frac{1}{2}w^\top A_T w
    +
    \gamma J
    + \text{const},
\end{align}
where
\[
g_{T,j}
=
\frac{b_j}{n}
+
\frac{\rho}{n}z\,s_j,
\qquad
A_T
=
\diag\!\left(\frac{n_1}{n}+\lambda,\dots,\frac{n_J}{n}+\lambda\right)
+
\frac{\rho}{n^2}ss^\top.
\]
Hence the optimal leaf weights satisfy
\begin{align}
    \label{eq:boosting_cov2_exact_leaf_system}
    A_T w^\star = -g_T,
    \qquad
    w^\star = -A_T^{-1}g_T,
\end{align}
whenever \(A_T\) is invertible. The associated gain for the fixed structure \(T\) is
\begin{align}
    \label{eq:boosting_cov2_exact_gain}
    \operatorname{Gain}(T)
    =
    \frac{1}{2}g_T^\top A_T^{-1}g_T - \gamma J.
\end{align}
This exact route preserves the full curvature of the covariance penalty, but requires a custom tree-construction procedure rather than the usual observation-wise Hessian interface.

\paragraph{Surrogate penalty as a sample-additive boosting objective.}
For the surrogate
\[
\Psi^{\mathrm{surr}}(f)
=
\frac{1}{n}\sum_{i=1}^n e_i^2 c_i^2,
\]
the gradient and Hessian remain observation-wise separable:
\begin{align}
    \label{eq:boosting_surr_grad_hess}
    g_i^{\mathrm{surr}}
    &=
    \frac{e_i}{n}
    +
    \frac{2\rho}{n}e_i c_i^2,
    \\
    h_i^{\mathrm{surr}}
    &=
    \frac{1}{n}
    +
    \frac{2\rho}{n}c_i^2.
\end{align}
Accordingly, the surrogate penalty is particularly well suited to off-the-shelf second-order boosting implementations: once the weights \(c_i^2\) are precomputed, the penalized objective can be supplied through standard custom gradient and Hessian routines.

\paragraph{Implementation summary.}
The squared covariance penalty is the most direct implementation of the covariance-control idea, but it introduces dense curvature (Hessian); in practice, it can be handled either through a diagonal Hessian approximation in standard boosting libraries (XGBoost or LightGBM), or through an exact custom leaf solver of the form \eqref{eq:boosting_cov2_exact_leaf_system}. The surrogate penalty is fully sample-additive, and is therefore the most convenient to implement through standard observation-wise gradient and Hessian interfaces.

The implementation details in this section show that the covariance-based framework is compatible with both linear estimators and high-performing nonlinear learners. This model-specific flexibility is important for the practical goal of the paper: regressivity correction can be layered onto standard predictive tools already used in assessment workflows, rather than requiring a separate valuation architecture.

\paragraph{Remark.}
For completeness, Appendix~\ref{app:boosting_algorithms} provides pseudocode for an exact Newton-style tree-boosting implementation of the squared covariance penalty. 

% Together, the learner--penalty combinations developed in this section define the candidate model configurations later compared in the constrained calibration procedure of Section~\ref{sec:model_selection}.

% \section{Covariance-Based Penalizations in Standard Learners}
% \label{sec:covariance_penalty_learning_in_learners}

% \subsection{Covariance-Based Penalization Terms}
% \label{sec:covariance_penalty_learning}

% We now instantiate the general penalized learning problem \eqref{opt:general_penalized_model} by specifying concrete choices of the functional $\Psi$ based on empirical covariance-type measures. For simplicity, let
% \[
% \H := \{f_\theta : \theta \in \Theta\}
% \]
% denote the hypothesis class induced by the parameter space $\Theta$, and for $f \in \H$ write
% \[
% \hat y_i(f) = f(x_i),
% \qquad
% e_i(f) = \hat y_i(f) - y_i,
% \qquad
% r_i(f) = \exp(e_i(f)).
% \]
% By \eqref{eq:log_ratio_identity}, $e_i(f)=\log r_i(f)$, while $y_i=\log P_i$. Therefore, penalizing dependence between $e(f):=(e_1(f),\dots,e_n(f))$ and $y:=(y_1,\dots,y_n)$ is equivalent to penalizing dependence between fitted log-assessment ratios and log-sale prices. Since regressivity corresponds to a systematic decrease of assessment ratios with property values, covariance-based penalties provide a natural way to instantiate the abstract correction class introduced in the previous section.

% \subsubsection{Empirical covariance target}

% Let
% \[
% \bar e(f) := \frac{1}{n}\sum_{i=1}^n e_i(f),
% \qquad
% \bar y := \frac{1}{n}\sum_{i=1}^n y_i,
% \]
% and define the empirical covariance
% \begin{align}
%     \label{eq:empirical_covariance}
%     \widehat{\Cov}_n(e(f),y)
%     :=
%     \frac{1}{n}\sum_{i=1}^n
%     \bigl(e_i(f)-\bar e(f)\bigr)\bigl(y_i-\bar y\bigr).
% \end{align}
% Recall from Section \ref{subsec:dependence_problem} that this quantity is directly connected to the slope of the least-squares regression of $e_i(f)$ on $y_i$, since
% \[
% \widehat{\beta}_{e\sim y}(f)
% =
% \frac{\widehat{\Cov}_n(e(f),y)}{\widehat{\Var}_n(y)},
% \]
% and $\widehat{\Var}_n(y)$ does not depend on the model choice. Accordingly, controlling the empirical covariance is a natural way to control first-order value-dependent trends in the fitted assessment ratios.

% \subsubsection{Squared covariance penalty}

% A first class of penalty is obtained by using the empirical covariance in \eqref{eq:empirical_covariance} directly. The most symmetric choice is to penalize the squared covariance:
% \begin{align}
%     \label{eq:cov2_penalty}
%     \Psi^{\mathrm{cov}^2}(f)
%     :=
%     \frac{1}{2}\,
%     \widehat{\Cov}_n(e(f),y)^2.
% \end{align}
% Substituting \eqref{eq:cov2_penalty} into \eqref{opt:general_penalized_model} yields
% \begin{align}
%     \label{opt:cov2_penalized_model}
%     f_{\rho}^{\mathrm{cov}^2}
%     \in
%     \arg\min_{f\in\H}
%     \left\{
%         L(f)+\Omega(f)
%         +
%         \frac{\rho}{2}\,\widehat{\Cov}_n(e(f),y)^2
%     \right\}.
% \end{align}
% This is the most direct instantiation of the dependence-control idea developed in the previous section. It penalizes any linear dependence between $e(f)$ and $y$, regardless of sign. In particular, it suppresses both regressive and progressive linear trends, and is therefore appropriate when the goal is to drive the first-order dependence toward zero. For this reason, we view \eqref{opt:cov2_penalized_model} as the canonical covariance-based penalization.

% \subsubsection{A separable surrogate penalty}

% While the direct covariance penalty above is conceptually natural, it has two important limitations. First, covariance is a first-order summary and may fail to reflect more complex dependence patterns when positive and negative local trends cancel each other out. Second, the term $\widehat{\Cov}_n(e(f),y)^2$ is non-separable across observations, which can be inconvenient in optimization methods that rely on observation-wise gradients or second-order approximations.

% To obtain a computationally simpler covariance-like regularizer, we upper-bound the squared empirical covariance as
% \begin{align}
%     \label{eq:covariance_upper_bound}
%     \widehat{\Cov}_n(e(f),y)^2
%     &=
%     \left(
%         \frac{1}{n}\sum_{i=1}^n
%         \bigl(e_i(f)-\bar e(f)\bigr)\bigl(y_i-\bar y\bigr)
%     \right)^2
%     \le
%     \frac{1}{n}\sum_{i=1}^n
%     \bigl(e_i(f)-\bar e(f)\bigr)^2
%     \bigl(y_i-\bar y\bigr)^2,
% \end{align}
% where the inequality follows from Jensen's inequality. This upper bound is still coupled across observations because it depends on the sample mean $\bar e(f)$. To obtain a separable approximation that is easier to use in optimization routines built around observation-wise objective contributions, we drop the centering of the residuals and consider the simplified surrogate
% \begin{align}
%     \label{eq:surrogate_penalty}
%     \Psi^{\mathrm{surr}}(f)
%     :=
%     \frac{1}{n}\sum_{i=1}^n
%     e_i(f)^2\bigl(y_i-\bar y\bigr)^2.
% \end{align}
% Substituting \eqref{eq:surrogate_penalty} into \eqref{opt:general_penalized_model} gives the surrogate-penalized problem
% \begin{align}
%     \label{opt:surrogate_penalized_model}
%     f_{\rho}^{\mathrm{surr}}
%     \in
%     \arg\min_{f\in\H}
%     \left\{
%         L(f)+\Omega(f)
%         +
%         \rho\,\Psi^{\mathrm{surr}}(f)
%     \right\}.
% \end{align}
% Thus, $\Psi^{\mathrm{surr}}$ should be viewed as a separable approximation motivated by covariance control, rather than as exact control of $\widehat{\Cov}_n(e(f),y)^2$ itself.

% The surrogate \eqref{eq:surrogate_penalty} has two useful features. First, it remains anchored to the same underlying idea as the covariance-based penalties: residual distortions are weighted more heavily when they occur far from the center of the value distribution. Second, unlike \eqref{eq:cov2_penalty}, it is separable across observations once the weights $(y_i-\bar y)^2$ are fixed. This makes \eqref{opt:surrogate_penalized_model} particularly attractive for model classes and optimization procedures in which observation-wise objective contributions are computationally convenient.

% \subsubsection{Penalty comparison and methodological roles}

% The two formulations above play different roles in the proposed framework. The squared covariance penalty \eqref{eq:cov2_penalty} is the closest direct instantiation of the paper's main dependence-control idea: it suppresses first-order linear dependence between log-ratios and log-values in either direction. The surrogate penalty \eqref{eq:surrogate_penalty} sacrifices exact covariance control in exchange for observation-wise separability, making it especially attractive in learning algorithms whose optimization interfaces are naturally sample-additive.

% In this sense, the two penalties occupy different points along a trade-off between conceptual fidelity to covariance control and computational convenience. The remainder of this section shows how these covariance-based objectives can be instantiated in model classes used in practice by assessors, namely regularized linear models and boosted trees.

% \subsection{Model-Specific Implementations}
% \label{sec:model_specific_implementations}

% We now show how the covariance-based penalizations introduced in Section~\ref{sec:covariance_penalty_learning} can be instantiated in specific model classes. The goal of this subsection is not to introduce new regressivity criteria, but to show that the previously defined covariance-based corrections can be embedded into predictive tools already familiar in mass-appraisal workflows. We focus on two examples that are especially relevant in practice: linear models and second-order tree boosting. The former provides a transparent analytical benchmark that clarifies the algebraic effect of each penalty, while the latter covers widely used high-performing nonlinear learners such as XGBoost- and LightGBM-type methods. These implemented learners will later define the candidate model configurations compared in the constrained calibration procedure of Section~\ref{sec:model_selection}.

% \subsubsection{Linear Covariance-Penalized Models}
% \label{subsec:linear_models}

% Consider the linear hypothesis class
% \[
% \H_{\mathrm{lin}} := \{f_\theta : x \mapsto x^\top \theta \;:\; \theta \in \Theta \subseteq \R^p\},
% \]
% and let $X \in \R^{n\times p}$ denote the design matrix. For a coefficient vector $\theta \in \Theta$, define
% \[
% f_\theta(x_i) = x_i^\top \theta,
% \qquad
% e := X\theta - y.
% \]
% As in the previous subsection, let
% \[
% H := I - \frac{1}{n}\mathbf{1}\mathbf{1}^\top,
% \qquad
% c := Hy = y - \bar y\,\mathbf{1},
% \qquad
% D := \diag(c_1^2,\dots,c_n^2),
% \]
% so that
% \[
% \widehat{\Cov}_n(e,y)
% =
% \frac{1}{n}e^\top c.
% \]

% Throughout this subsection, we take the baseline loss to be
% \[
% L(f_\theta) = \frac{1}{2n}\|X\theta-y\|_2^2.
% \]

% \paragraph{Squared covariance penalty in linear models.}
% For the symmetric penalty
% \[
% \Psi^{\mathrm{cov}^2}(f_\theta)
% =
% \frac{1}{2}\,\widehat{\Cov}_n(e,y)^2
% =
% \frac{1}{2n^2}\bigl(e^\top c\bigr)^2,
% \]
% the penalized estimator solves
% \begin{align}
%     \label{opt:linear_cov2}
%     \theta_{\rho}^{\mathrm{cov}^2}
%     \in
%     \arg\min_{\theta\in\Theta}
%     \left\{
%         \frac{1}{2n}\|X\theta-y\|_2^2
%         +
%         \frac{\rho}{2n^2}\bigl((X\theta-y)^\top c\bigr)^2
%     \right\}.
% \end{align}
% Define
% \[
% a := X^\top c,
% \qquad
% b := y^\top c = c^\top c.
% \]
% Then the first-order optimality condition for \eqref{opt:linear_cov2} is
% \[
% \left(
% \frac{1}{n}X^\top X + \frac{\rho}{n^2}aa^\top
% \right)\theta
% =
% \frac{1}{n}X^\top y + \frac{\rho}{n^2}ba.
% \]
% Hence, whenever the displayed inverse exists,
% \begin{align}
%     \label{eq:linear_cov2_closed_form}
%     \theta_{\rho}^{\mathrm{cov}^2}
%     =
%     \left(
%         \frac{1}{n}X^\top X
%         +
%         \frac{\rho}{n^2}aa^\top
%     \right)^{-1}
%     \left(
%         \frac{1}{n}X^\top y
%         +
%         \frac{\rho}{n^2}ba
%     \right).
% \end{align}
% This shows that the squared covariance penalty induces a rank-one modification of the usual normal equations.

% \paragraph{Surrogate penalty in linear models.}
% For the separable surrogate
% \[
% \Psi^{\mathrm{surr}}(f_\theta)
% =
% \frac{1}{n}\sum_{i=1}^n e_i^2 c_i^2
% =
% \frac{1}{n}e^\top D\,e,
% \]
% the penalized problem is
% \begin{align}
%     \label{opt:linear_surr}
%     \theta_{\rho}^{\mathrm{surr}}
%     \in
%     \arg\min_{\theta\in\Theta}
%     \left\{
%         \frac{1}{2n}\|X\theta-y\|_2^2
%         +
%         \frac{\rho}{n}e^\top D\,e
%     \right\}.
% \end{align}
% The corresponding optimality condition is
% \[
% \left(
% \frac{1}{n}X^\top X
% +
% \frac{2\rho}{n}X^\top D X
% \right)\theta
% =
% \frac{1}{n}X^\top y
% +
% \frac{2\rho}{n}X^\top D y,
% \]
% and therefore, whenever the displayed inverse exists,
% \begin{align}
%     \label{eq:linear_surr_closed_form}
%     \theta_{\rho}^{\mathrm{surr}}
%     =
%     \left(
%         \frac{1}{n}X^\top X
%         +
%         \frac{2\rho}{n}X^\top D X
%     \right)^{-1}
%     \left(
%         \frac{1}{n}X^\top y
%         +
%         \frac{2\rho}{n}X^\top D y
%     \right).
% \end{align}

% \paragraph{Extension to regularized linear models.}
% The same constructions extend immediately to penalized linear estimators of the form
% \[
% \min_{\theta\in\Theta}
% \left\{
%     L(f_\theta) + \mathcal{R} + \rho\,\Psi(f_\theta)
% \right\},
% \]
% where \(\mathcal{R}\) is a complexity regularizer on the coefficients. In the ridge case, the resulting systems remain quadratic and admit closed-form solutions after adding the usual ridge term to the coefficient matrix. In contrast, \(\ell_1\)-regularized variants such as the LASSO remain convex but no longer admit closed-form expressions, and must instead be solved numerically.

% \paragraph{What the linear case teaches.}
% Linear models clarify the algebraic effect of the penalties and show that the method can be integrated into standard regularized regression with minimal conceptual disruption.

% \subsubsection{Covariance-Penalized Second-Order Tree Boosting}
% \label{subsec:tree_boosting}

% We next consider second-order additive tree boosting under squared loss, as used in gradient-boosted tree systems such as XGBoost and LightGBM. Let \(f^{(m-1)}\) denote the current prediction function at boosting iteration \(m\), and write
% \[
% f_i := f^{(m-1)}(x_i),
% \qquad
% e_i := f_i - y_i,
% \qquad
% c_i := y_i - \bar y,
% \qquad
% z := \widehat{\Cov}_n(e,y) = \frac{1}{n}\sum_{i=1}^n e_i c_i.
% \]
% At step \(m\), one adds a regression tree \(t\) to the current predictor, so that
% \[
% f^{(m)}(x) = f^{(m-1)}(x) + \nu\,t(x),
% \]
% where \(\nu\in(0,1]\) is the learning rate.

% \paragraph{Adding dependence penalties into the boosting objective.}
% For standard second-order boosting, the objective is approximated locally by
% \begin{align}
%     \label{eq:boosting_local_objective}
%     \widetilde{\mathcal J}(t)
%     =
%     \sum_{i=1}^n
%     \left(
%         g_i\,t_i + \frac{1}{2}h_i\,t_i^2
%     \right)
%     +
%     \Omega(t),
% \end{align}
% where \(t_i := t(x_i)\), \(\Omega(t)\) is the tree-complexity penalty, and
% \[
% g_i := \frac{\partial}{\partial f_i}
% \left(
% L(f) + \rho\,\Psi(f)
% \right),
% \qquad
% h_i := \frac{\partial^2}{\partial f_i^2}
% \left(
% L(f) + \rho\,\Psi(f)
% \right)
% \]
% are evaluated at the current predictions. Under squared loss,
% \[
% \frac{\partial L(f)}{\partial f_i} = \frac{e_i}{n},
% \qquad
% \frac{\partial^2 L(f)}{\partial f_i^2} = \frac{1}{n}.
% \]
% The covariance-based penalties modify these quantities as follows.

% \paragraph{Squared covariance penalty in boosting objective.}
% For \(\Psi^{\mathrm{cov}^2}(f)=\tfrac12 z^2\), the gradient is
% \begin{align}
%     \label{eq:boosting_cov2_gradient}
%     g_i^{\mathrm{cov}^2}
%     =
%     \frac{e_i}{n}
%     +
%     \frac{\rho}{n}z\,c_i.
% \end{align}
% The exact Hessian of the penalized objective with respect to the prediction vector \(f=(f_1,\dots,f_n)\) is
% \begin{align}
%     \label{eq:boosting_cov2_exact_hessian}
%     \nabla_f^2
%     \bigl(L(f)+\rho\,\Psi^{\mathrm{cov}^2}(f)\bigr)
%     =
%     \frac{1}{n}I
%     +
%     \frac{\rho}{n^2}cc^\top.
% \end{align}
% Unlike the separable boosting objectives used in standard libraries, \eqref{eq:boosting_cov2_exact_hessian} is dense: each observation is coupled to every other observation through the rank-one term \(cc^\top\). Consequently, there are two natural implementation routes.

% The first is a practical approximation that keeps only the diagonal terms,
% \begin{align}
%     \label{eq:boosting_cov2_diag_hessian}
%     h_i^{\mathrm{cov}^2,\mathrm{diag}}
%     =
%     \frac{1}{n}
%     +
%     \frac{\rho}{n^2}c_i^2.
% \end{align}
% This produces a custom objective compatible with standard second-order boosting APIs, at the price of ignoring the cross-observation curvature induced by the covariance penalty.

% The second is an exact leaf-wise implementation for a fixed tree structure. Suppose a candidate tree \(t\) has leaves \(I_1,\dots,I_J\) and leaf weights \(w=(w_1,\dots,w_J)\), so that \(t_i=w_j\) for \(i\in I_j\). Let
% \[
% n_j := |I_j|,
% \qquad
% b_j := \sum_{i\in I_j} e_i,
% \qquad
% s_j := \sum_{i\in I_j} c_i,
% \]
% and let \(s=(s_1,\dots,s_J)^\top\). If \(\Omega(t)=\gamma J + \frac{\lambda}{2}\|w\|_2^2\), then the exact local objective for the squared covariance penalty can be written as
% \begin{align}
%     \label{eq:boosting_cov2_exact_leaf_objective}
%     \widetilde{\mathcal J}_{\mathrm{cov}^2}(w)
%     =
%     g_T^\top w
%     +
%     \frac{1}{2}w^\top A_T w
%     +
%     \gamma J
%     + \text{const},
% \end{align}
% where
% \[
% g_{T,j}
% =
% \frac{b_j}{n}
% +
% \frac{\rho}{n}z\,s_j,
% \qquad
% A_T
% =
% \diag\!\left(\frac{n_1}{n}+\lambda,\dots,\frac{n_J}{n}+\lambda\right)
% +
% \frac{\rho}{n^2}ss^\top.
% \]
% Hence the optimal leaf weights satisfy
% \begin{align}
%     \label{eq:boosting_cov2_exact_leaf_system}
%     A_T w^\star = -g_T,
%     \qquad
%     w^\star = -A_T^{-1}g_T,
% \end{align}
% whenever \(A_T\) is invertible. The associated gain for the fixed structure \(T\) is
% \begin{align}
%     \label{eq:boosting_cov2_exact_gain}
%     \operatorname{Gain}(T)
%     =
%     \frac{1}{2}g_T^\top A_T^{-1}g_T - \gamma J.
% \end{align}
% This exact route preserves the full curvature of the covariance penalty, but requires a custom tree-construction procedure rather than the usual observation-wise Hessian interface.

% \paragraph{Surrogate penalty as a sample-additive boosting objective.}
% For the surrogate
% \[
% \Psi^{\mathrm{surr}}(f)
% =
% \frac{1}{n}\sum_{i=1}^n e_i^2 c_i^2,
% \]
% the gradient and Hessian remain observation-wise separable:
% \begin{align}
%     \label{eq:boosting_surr_grad_hess}
%     g_i^{\mathrm{surr}}
%     &=
%     \frac{e_i}{n}
%     +
%     \frac{2\rho}{n}e_i c_i^2,
%     \\
%     h_i^{\mathrm{surr}}
%     &=
%     \frac{1}{n}
%     +
%     \frac{2\rho}{n}c_i^2.
% \end{align}
% Accordingly, the surrogate penalty is particularly well suited to off-the-shelf second-order boosting implementations: once the weights \(c_i^2\) are precomputed, the penalized objective can be supplied through standard custom gradient and Hessian routines.

% \paragraph{Implementation summary.}
% The two penalties above lead to distinct computational regimes in boosted trees. The squared covariance penalty is the most direct implementation of the covariance-control idea, but it introduces dense curvature; in practice, it can be handled either through a diagonal Hessian approximation in standard boosting libraries or through an exact custom leaf solver of the form \eqref{eq:boosting_cov2_exact_leaf_system}. The surrogate penalty is fully sample-additive, and is therefore the most convenient to implement through standard observation-wise gradient and Hessian interfaces.

% The formulas in this section show that the covariance-based framework is compatible with both transparent linear estimators and high-performing nonlinear learners. This model-specific flexibility is important for the practical goal of the paper: regressivity correction can be layered onto standard predictive tools already used in assessment workflows, rather than requiring a separate valuation architecture.

% \paragraph{Algorithmic details.}
% For completeness, Appendix~\ref{app:boosting_algorithms} provides pseudocode for an exact Newton-style tree-boosting implementation of the squared covariance penalty, together with an optional first-order approximation. In the main text, we focus on the gradient, Hessian, and leaf-system formulas that define the method and clarify its relation to standard boosted-tree implementations.

% Together, the learner--penalty combinations developed in this section define the candidate model configurations later compared in the constrained calibration procedure of Section~\ref{sec:model_selection}.



% \section{Covariance-Based Penalizations in Standard Learners}
% \label{sec:covariance_penalty_learning_in_learners}

% \subsection{Covariance-Based Penalization Terms}
% \label{sec:covariance_penalty_learning}

% We now instantiate the general penalized learning problem \eqref{opt:general_penalized_model} by specifying concrete choices of the functional $\Psi$ based on empirical covariance-type measures. For simplicity, let
% \[
% \H := \{f_\theta : \theta \in \Theta\}
% \]
% denote the hypothesis class induced by the parameter space $\Theta$, and for $f \in \H$ write
% \[
% \hat y_i(f) = f(x_i),
% \qquad
% e_i(f) = \hat y_i(f) - y_i,
% \qquad
% r_i(f) = \exp(e_i(f)).
% \]
% By \eqref{eq:log_ratio_identity}, $e_i(f)=\log r_i(f)$, while $y_i=\log P_i$. Therefore, penalizing dependence between $e(f):=(e_1(f),\dots,e_n(f))$ and $y:=(y_1,\dots,y_n)$ is equivalent to penalizing dependence between fitted log-assessment ratios and log-sale prices. Since regressivity corresponds to a systematic decrease of assessment ratios with property values, covariance-based penalties provide a natural way to instantiate the abstract correction class introduced in the previous section.

% \subsubsection{Empirical covariance target}

% Let
% \[
% \bar e(f) := \frac{1}{n}\sum_{i=1}^n e_i(f),
% \qquad
% \bar y := \frac{1}{n}\sum_{i=1}^n y_i,
% \]
% and define the empirical covariance
% \begin{align}
%     \label{eq:empirical_covariance}
%     \widehat{\Cov}_n(e(f),y)
%     :=
%     \frac{1}{n}\sum_{i=1}^n
%     \bigl(e_i(f)-\bar e(f)\bigr)\bigl(y_i-\bar y\bigr).
% \end{align}
% Recall from Section \ref{subsec:dependence_problem} that this quantity is directly connected to the slope of the least-squares regression of $e_i(f)$ on $y_i$, since
% \[
% \widehat{\beta}_{e\sim y}(f)
% =
% \frac{\widehat{\Cov}_n(e(f),y)}{\widehat{\Var}_n(y)},
% \]
% and $\widehat{\Var}_n(y)$ does not depend on the model choice. Accordingly, controlling the empirical covariance is a natural way to control first-order value-dependent trends in the fitted assessment ratios.

% \subsubsection{Squared covariance penalty}

% A first class of penalty is obtained by using the empirical covariance in \eqref{eq:empirical_covariance} directly. The most symmetric choice is to penalize the squared covariance:
% \begin{align}
%     \label{eq:cov2_penalty}
%     \Psi^{\mathrm{cov}^2}(f)
%     :=
%     \frac{1}{2}\,
%     \widehat{\Cov}_n(e(f),y)^2.
% \end{align}
% Substituting \eqref{eq:cov2_penalty} into \eqref{opt:general_penalized_model} yields
% \begin{align}
%     \label{opt:cov2_penalized_model}
%     f_{\rho}^{\mathrm{cov}^2}
%     \in
%     \arg\min_{f\in\H}
%     \left\{
%         L(f)+\Omega(f)
%         +
%         \frac{\rho}{2}\,\widehat{\Cov}_n(e(f),y)^2
%     \right\}.
% \end{align}
% This is the most direct instantiation of the dependence-control idea developed in the previous section. It penalizes any linear dependence between $e(f)$ and $y$, regardless of sign. In particular, it suppresses both regressive and progressive linear trends, and is therefore appropriate when the goal is to drive the first-order dependence toward zero. For this reason, we view \eqref{opt:cov2_penalized_model} as the canonical covariance-based penalization.

% \subsubsection{A separable surrogate penalty}

% While the direct covariance penalty above is conceptually natural, it has two important limitations. First, covariance is a first-order summary and may fail to reflect more complex dependence patterns when positive and negative local trends cancel each other out. Second, the term $\widehat{\Cov}_n(e(f),y)^2$ is non-separable across observations, which can be inconvenient in optimization methods that rely on observation-wise gradients or second-order approximations.

% To obtain a computationally simpler covariance-like regularizer, we upper-bound the squared empirical covariance as
% \begin{align}
%     \label{eq:covariance_upper_bound}
%     \widehat{\Cov}_n(e(f),y)^2
%     &=
%     \left(
%         \frac{1}{n}\sum_{i=1}^n
%         \bigl(e_i(f)-\bar e(f)\bigr)\bigl(y_i-\bar y\bigr)
%     \right)^2
%     \le
%     \frac{1}{n}\sum_{i=1}^n
%     \bigl(e_i(f)-\bar e(f)\bigr)^2
%     \bigl(y_i-\bar y\bigr)^2,
% \end{align}
% where the inequality follows from Jensen's inequality. This upper bound is still coupled across observations because it depends on the sample mean $\bar e(f)$. To obtain a separable approximation that is easier to use in optimization routines built around observation-wise objective contributions, we drop the centering of the residuals and consider the simplified surrogate
% \begin{align}
%     \label{eq:surrogate_penalty}
%     \Psi^{\mathrm{surr}}(f)
%     :=
%     \frac{1}{n}\sum_{i=1}^n
%     e_i(f)^2\bigl(y_i-\bar y\bigr)^2.
% \end{align}
% Substituting \eqref{eq:surrogate_penalty} into \eqref{opt:general_penalized_model} gives the surrogate-penalized problem
% \begin{align}
%     \label{opt:surrogate_penalized_model}
%     f_{\rho}^{\mathrm{surr}}
%     \in
%     \arg\min_{f\in\H}
%     \left\{
%         L(f)+\Omega(f)
%         +
%         \rho\,\Psi^{\mathrm{surr}}(f)
%     \right\}.
% \end{align}
% Thus, $\Psi^{\mathrm{surr}}$ should be viewed as a separable approximation motivated by covariance control, rather than as exact control of $\widehat{\Cov}_n(e(f),y)^2$ itself.

% The surrogate \eqref{eq:surrogate_penalty} has two useful features. First, it remains anchored to the same underlying idea as the covariance penalty: residual distortions are weighted more heavily when they occur far from the center of the value distribution. Second, unlike \eqref{eq:cov2_penalty}, it is separable across observations once the weights $(y_i-\bar y)^2$ are fixed. This makes \eqref{opt:surrogate_penalized_model} particularly attractive for model classes and optimization procedures in which observation-wise objective contributions are computationally convenient.

% \subsubsection{Penalty comparison and methodological roles}

% The two formulations above play different roles in the proposed framework. The squared covariance penalty \eqref{eq:cov2_penalty} is the closest direct instantiation of the paper's main dependence-control idea: it suppresses first-order linear dependence between log-ratios and log-values in either direction. The surrogate penalty \eqref{eq:surrogate_penalty} sacrifices exact covariance control in exchange for observation-wise separability, making it especially attractive in learning algorithms whose optimization interfaces are naturally sample-additive.

% In this sense, the two penalties occupy different points along a trade-off between conceptual fidelity to covariance control and computational convenience. The remainder of this section shows how these two covariance-based objectives can be instantiated in model classes used in practice by assessors, namely regularized linear models and boosted trees.

% \subsection{Model-Specific Implementations}
% \label{sec:model_specific_implementations}

% We now show how the two covariance-based penalizations introduced above can be instantiated in specific model classes. The goal of this subsection is not to introduce new regressivity criteria, but to show that the previously defined covariance-based corrections can be embedded into predictive tools already familiar in mass-appraisal workflows. We focus on two examples that are especially relevant in practice: linear models and second-order tree boosting. The former provides a transparent analytical benchmark that clarifies the algebraic effect of the penalties, while the latter covers widely used high-performing nonlinear learners such as XGBoost- and LightGBM-type methods.

% \subsubsection{Linear Covariance-Penalized Models}
% \label{subsec:linear_models}

% Consider the linear hypothesis class
% \[
% \H_{\mathrm{lin}} := \{f_\theta : x \mapsto x^\top \theta \;:\; \theta \in \Theta \subseteq \R^p\},
% \]
% and let $X \in \R^{n\times p}$ denote the design matrix. For a coefficient vector $\theta \in \Theta$, define
% \[
% f_\theta(x_i) = x_i^\top \theta,
% \qquad
% e := X\theta - y.
% \]
% As in the previous subsection, let
% \[
% H := I - \frac{1}{n}\mathbf{1}\mathbf{1}^\top,
% \qquad
% c := Hy = y - \bar y\,\mathbf{1},
% \qquad
% D := \diag(c_1^2,\dots,c_n^2),
% \]
% so that
% \[
% \widehat{\Cov}_n(e,y)
% =
% \frac{1}{n}e^\top c.
% \]

% Throughout this subsection, we take the baseline loss to be
% \[
% L(f_\theta) = \frac{1}{2n}\|X\theta-y\|_2^2.
% \]

% \paragraph{Squared covariance penalty in linear models.}
% For the symmetric penalty
% \[
% \Psi^{\mathrm{cov}^2}(f_\theta)
% =
% \frac{1}{2}\,\widehat{\Cov}_n(e,y)^2
% =
% \frac{1}{2n^2}\bigl(e^\top c\bigr)^2,
% \]
% the penalized estimator solves
% \begin{align}
%     \label{opt:linear_cov2}
%     \theta_{\rho}^{\mathrm{cov}^2}
%     \in
%     \arg\min_{\theta\in\Theta}
%     \left\{
%         \frac{1}{2n}\|X\theta-y\|_2^2
%         +
%         \frac{\rho}{2n^2}\bigl((X\theta-y)^\top c\bigr)^2
%     \right\}.
% \end{align}
% Define
% \[
% a := X^\top c,
% \qquad
% b := y^\top c = c^\top c.
% \]
% Then the first-order optimality condition for \eqref{opt:linear_cov2} is
% \[
% \left(
% \frac{1}{n}X^\top X + \frac{\rho}{n^2}aa^\top
% \right)\theta
% =
% \frac{1}{n}X^\top y + \frac{\rho}{n^2}ba.
% \]
% Hence, whenever the displayed inverse exists,
% \begin{align}
%     \label{eq:linear_cov2_closed_form}
%     \theta_{\rho}^{\mathrm{cov}^2}
%     =
%     \left(
%         \frac{1}{n}X^\top X
%         +
%         \frac{\rho}{n^2}aa^\top
%     \right)^{-1}
%     \left(
%         \frac{1}{n}X^\top y
%         +
%         \frac{\rho}{n^2}ba
%     \right).
% \end{align}
% This shows that the squared covariance penalty induces a rank-one modification of the usual normal equations.

% \paragraph{Surrogate penalty in linear models.}
% For the separable surrogate
% \[
% \Psi^{\mathrm{surr}}(f_\theta)
% =
% \frac{1}{n}\sum_{i=1}^n e_i^2 c_i^2
% =
% \frac{1}{n}e^\top D\,e,
% \]
% the penalized problem is
% \begin{align}
%     \label{opt:linear_surr}
%     \theta_{\rho}^{\mathrm{surr}}
%     \in
%     \arg\min_{\theta\in\Theta}
%     \left\{
%         \frac{1}{2n}\|X\theta-y\|_2^2
%         +
%         \frac{\rho}{n}e^\top D\,e
%     \right\}.
% \end{align}
% The corresponding optimality condition is
% \[
% \left(
% \frac{1}{n}X^\top X
% +
% \frac{2\rho}{n}X^\top D X
% \right)\theta
% =
% \frac{1}{n}X^\top y
% +
% \frac{2\rho}{n}X^\top D y,
% \]
% and therefore, whenever the displayed inverse exists,
% \begin{align}
%     \label{eq:linear_surr_closed_form}
%     \theta_{\rho}^{\mathrm{surr}}
%     =
%     \left(
%         \frac{1}{n}X^\top X
%         +
%         \frac{2\rho}{n}X^\top D X
%     \right)^{-1}
%     \left(
%         \frac{1}{n}X^\top y
%         +
%         \frac{2\rho}{n}X^\top D y
%     \right).
% \end{align}

% \paragraph{Extension to regularized linear models.}
% The same constructions extend immediately to penalized linear estimators of the form
% \[
% \min_{\theta\in\Theta}
% \left\{
%     L(f_\theta) + \mathcal{R} + \rho\,\Psi(f_\theta)
% \right\},
% \]
% where \(\mathcal{R}\) is a complexity regularizer on the coefficients. In the ridge case, the resulting systems remain quadratic and admit closed-form solutions after adding the usual ridge term to the coefficient matrix. In contrast, \(\ell_1\)-regularized variants such as the LASSO remain convex but no longer admit closed-form expressions, and must instead be solved numerically.

% \paragraph{What the linear case teaches.}
% Linear models clarify the algebraic effect of the penalties and show that the method can be integrated into standard regularized regression with minimal conceptual disruption.

% \subsubsection{Covariance-Penalized Second-Order Tree Boosting}
% \label{subsec:tree_boosting}

% We next consider second-order additive tree boosting under squared loss, as used in gradient-boosted tree systems such as XGBoost and LightGBM. Let \(f^{(m-1)}\) denote the current prediction function at boosting iteration \(m\), and write
% \[
% f_i := f^{(m-1)}(x_i),
% \qquad
% e_i := f_i - y_i,
% \qquad
% c_i := y_i - \bar y,
% \qquad
% z := \widehat{\Cov}_n(e,y) = \frac{1}{n}\sum_{i=1}^n e_i c_i.
% \]
% At step \(m\), one adds a regression tree \(t\) to the current predictor, so that
% \[
% f^{(m)}(x) = f^{(m-1)}(x) + \nu\,t(x),
% \]
% where \(\nu\in(0,1]\) is the learning rate.

% \paragraph{Adding dependence penalties into the boosting objective.}
% For standard second-order boosting, the objective is approximated locally by
% \begin{align}
%     \label{eq:boosting_local_objective}
%     \widetilde{\mathcal J}(t)
%     =
%     \sum_{i=1}^n
%     \left(
%         g_i\,t_i + \frac{1}{2}h_i\,t_i^2
%     \right)
%     +
%     \Omega(t),
% \end{align}
% where \(t_i := t(x_i)\), \(\Omega(t)\) is the tree-complexity penalty, and
% \[
% g_i := \frac{\partial}{\partial f_i}
% \left(
% L(f) + \rho\,\Psi(f)
% \right),
% \qquad
% h_i := \frac{\partial^2}{\partial f_i^2}
% \left(
% L(f) + \rho\,\Psi(f)
% \right)
% \]
% are evaluated at the current predictions. Under squared loss,
% \[
% \frac{\partial L(f)}{\partial f_i} = \frac{e_i}{n},
% \qquad
% \frac{\partial^2 L(f)}{\partial f_i^2} = \frac{1}{n}.
% \]
% The covariance-based penalties modify these quantities as follows.

% \paragraph{Squared covariance penalty in boosting objective.}
% For \(\Psi^{\mathrm{cov}^2}(f)=\tfrac12 z^2\), the gradient is
% \begin{align}
%     \label{eq:boosting_cov2_gradient}
%     g_i^{\mathrm{cov}^2}
%     =
%     \frac{e_i}{n}
%     +
%     \frac{\rho}{n}z\,c_i.
% \end{align}
% The exact Hessian of the penalized objective with respect to the prediction vector \(f=(f_1,\dots,f_n)\) is
% \begin{align}
%     \label{eq:boosting_cov2_exact_hessian}
%     \nabla_f^2
%     \bigl(L(f)+\rho\,\Psi^{\mathrm{cov}^2}(f)\bigr)
%     =
%     \frac{1}{n}I
%     +
%     \frac{\rho}{n^2}cc^\top.
% \end{align}
% Unlike the separable boosting objectives used in standard libraries, \eqref{eq:boosting_cov2_exact_hessian} is dense: each observation is coupled to every other observation through the rank-one term \(cc^\top\). Consequently, there are two natural implementation routes.

% The first is a practical approximation that keeps only the diagonal terms,
% \begin{align}
%     \label{eq:boosting_cov2_diag_hessian}
%     h_i^{\mathrm{cov}^2,\mathrm{diag}}
%     =
%     \frac{1}{n}
%     +
%     \frac{\rho}{n^2}c_i^2.
% \end{align}
% This produces a custom objective compatible with standard second-order boosting APIs, at the price of ignoring the cross-observation curvature induced by the covariance penalty.

% The second is an exact leaf-wise implementation for a fixed tree structure. Suppose a candidate tree \(t\) has leaves \(I_1,\dots,I_J\) and leaf weights \(w=(w_1,\dots,w_J)\), so that \(t_i=w_j\) for \(i\in I_j\). Let
% \[
% n_j := |I_j|,
% \qquad
% b_j := \sum_{i\in I_j} e_i,
% \qquad
% s_j := \sum_{i\in I_j} c_i,
% \]
% and let \(s=(s_1,\dots,s_J)^\top\). If \(\Omega(t)=\gamma J + \frac{\lambda}{2}\|w\|_2^2\), then the exact local objective for the squared covariance penalty can be written as
% \begin{align}
%     \label{eq:boosting_cov2_exact_leaf_objective}
%     \widetilde{\mathcal J}_{\mathrm{cov}^2}(w)
%     =
%     g_T^\top w
%     +
%     \frac{1}{2}w^\top A_T w
%     +
%     \gamma J
%     + \text{const},
% \end{align}
% where
% \[
% g_{T,j}
% =
% \frac{b_j}{n}
% +
% \frac{\rho}{n}z\,s_j,
% \qquad
% A_T
% =
% \diag\!\left(\frac{n_1}{n}+\lambda,\dots,\frac{n_J}{n}+\lambda\right)
% +
% \frac{\rho}{n^2}ss^\top.
% \]
% Hence the optimal leaf weights satisfy
% \begin{align}
%     \label{eq:boosting_cov2_exact_leaf_system}
%     A_T w^\star = -g_T,
%     \qquad
%     w^\star = -A_T^{-1}g_T,
% \end{align}
% whenever \(A_T\) is invertible. The associated gain for the fixed structure \(T\) is
% \begin{align}
%     \label{eq:boosting_cov2_exact_gain}
%     \operatorname{Gain}(T)
%     =
%     \frac{1}{2}g_T^\top A_T^{-1}g_T - \gamma J.
% \end{align}
% This exact route preserves the full curvature of the covariance penalty, but requires a custom tree-construction procedure rather than the usual observation-wise Hessian interface.

% \paragraph{Surrogate penalty as a sample-additive boosting objective.}
% For the surrogate
% \[
% \Psi^{\mathrm{surr}}(f)
% =
% \frac{1}{n}\sum_{i=1}^n e_i^2 c_i^2,
% \]
% the gradient and Hessian remain observation-wise separable:
% \begin{align}
%     \label{eq:boosting_surr_grad_hess}
%     g_i^{\mathrm{surr}}
%     &=
%     \frac{e_i}{n}
%     +
%     \frac{2\rho}{n}e_i c_i^2,
%     \\
%     h_i^{\mathrm{surr}}
%     &=
%     \frac{1}{n}
%     +
%     \frac{2\rho}{n}c_i^2.
% \end{align}
% Accordingly, the surrogate penalty is particularly well suited to off-the-shelf second-order boosting implementations: once the weights \(c_i^2\) are precomputed, the penalized objective can be supplied through standard custom gradient and Hessian routines.

% \subsubsection{Practical implementation remarks}

% The two penalties above lead to two distinct computational regimes in boosted trees. The squared covariance penalty is the most direct implementation of the covariance-control idea, but introduces dense curvature; in practice, it can be handled either through a diagonal Hessian approximation in standard boosting libraries or through an exact custom leaf solver of the form \eqref{eq:boosting_cov2_exact_leaf_system}. The surrogate penalty is fully sample-additive, and is therefore the most convenient to implement through standard observation-wise gradient and Hessian interfaces.

% The formulas in this section show that the covariance-based framework is compatible with both transparent linear estimators and high-performing nonlinear learners. This model-specific flexibility is important for the practical goal of the paper: regressivity correction can be layered onto standard predictive tools already used in assessment workflows, rather than requiring a separate valuation architecture.

% \paragraph{Algorithmic details.}
% For completeness, Appendix~\ref{app:boosting_algorithms} provides pseudocode for an exact Newton-style tree-boosting implementation of the squared covariance penalty, together with an optional first-order approximation. In the main text, we focus on the gradient, Hessian, and leaf-system formulas that define the method and clarify its relation to standard boosted-tree implementations.

% Together, the learner--penalty combinations developed in this section define the candidate model configurations later compared in the constrained calibration procedure of Section~\ref{sec:model_selection}.







\section{Empirical Study: Cook County Assessor's Office}
\label{sec:empirical}

\subsection{Scope, Pipeline, and Evaluation Design}

Our empirical study is based on the residential valuation workflow developed by the Cook County Assessor's Office (CCAO) \cite{ccao_data_2026}. The scope is restricted to \emph{residential} properties in Cook County, Illinois, covering single-family homes, townhomes, small multi-family buildings, and condominiums. For the experiments, we follow the data preparation, preprocessing, and feature-engineering structure of the \texttt{model-res-avm} codebase, and in particular the workflow corresponding to the 2025 assessment cycle. The original workflow is largely implemented in R, whereas our experiments are run in Python~3.10.14 as a translation of that pipeline. The temporal split is chosen to mimic a realistic assessor setting: 2016--2022 are used for model development, 2023 is used as an out-of-time test year, and 2024 is reserved as the current assessment year for later deployment analysis. Unless otherwise stated, the empirical claims in this section are based on the validation and test components of this translated pipeline.

Because the implementation is a translation rather than a direct execution of the original production code, the results below should be interpreted as evidence within the translated CCAO workflow. This distinction matters for reproducibility, but it does not change the main empirical question of the section: whether a covariance-based correction can improve the accuracy--equity tradeoff within a realistic assessor pipeline.

\subsection{Data Description}
\label{subsec:data_description}

\paragraph{Sample construction.} The dataset consists of verified residential property sales in Cook County, drawn from the CCAO's 2025 assessment-cycle data. Following the CCAO pipeline, we exclude multicard PINs (properties with multiple assessment cards) and sales flagged as outliers by the office's verification process. After these filters, the modeling-ready sample contains approximately 349{,}661 observations. Table~\ref{tab:data_splits} reports the temporal splits used throughout the study.

\begin{table}[!ht]
\centering
\caption{Temporal data splits and sample sizes.}
\label{tab:data_splits}
\begin{tabular}{llr}
\toprule
\textbf{Split} & \textbf{Period} & \textbf{Observations} \\
\midrule
Train / Validation  & Jan.\ 2016 -- Dec.\ 2022 & 314{,}694 \\
Test (held-out)     & Jan.\ 2023 -- Dec.\ 2023 & 34{,}967 \\
\midrule
\textit{Total (pre-2024)} & & \textit{349{,}661} \\
\bottomrule
\end{tabular}
\end{table}

\paragraph{Predictors.} The model uses 95 predictor features, of which 23 are treated as categorical. Table~\ref{tab:feature_groups} summarizes the features by category. All features are constructed upstream by the CCAO pipeline; our implementation consumes the already-built feature set without additional feature engineering for LightGBM models.

\begin{table}[!ht]
\centering
\caption{Summary of predictor features by category.}
\label{tab:feature_groups}
\begin{tabular}{lrl}
\toprule
\textbf{Category} & \textbf{Count} & \textbf{Examples} \\
\midrule
Property characteristics & 25 & Year built, bedrooms, baths, building/land sqft \\
Proximity                & 20 & Distance to transit, parks, hospitals, highways \\
ACS Census (5-year)      & 16 & Median income, education, poverty rate \\
Location                 & 8  & Lat/lon, flood factor, walkability, school districts \\
Time                     & 8  & Sale year, quarter, day-of-week, post-COVID indicator \\
Parcel shape             & 6  & Vertex count, edge length SD, bounding rectangle ratios \\
Other / identifiers      & 12 & Township, neighborhood code, tax bill rate \\
\midrule
\textbf{Total}           & \textbf{95} & \\
\bottomrule
\end{tabular}
\end{table}

\paragraph{Validation protocol.} Within the 2016--2022 development sample, we construct 8 rolling-origin validation folds using an expanding-window protocol. Each fold uses a fixed validation block of approximately 31{,}469 observations (10\% of the development sample), while the training window expands chronologically from approximately 35{,}054 to 255{,}337 observations across folds. Validation windows step forward in 9-month increments, covering the period from late 2016 through mid-2022. Table~\ref{tab:fold_structure} reports the fold structure.

\begin{table}[!ht]
\centering
\caption{Rolling-origin validation fold structure. Training windows expand chronologically; validation blocks are of fixed size.}
\label{tab:fold_structure}
\resizebox{0.85\textwidth}{!}{%
\begin{tabular}{c|cc|cc|r}
\toprule
\textbf{Fold} & \textbf{Train start} & \textbf{Train end} & \textbf{Val.\ start} & \textbf{Val.\ end} & \textbf{Train size} \\
\midrule
0 & Jan.\ 2016 & Sep.\ 2016 & Oct.\ 2016 & Jul.\ 2017 & 35{,}054 \\
1 & Jan.\ 2016 & Jul.\ 2017 & Jul.\ 2017 & Mar.\ 2018 & 66{,}523 \\
2 & Jan.\ 2016 & Mar.\ 2018 & Mar.\ 2018 & Nov.\ 2018 & 97{,}992 \\
3 & Jan.\ 2016 & Nov.\ 2018 & Nov.\ 2018 & Aug.\ 2019 & 129{,}461 \\
4 & Jan.\ 2016 & Aug.\ 2019 & Aug.\ 2019 & May 2020 & 160{,}930 \\
5 & Jan.\ 2016 & May 2020 & May 2020 & Jan.\ 2021 & 192{,}399 \\
6 & Jan.\ 2016 & Jan.\ 2021 & Jan.\ 2021 & Aug.\ 2021 & 223{,}868 \\
7 & Jan.\ 2016 & Aug.\ 2021 & Aug.\ 2021 & May 2022 & 255{,}337 \\
\bottomrule
\multicolumn{6}{l}{\footnotesize Validation block size: $\approx$31{,}469 observations per fold.}
\end{tabular}}
\end{table}

\subsection{Baseline Tension: Linear Versus Nonlinear Models}

We first compare two baseline learners used in the CCAO workflow: linear regression and LightGBM. Table~\ref{tab:baseline_models} reports held-out test-set results. LightGBM substantially improves predictive fit and horizontal-equity diagnostics relative to the linear model, raising $R^2$ and lowering MAE, COD, and COV. At the same time, PRD, PRB, and VEI move sharply in the regressive direction. The assessment-level columns confirm the pattern: Linear Regression attains a median ratio near unity ($1.001$), whereas the LightGBM median ratio drops to $0.944$, indicating systematic under-assessment at the typical property value.

This is the central empirical motivation for the paper. A more flexible learner captures market structure better but induces a stronger value-dependent pattern in the assessment ratios. The proposed correction is designed to preserve the predictive advantages of LightGBM while reducing this regressivity.





\begin{table}[!ht]
\centering
\caption{Baseline comparison on the 2023 test set. Lower is better for MAE, COD, and COV. For PRD, median $r$, and mean $r$, values closer to one are preferred; for PRB and VEI, values closer to zero are preferred.}
\label{tab:baseline_models}
\resizebox{0.9\textwidth}{!}{%
\begin{tabular}{l|rr|rr|rr|rrr}
\toprule
Model & $R^2$ & MAE & Med.\ $r$ & Mean $r$ & COD & COV & VEI & PRD & PRB \\
\midrule
Linear Regression  & 0.762 & \$86{,}257 & 1.001 & 1.051 & 25.013\% & 36.8\% & -13.474\% & 1.039 & -0.011 \\
LGBMRegressor      & 0.883 & \$72{,}064 & 0.944 & 1.021 & 23.106\% & 35.4\% & -36.816\% & 1.085 & -0.118 \\
\bottomrule
\end{tabular}}
\end{table}

% \begin{figure}[!ht]
% \centering
% \includegraphics[width=0.72\linewidth]{Fairness Price/GE Paper/img/motivation/Baseline LightGBM_0_motivation_squared_only.pdf}
% \caption{Baseline LightGBM regressivity pattern on the test set. The assessment ratio \(r=\hat P/P\) declines with property value, indicating that lower-valued properties tend to be over-assessed and higher-valued properties tend to be under-assessed.}
% \label{fig:baseline_regressivity}
% \end{figure}

\subsection{Effect of Covariance Penalization}

% We next evaluate whether the proposed covariance-based penalties correct the value-dependent ratio trend induced by the nonlinear baseline. 

The regressive pattern is visible in the ratio--value plots of Figure~\ref{fig:penalty_mechanism}. Under the baseline LightGBM, the fitted ratios exhibit a strong negative trend with linear slope $\approx -0.2156$ and ratio--price correlation $\approx -0.2666$. The squared covariance penalty moves these to $\approx -0.1263$ and $\approx -0.1224$; the surrogate penalty moves them further toward zero, to $\approx -0.0222$ and $\approx -0.0239$. These slope and correlation summaries are not the assessor-facing evaluation criteria, but they confirm that the training-time penalty acts in the intended direction: fitted assessment ratios become less systematically related to property value as $\rho$ grows.

% \begin{figure}[!htbp]
%     \centering
%     \begin{subfigure}{0.5\textwidth}
%         \centering
%         \includegraphics[width=\linewidth, trim={0cm 0mm 0cm 0mm}]{Fairness Price/GE Paper/img/Baseline LightGBM_0_motivation.pdf}
%         \caption{\scriptsize Baseline LightGBM}
%     \end{subfigure}\hfill
%     \begin{subfigure}{0.48\textwidth}
%         \centering
%         \includegraphics[width=\linewidth, trim={0cm 0mm 0cm 0mm}]{Fairness Price/CCAO meeting/img/Covariance-Penalized LightGBM_2_motivation.pdf}
%         \caption{\scriptsize Squared-covariance-penalized LightGBM}
%     \end{subfigure}

%     \caption{Mechanism view of covariance penalization. Relative to the baseline model, the penalized learner visibly flattens the ratio--value trend. In the current experiment, the linear slope moves from approximately \(-0.2156\) to \(-0.1263\), while the ratio--price correlation moves from approximately \(-0.2666\) to \(-0.1224\).}
%     \label{fig:penalty_mechanism}
% \end{figure}

\begin{figure}[!htbp]
    \centering
    
    % --- TOP ROW ---
    \begin{subfigure}{0.48\textwidth}
        \centering
        \includegraphics[width=\linewidth, trim={0cm 0mm 0cm 0mm}]{Fairness Price/CCAO meeting/img/Baseline LightGBM_0_motivation.pdf}
        \caption{\scriptsize Baseline LightGBM}
    \end{subfigure}\hfill
    \begin{subfigure}{0.48\textwidth}
        \centering
        \includegraphics[width=\linewidth, trim={0cm 0mm 0cm 0mm}]{Fairness Price/GE Paper/img/Covariance-Penalized LightGBM_2_motivation.pdf}
        \caption{\scriptsize Squared-covariance-penalized LightGBM}
    \end{subfigure}
    
    \vspace{1em} % Adds a nice gap between the rows
    
    % --- BOTTOM ROW (Centered) ---
    \begin{subfigure}{0.48\textwidth}
        \centering
        % Update the path and filename below for your third plot
        \includegraphics[width=\linewidth, trim={0cm 0mm 0cm 0mm}]{Fairness Price/GE Paper/img/Covariance-Penalized LightGBM_3_motivation.pdf}
        \caption{\scriptsize Surrogate-penalized LightGBM}
    \end{subfigure}

    \caption{Mechanism view of covariance penalization. Relative to the baseline model, the penalized learners visibly flatten the ratio--value trend. In the current experiment, the linear slope moves from approximately \(-0.2156\) to \(-0.1263\), while the ratio--price correlation moves from approximately \(-0.2666\) to \(-0.1224\) for the squared-covariance penalization. For the surrogate penalization, the linear slope moves to \(-0.0222\) and the ratio--price correlation to \(-0.0239\).}
    \label{fig:penalty_mechanism}
\end{figure}

\subsection{Penalized LightGBM Performance on the Test Set}

Table~\ref{tab:penalized_models} reports held-out test results for representative penalized LightGBM configurations. Relative to the unconstrained LightGBM baseline, the penalized models move vertical-equity diagnostics sharply toward assessor-preferred ranges. The best squared-covariance and surrogate configurations remain competitive on predictive accuracy and horizontal equity, and in several cases improve them.

A representative configuration is the surrogate-penalized LightGBM at $\rho=6.572$. Relative to the LightGBM baseline, $R^2$ stays at $0.882$, COD improves from $23.106$ to $22.806$, and COV improves from $35.4\%$ to $33.4\%$. MAE rises modestly from \$72{,}064 to \$74{,}428, a cost of roughly $3.3\%$ in predictive error. Over the same configuration, VEI moves from $-36.8\%$ to $-8.5\%$, PRD from $1.085$ to $1.038$, and PRB from $-0.118$ to $-0.011$. The PRB and VEI reach IAAO-consistent ranges; PRD is close to, but still outside, the conventional $[0.98,1.03]$ tolerance. Larger $\rho$ values (e.g., $\rho=26.519$) push PRD inside the tolerance at the cost of a slightly larger drop in $R^2$, illustrating the equity--accuracy tradeoff the penalty is designed to expose.
\begin{table}[!htbp]
\centering
\caption{Held-out test-set performance for representative penalized LightGBM configurations. The penalty parameter $\rho$ is penalty-specific and should not be compared across different penalty definitions.}
\label{tab:penalized_models}
\resizebox{\textwidth}{!}{%
\begin{tabular}{lc|cc|cc|cc|ccc}
\toprule
Model Family & $\rho$ & $R^2$ & MAE & Med.\ $r$ & Mean $r$ & COD & COV & VEI & PRD & PRB \\
\midrule
Linear Regression  & -- & 0.762 & \$86{,}256.879 & 1.001 & 1.051 & 25.013\% & 36.825\% & -13.474\% & 1.039 & -0.011 \\
LGBMRegressor      & -- & 0.883 & \$72{,}063.853 & 0.944 & 1.021 & 23.106\% & 35.398\% & -36.816\% & 1.085 & -0.118 \\
\midrule
LGBSurrogatePenalty & 0.100 & 0.881 & \$73{,}271.882 & 0.936 & 1.011 & 23.125\% & 35.236\% & -36.367\% & 1.084 & -0.115 \\
LGBSurrogatePenalty & 0.189 & 0.881 & \$73{,}208.883 & 0.936 & 1.010 & 23.042\% & 35.048\% & -33.608\% & 1.080 & -0.107 \\
LGBSurrogatePenalty & 0.404 & 0.882 & \$73{,}110.847 & 0.937 & 1.006 & 22.924\% & 34.782\% & -29.878\% & 1.075 & -0.095 \\
LGBSurrogatePenalty & 0.761 & 0.883 & \$72{,}974.386 & 0.937 & 1.001 & 22.749\% & 34.189\% & -23.789\% & 1.067 & -0.076 \\
LGBSurrogatePenalty & 1.628 & 0.882 & \$73{,}233.014 & 0.938 & 0.994 & 22.645\% & 33.664\% & -16.431\% & 1.057 & -0.052 \\
LGBSurrogatePenalty & 3.070 & 0.881 & \$73{,}663.987 & 0.937 & 0.986 & 22.667\% & 33.620\% & -12.630\% & 1.048 & -0.032 \\
LGBSurrogatePenalty & 6.572 & 0.882 & \$74{,}428.114 & 0.937 & 0.976 & 22.806\% & 33.386\% & -8.545\% & 1.038 & -0.011 \\
LGBSurrogatePenalty & 12.390 & 0.879 & \$75{,}846.355 & 0.935 & 0.968 & 23.249\% & 33.684\% & -7.446\% & 1.032 & 0.003 \\
LGBSurrogatePenalty & 26.519 & 0.879 & \$77{,}113.783 & 0.933 & 0.960 & 23.719\% & 33.959\% & -4.484\% & 1.024 & 0.020 \\
LGBSurrogatePenalty & 50.000 & 0.878 & \$78{,}155.478 & 0.932 & 0.956 & 24.092\% & 34.336\% & -4.883\% & 1.021 & 0.027 \\
\midrule
LGBCovPenalty & 0.100 & 0.879 & \$73{,}459.227 & 0.937 & 1.014 & 23.210\% & 35.551\% & -38.368\% & 1.087 & -0.122 \\
LGBCovPenalty & 0.189 & 0.880 & \$73{,}311.803 & 0.939 & 1.015 & 23.221\% & 35.443\% & -37.671\% & 1.086 & -0.119 \\
LGBCovPenalty & 0.404 & 0.882 & \$72{,}947.453 & 0.941 & 1.017 & 23.224\% & 35.432\% & -36.501\% & 1.083 & -0.115 \\
LGBCovPenalty & 0.761 & 0.883 & \$72{,}607.948 & 0.944 & 1.019 & 23.157\% & 35.246\% & -34.017\% & 1.078 & -0.106 \\
LGBCovPenalty & 1.628 & 0.886 & \$72{,}198.279 & 0.949 & 1.023 & 23.176\% & 35.086\% & -30.427\% & 1.071 & -0.096 \\
LGBCovPenalty & 3.070 & 0.884 & \$72{,}293.558 & 0.955 & 1.027 & 23.216\% & 35.053\% & -26.850\% & 1.065 & -0.085 \\
LGBCovPenalty & 6.572 & 0.881 & \$72{,}871.805 & 0.961 & 1.030 & 23.282\% & 35.040\% & -20.525\% & 1.055 & -0.068 \\
LGBCovPenalty & 12.390 & 0.878 & \$73{,}449.207 & 0.963 & 1.030 & 23.261\% & 34.757\% & -16.050\% & 1.048 & -0.057 \\
LGBCovPenalty & 26.519 & 0.870 & \$74{,}741.535 & 0.963 & 1.028 & 23.310\% & 34.542\% & -13.572\% & 1.041 & -0.046 \\
LGBCovPenalty & 50.000 & 0.866 & \$75{,}574.792 & 0.960 & 1.023 & 23.357\% & 34.522\% & -11.105\% & 1.037 & -0.039 \\
\bottomrule
\end{tabular}}
\end{table}

% \begin{table}[!htbp]
% \centering
% \caption{Held-out test-set performance for representative penalized LightGBM configurations. The penalty parameter $\rho$ is penalty-specific and should not be compared across different penalty definitions.}
% \label{tab:penalized_models}
% \resizebox{\textwidth}{!}{%
% \begin{tabular}{lc|cc|cc|ccc}
% \toprule
% Model & $\rho$ & $R^2$ & MAE & COD & COV & VEI & PRD & PRB \\
% \midrule
% Linear Regression  & -- & 0.762 & \$86{,}257 & 25.013\% & 36.8\% & -13.474\% & 1.039 & -0.011 \\
% LGBMRegressor      & -- & 0.883 & \$72{,}064 & 23.106\% & 35.4\% & -36.816\% & 1.085 & -0.118 \\
% \midrule
% LGBCovPenalty & 0.100 & 0.879 & \$73{,}459 & 23.210\% & 35.6\% & -38.368\% & 1.087 & -0.122 \\
% LGBCovPenalty & 0.189 & 0.880 & \$73{,}312 & 23.221\% & 35.4\% & -37.671\% & 1.086 & -0.119 \\
% LGBCovPenalty & 0.404 & 0.882 & \$72{,}947 & 23.224\% & 35.4\% & -36.501\% & 1.083 & -0.115 \\
% LGBCovPenalty & 0.761 & 0.883 & \$72{,}608 & 23.157\% & 35.2\% & -34.017\% & 1.078 & -0.106 \\
% LGBCovPenalty & 1.628 & 0.886 & \$72{,}198 & 23.176\% & 35.1\% & -30.427\% & 1.071 & -0.096 \\
% LGBCovPenalty & 3.070 & 0.884 & \$72{,}294 & 23.216\% & 35.1\% & -26.850\% & 1.065 & -0.085 \\
% LGBCovPenalty & 6.572 & 0.881 & \$72{,}872 & 23.282\% & 35.0\% & -20.525\% & 1.055 & -0.068 \\
% LGBCovPenalty & 12.390 & 0.878 & \$73{,}449 & 23.261\% & 34.8\% & -16.050\% & 1.048 & -0.057 \\
% LGBCovPenalty & 26.519 & 0.870 & \$74{,}742 & 23.310\% & 34.5\% & -13.572\% & 1.041 & -0.046 \\
% LGBCovPenalty & 50.000 & 0.866 & \$75{,}575 & 23.357\% & 34.5\% & -11.105\% & 1.037 & -0.039 \\
% \midrule
% LGBSurrogatePenalty & 0.100 & 0.881 & \$73{,}272 & 23.125\% & 35.2\% & -36.367\% & 1.084 & -0.115 \\
% LGBSurrogatePenalty & 0.189 & 0.881 & \$73{,}209 & 23.042\% & 35.0\% & -33.608\% & 1.080 & -0.107 \\
% LGBSurrogatePenalty & 0.404 & 0.882 & \$73{,}111 & 22.924\% & 34.8\% & -29.878\% & 1.075 & -0.095 \\
% LGBSurrogatePenalty & 0.761 & 0.883 & \$72{,}974 & 22.749\% & 34.2\% & -23.789\% & 1.067 & -0.076 \\
% LGBSurrogatePenalty & 1.628 & 0.882 & \$73{,}233 & 22.645\% & 33.7\% & -16.431\% & 1.057 & -0.052 \\
% LGBSurrogatePenalty & 3.070 & 0.881 & \$73{,}664 & 22.667\% & 33.6\% & -12.630\% & 1.048 & -0.032 \\
% LGBSurrogatePenalty & 6.572 & 0.882 & \$74{,}428 & 22.806\% & 33.4\% & -8.545\% & 1.038 & -0.011 \\
% LGBSurrogatePenalty & 12.390 & 0.879 & \$75{,}846 & 23.249\% & 33.7\% & -7.446\% & 1.032 & 0.003 \\
% LGBSurrogatePenalty & 26.519 & 0.879 & \$77{,}114 & 23.719\% & 34.0\% & -4.484\% & 1.024 & 0.020 \\
% LGBSurrogatePenalty & 50.000 & 0.878 & \$78{,}155 & 24.092\% & 34.3\% & -4.883\% & 1.021 & 0.027 \\
% \bottomrule
% \end{tabular}}
% \end{table}


\subsection{Accuracy--Equity Tradeoffs}

Increasing $\rho$ traces out a nontrivial accuracy--equity frontier rather than a steep linear tradeoff. In PRD--$R^2$ space, the penalized models move away from the unconstrained LightGBM point toward the assessor compliance region while remaining close to the upper edge of the accuracy frontier. The same pattern holds in VEI--$R^2$ space. This is the empirical justification for treating calibration as a constrained model-selection problem: once multiple diagnostics matter simultaneously, the operational question is not which model ranks best on a single score, but which model delivers the strongest predictive performance among those that satisfy explicit ratio-study tolerances.

\begin{figure}[!htbp]
    \centering
    \begin{subfigure}[b]{0.48\textwidth}
        \centering
        \includegraphics[width=\textwidth, trim={2mm 0 4mm 2mm}]{Fairness Price/GE Paper/img/tradeoffs/PRD_vs_R2_band.pdf}
        \caption{PRD versus \(R^2\)}
        \label{fig:prd_tradeoff}
    \end{subfigure}
    \hfill
    \begin{subfigure}[b]{0.48\textwidth}
        \centering
        \includegraphics[width=\textwidth, trim={4mm 0 0 2mm}]{Fairness Price/GE Paper/img/tradeoffs/VEI_vs_R2_band.pdf}
        \caption{VEI versus \(R^2\)}
        \label{fig:vei_tradeoff}
    \end{subfigure}

    \caption{Accuracy--equity tradeoffs induced by covariance penalization on the test set. As \(\rho\) increases, the penalized models move away from the unconstrained LightGBM point toward assessor-preferred vertical-equity ranges while remaining near the upper part of the accuracy frontier.}
    \label{fig:pareto_frontier}. 
\end{figure}





\subsection{Multi-Objective Model Selection Across Folds}
\label{subsec:model_selection}



Once the implemented model families are fit (with or without covariance-based penalization), a final candidate must be chosen for deployment. Because different configurations achieve different tradeoffs across predictive error, vertical-equity diagnostics, and dispersion, calibration is naturally formulated as a constrained multi-objective selection problem over the candidate set.

% This is illustrated in Figure \label{fig:tradeoffs} \todo{specify?}

% \begin{figure}[!htp]
%     \centering
%     \includegraphics[width=0.5\linewidth]{Fairness Price/GE Paper/img/model-selection/extreme_case_3_constraint_O1_0.90_edit.pdf}
%     \caption{Caption}
%     \label{fig:tradeoffs}
% \end{figure}


\subsubsection{Multi-Objective Model-Selection Setting}


% \subsection{Temporal Evaluation and Validation Setup}
% \label{subsec:oos_protocol}

% Mass appraisal is a temporally-deployed prediction problem: models are trained on historical property sales, and then used to predict future market values. Hence, out-of-sample performance is chronological rather than purely random-split based. We use rolling-origin is standard, time ordered, evaluation and calibration setting \cite{tashman2000out}.

% Let the full sold-data sample be ordered by sale date and partitioned into two chronologically separated parts:
% $
% \mathcal D^{\mathrm{dev}}
% \text{ and }
% \mathcal D^{\mathrm{test}},
% $
% where $\mathcal D^{\mathrm{dev}}$ is the \emph{development} sample used for training and calibration, and $\mathcal D^{\mathrm{test}}$ is a final holdout period used only for one-time reporting after model selection is complete.

% Within the development sample $\mathcal D^{\mathrm{dev}}$, we construct a sequence of rolling-origin windows
% $
% \{(\mathcal T_v,\mathcal V_v)\}_{v=1}^V,
% $
% where each training set $\mathcal T_v$ precedes its corresponding validation set $\mathcal V_v$ in time. The windows may be expanding or otherwise office-specific, but they always satisfy the chronological requirement
% $
% \max_{i\in \mathcal T_v} t_i
% <
% \min_{j\in \mathcal V_v} t_j.
% $

% For each validation window, predictive performance may be evaluated on either the log-price or price scale, depending on the metric of interest. Equity diagnostics, however, are always computed after back-transforming predictions to the original price scale and forming the assessment ratios in \eqref{eq:ratio_def}. This distinction is maintained throughout the paper: models are estimated on log prices, but assessor-facing fairness is judged on price-scale ratios.

% These rolling-origin validation windows provide the out-of-time predictions used in Section\todo{reference} for constrained calibration and for final evaluation on the holdout test period.





% \paragraph{Single-model calibration.}

Let $\mathcal V=\{1,\dots,V\}$ index the validation folds (or rolling-origin validation windows), and let $V_v$ denote the validation sample associated with fold $v\in\mathcal V$. Let $\mathcal C$ denote a finite set of candidate model configurations. Each configuration $c\in\mathcal C$ specifies
(i) a model family from the implemented classes,
(ii) a penalty choice in
$
\Psi \in \{0,\Psi^{\mathrm{cov}^2},\Psi^{\mathrm{surr}}\},
$
where $\Psi=0$ denotes the unpenalized baseline,
and
(iii) all associated tuning parameters, including the penalty level $\rho\ge 0$ when applicable.

For each fold $v\in\mathcal V$, fitting configuration $c$ on the corresponding training portion yields a predictor $f_{c,v}$. For each observation $i\in V_v$, let
$
\hat y_i^{\,v}(c):=f_{c,v}(x_i^{\,v})$, $
\hat P_i^{\,v}(c):=\exp\!\bigl(\hat y_i^{\,v}(c)\bigr),$
\text{ and }
$r_i^{\,v}(c):=\frac{\hat P_i^{\,v}(c)}{P_i^{\,v}},
$
where $P_i^{\,v}$ is the observed sale price. 

From the validation predictions on fold $v$, we compute a predictive performance error score $\mathrm{E}_v(c)$ (e.g., mean-squared-error), together with a collection of fold-level assessment diagnostics $Q_{m,v}(c)$ indexed by $m$ (e.g., PRD, VEI, and COD). 

% If one starts from an error metric $E_v(c)$ that is minimized, such as MSE, RMSE, or MAE, one may simply set $\mathrm{Perf}_v(c):=-E_v(c)$.

To aggregate results across folds, let $\alpha_v\ge 0$ satisfy $\sum_{v\in\mathcal V}\alpha_v=1$, and define
\[
\bar{\mathrm{Perf}}(c):=\sum_{v\in\mathcal V}\alpha_v\,\mathrm{Perf}_v(c),
\qquad
\bar Q_m(c):=\sum_{v\in\mathcal V}\alpha_v\,Q_{m,v}(c).
\]
The default choice is uniform weighting, $\alpha_v=1/V$, although the same formulation also allows non-uniform weights when one wishes to emphasize, for example, more recent validation windows.


\paragraph{Model Selection Formulation.} 
We select a single deployed model by solving
\begin{align}
    \label{opt:model_selection_single}
    c^\star
    \in
    \arg\max_{c\in\mathcal C}
    \quad
    & \bar{\mathrm{Perf}}(c)
    \\
    \text{s.t.}\quad
    & \underline{\kappa}_m \le \bar Q_m(c) \le \overline{\kappa}_m,
    \qquad \forall m\in\mathcal M_{\mathrm{int}},
    \nonumber
    % & \bar Q_m(c) \le \overline{\kappa}_m,
    % \qquad \forall m\in\mathcal M_{\mathrm{up}}.
    % \nonumber
\end{align}
where $\mathcal M$ indexes diagnostics constrained to lie in admissible intervals.

Formulation \eqref{opt:model_selection_single} avoids collapsing the validation criteria into a one-dimensional weighted score: predictive performance is the objective, while ratio-study, regressivity, and dispersion diagnostics enter as feasibility constraints. This is the direct analogue of the compliance logic assessors already apply to candidate models, under which the selected model is the one with the best predictive performance among those that satisfy the chosen ratio-study tolerances.

% \paragraph{How the constrained calibration differs from standard accuracy-only selection.}
% The single-model procedure can be read as the following practical recipe.
% \begin{enumerate}
%     \item Specify a finite candidate set $\mathcal C$ of model configurations, including baseline and covariance-penalized variants.
%     \item For each $c\in\mathcal C$ and each validation fold $v\in\mathcal V$, fit the model on the corresponding training portion and compute validation predictions.
%     \item On each fold, compute the chosen predictive score $\mathrm{Perf}_v(c)$ together with the selected assessment diagnostics $Q_{m,v}(c)$.
%     \item Aggregate these fold-level quantities into $\bar{\mathrm{Perf}}(c)$ and $\bar Q_m(c)$.
%     \item Discard all configurations that violate any active assessment-quality constraint.
%     \item Rank the remaining feasible configurations by aggregated predictive performance and select the best one.
%     \item Refit the selected configuration $c^\star$ on the full pre-assessment development sample and deploy that final predictor.
% \end{enumerate}
% The only substantive departure from the usual accuracy-only workflow lies in Steps 5--6: instead of ranking all candidates solely by predictive error, the constrained version first screens for feasibility with respect to the chosen assessor-facing diagnostics and then ranks only the feasible candidates by predictive performance.

The formulation above uses aggregated fold-level diagnostics as the default choice. When temporal stability is especially important, one may replace some aggregated constraints with foldwise requirements of the form
\[
Q_{m,v}(c)\in[\underline{\kappa}_m,\overline{\kappa}_m]
\qquad
\forall v\in\mathcal V, \forall m \in \mathcal M
\]
or, more generally, use a more conservative fold aggregator. We keep the weighted-average formulation as the main presentation because it yields a transparent calibration problem over a finite candidate set and aligns naturally with standard validation practice.


\subsubsection{Alternative Pareto Criteria and a Stacking Extension}
\label{subsubsec:pareto_alternatives}

The constrained formulation in \eqref{opt:model_selection_single} has three known limitations: the constraints are enforced on validation data and need not hold out-of-sample, they may be jointly infeasible over the candidate set, and the predictive cost of tight tolerances can be large. To cover these cases we also consider two Pareto-based selection rules that do not require explicit tolerances: a \emph{closest-to-Utopian} rule that picks the configuration closest (after normalization) to the ideal of the best attained value on each metric, and a \emph{Nash equilibrium / maximum-volume} rule that picks the configuration whose aggregated-metric vector sacrifices proportionally the least on each dimension. Figure~\ref{fig:pareto_solutions} illustrates the three rules on a two-metric example. We treat these rules as supplements to \eqref{opt:model_selection_single} rather than as coequal contributions; both are fully specified, with the convex-stacking extension, in Appendix~\ref{app:pareto_and_stacking}.

\begin{figure}[!htp]
    \centering
    \includegraphics[width=0.7\linewidth]{Fairness Price/GE Paper/img/model-selection/extreme_case_3_constraint_O1_0.90_edit.pdf}
    \caption{The three Pareto-optimal solutions on a two-objective convex frontier: constrained, closest-to-Utopian, and Nash / maximum-volume. The constrained rule is threshold-dependent; the other two do not require explicit tolerances.}
    \label{fig:pareto_solutions}
\end{figure}

\paragraph{Stacking extension.}
A natural further extension of \eqref{opt:model_selection_single} replaces discrete candidate selection by a convex combination of candidate predictions, which preserves the objective--constraint logic and can only enlarge the attainable accuracy--equity frontier. Because the assessor diagnostics are defined in terms of price-scale ratios, stacking is carried out on the price scale, and several ratio-study constraints admit an exact linear reformulation. The full stacking formulation and its convex-quadratic special case are deferred to Appendix~\ref{app:pareto_and_stacking}.



\subsubsection{Empirical Calibration Results}

Table~\ref{tab:test_results} reports the held-out performance of models selected under different subsets of metrics. The constrained rule uses the intervals in Table~\ref{tab:assessment_metrics_summary}. All selected configurations substantially improve regressivity, but the constrained versions become infeasible once non-regressivity metrics are added to the constraint set. This is consistent with the fact that our penalties target value-dependent trends rather than dispersion or central-tendency diagnostics, and it suggests that extending the penalty family to cover additional assessor-facing quantities is a natural next step.
% \subsubsection{Constrained Single-Model Selection}

\begin{table}[!htbp]
\centering
\caption{Held-out test set performance. The penalty parameter $\rho$ is specific to the penalty definition.}
\label{tab:test_results}
\resizebox{\textwidth}{!}{%
\begin{tabular}{llc|cc|cc|cc|ccc}
\toprule
Solution & Selected model & $\rho$ & $R^2$ & MAE & Med.\ $r$ & Mean $r$ & COD & COV & PRD & PRB & VEI \\
\midrule
\multicolumn{12}{l}{\textbf{Unconstrained baselines}} \\
\midrule
BASELINE & Linear Reg. & -- & 0.762 & \$86{,}244 & 1.001 & 1.050 & 25.01\% & 36.86\% & 1.039 & -0.011 & -13.53\% \\
BASELINE & LGBM & -- & 0.885 & \$71{,}912 & 0.944 & 1.021 & 23.07\% & 35.48\% & 1.085 & -0.118 & -37.64\% \\
\midrule
\multicolumn{12}{l}{\textbf{Constraints: PRD}} \\
\midrule
CONSTRAINED & LGBSurrogate & 5.96 & 0.881 & \$74{,}416 & 0.937 & 0.978 & 22.80\% & 33.64\% & 1.039 & -0.014 & -9.27\% \\
UTOPIA & LGBSurrogate & 100.00 & 0.877 & \$79{,}217 & 0.930 & 0.953 & 24.47\% & 34.89\% & 1.019 & 0.032 & -4.41\% \\
NASH & LGBSurrogate & 2.95 & 0.883 & \$73{,}395 & 0.938 & 0.988 & 22.63\% & 33.63\% & 1.048 & -0.034 & -12.20\% \\
\midrule
\multicolumn{12}{l}{\textbf{Constraints: PRD, PRB, and VEI}} \\
\midrule
CONSTRAINED & LGBSurrogate & 5.96 & 0.881 & \$74{,}416 & 0.937 & 0.978 & 22.80\% & 33.64\% & 1.039 & -0.014 & -9.27\% \\
UTOPIA & LGBSurrogate & 7.91 & 0.882 & \$74{,}939 & 0.935 & 0.974 & 23.02\% & 33.85\% & 1.036 & -0.007 & -8.36\% \\
NASH & LGBSurrogate & 5.96 & 0.881 & \$74{,}416 & 0.937 & 0.978 & 22.80\% & 33.64\% & 1.039 & -0.014 & -9.27\% \\
\midrule
\multicolumn{12}{l}{\textbf{Constraints: PRD, Mean $r$, and COD}} \\
\midrule
INFEASIBLE* & --- & --- & --- & --- & --- & --- & --- & --- & --- & --- & --- \\
UTOPIA & LGBSurrogate & 6.87 & 0.882 & \$74{,}449 & 0.937 & 0.977 & 22.84\% & 33.72\% & 1.038 & -0.011 & -9.38\% \\
NASH & LGBSurrogate & 2.95 & 0.883 & \$73{,}395 & 0.938 & 0.988 & 22.63\% & 33.63\% & 1.048 & -0.034 & -12.20\% \\
\midrule
\multicolumn{12}{l}{\textbf{Constraints: Full metric set$^\dagger$}} \\
\midrule
INFEASIBLE* & --- & --- & --- & --- & --- & --- & --- & --- & --- & --- & --- \\
UTOPIA & LGBSurrogate & 5.96 & 0.881 & \$74{,}416 & 0.937 & 0.978 & 22.80\% & 33.64\% & 1.039 & -0.014 & -9.27\% \\
NASH & LGBSurrogate & 5.96 & 0.881 & \$74{,}416 & 0.937 & 0.978 & 22.80\% & 33.64\% & 1.039 & -0.014 & -9.27\% \\
\bottomrule
\multicolumn{12}{l}{\footnotesize *INFEASIBLE abbreviates CONSTRAINED INFEASIBLE; no feasible selected solution.} \\
\multicolumn{12}{l}{\footnotesize $^\dagger$ PRD, PRB, VEI, COD, COV, Mean $r$, Median $r$, and Weighted Mean $r$.}
\end{tabular}}
\end{table}



% \begin{table}[!htbp]
% \centering
% \caption{Held-out test-set performance expressed as signed percentage change relative to the LGBM baseline. The penalty parameter $\rho$ is specific to the penalty definition.}
% \label{tab:test_results_pct_change}
% \resizebox{\textwidth}{!}{%
% \begin{tabular}{llc|cc|cc|cc|ccc}
% \toprule
% Solution & Selected model & $\rho$ & $R^2$ & MAE & Med.\ $r$ & Mean $r$ & COD & COV & PRD & PRB & VEI \\
% \midrule
% \multicolumn{12}{l}{\textbf{Unconstrained baselines}} \\
% \midrule
% BASELINE & Linear Reg. & -- & -13.9\% & +19.9\% & +6.0\% & +2.9\% & +8.4\% & +3.9\% & -4.3\% & -91.1\% & -64.1\% \\
% BASELINE & LGBM & -- & +0.0\% & +0.0\% & +0.0\% & +0.0\% & +0.0\% & +0.0\% & +0.0\% & +0.0\% & +0.0\% \\
% \midrule
% \multicolumn{12}{l}{\textbf{Constraints: PRD}} \\
% \midrule
% CONSTRAINED & LGBSurrogate & 5.96 & -0.4\% & +3.5\% & -0.8\% & -4.2\% & -1.2\% & -5.2\% & -4.2\% & -88.2\% & -75.4\% \\
% UTOPIA & LGBSurrogate & 100.00 & -0.9\% & +10.2\% & -1.5\% & -6.7\% & +6.1\% & -1.7\% & -6.1\% & -126.9\% & -88.3\% \\
% NASH & LGBSurrogate & 2.95 & -0.3\% & +2.1\% & -0.6\% & -3.2\% & -1.9\% & -5.2\% & -3.4\% & -71.0\% & -67.6\% \\
% \midrule
% \multicolumn{12}{l}{\textbf{Constraints: PRD, PRB, and VEI}} \\
% \midrule
% CONSTRAINED & LGBSurrogate & 5.96 & -0.4\% & +3.5\% & -0.8\% & -4.2\% & -1.2\% & -5.2\% & -4.2\% & -88.2\% & -75.4\% \\
% UTOPIA & LGBSurrogate & 7.91 & -0.4\% & +4.2\% & -1.0\% & -4.6\% & -0.2\% & -4.6\% & -4.5\% & -93.9\% & -77.8\% \\
% NASH & LGBSurrogate & 5.96 & -0.4\% & +3.5\% & -0.8\% & -4.2\% & -1.2\% & -5.2\% & -4.2\% & -88.2\% & -75.4\% \\
% \midrule
% \multicolumn{12}{l}{\textbf{Constraints: PRD, Mean $r$, and COD}} \\
% \midrule
% INFEASIBLE* & --- & --- & --- & --- & --- & --- & --- & --- & --- & --- & --- \\
% UTOPIA & LGBSurrogate & 6.87 & -0.3\% & +3.5\% & -0.8\% & -4.3\% & -1.0\% & -5.0\% & -4.3\% & -90.5\% & -75.1\% \\
% NASH & LGBSurrogate & 2.95 & -0.3\% & +2.1\% & -0.6\% & -3.2\% & -1.9\% & -5.2\% & -3.4\% & -71.0\% & -67.6\% \\
% \midrule
% \multicolumn{12}{l}{\textbf{Constraints: Full metric set$^\dagger$}} \\
% \midrule
% INFEASIBLE* & --- & --- & --- & --- & --- & --- & --- & --- & --- & --- & --- \\
% UTOPIA & LGBSurrogate & 5.96 & -0.4\% & +3.5\% & -0.8\% & -4.2\% & -1.2\% & -5.2\% & -4.2\% & -88.2\% & -75.4\% \\
% NASH & LGBSurrogate & 5.96 & -0.4\% & +3.5\% & -0.8\% & -4.2\% & -1.2\% & -5.2\% & -4.2\% & -88.2\% & -75.4\% \\
% \bottomrule
% \multicolumn{12}{l}{\footnotesize *INFEASIBLE abbreviates CONSTRAINED INFEASIBLE; no feasible selected solution.} \\
% \multicolumn{12}{l}{\footnotesize $^\dagger$ PRD, PRB, VEI, COD, COV, Mean $r$, Median $r$, and Weighted Mean $r$.} \\
% \multicolumn{12}{l}{\footnotesize Percent changes are computed as $100(x/x_{\mathrm{LGBM}}-1)$. For signed metrics such as PRB and VEI,} \\
% \multicolumn{12}{l}{\footnotesize the sign reflects the standard relative-change formula and may look counterintuitive when values move closer to zero.}
% \end{tabular}}
% \end{table}
% % \begin{table}[!htbp]
% % \centering
% % \caption{Held-out test set performance. The penalty parameter $\rho$ is specific to the penalty definition.}
% % \label{tab:test_results}
% % \resizebox{\textwidth}{!}{%
% % \begin{tabular}{llc|cc|cc|cc|ccc}
% % \toprule
% % Family & Solution & $\rho$ & $R^2$ & MAE & Med $r$ & Mean $r$ & COD & COV & VEI & PRD & PRB \\
% % \midrule
% % \multicolumn{12}{l}{\textbf{Baselines}} \\
% % \midrule
% % Linear Reg. & BASELINE & -- & 0.762 & \$86{,}244 & 1.001 & 1.050 & 25.01\% & 36.86\% & -13.53\% & 1.039 & -0.011 \\
% % LGBM & BASELINE & -- & 0.885 & \$71{,}912 & 0.944 & 1.021 & 23.07\% & 35.48\% & -37.64\% & 1.085 & -0.118 \\
% % \midrule
% % \multicolumn{12}{l}{\textbf{Constraints: PRD}} \\
% % \midrule
% % LGBSurrogate & CONSTRAINED & 5.96 & 0.881 & \$74{,}416 & 0.937 & 0.978 & 22.80\% & 33.64\% & -9.27\% & 1.039 & -0.014 \\
% % LGBSurrogate & UTOPIA & 100.00 & 0.877 & \$79{,}217 & 0.930 & 0.953 & 24.47\% & 34.89\% & -4.41\% & 1.019 & 0.032 \\
% % LGBSurrogate & NASH & 2.95 & 0.883 & \$73{,}395 & 0.938 & 0.988 & 22.63\% & 33.63\% & -12.20\% & 1.048 & -0.034 \\
% % \midrule
% % \multicolumn{12}{l}{\textbf{Constraints: PRD, PRB, and VEI}} \\
% % \midrule
% % LGBSurrogate & CONSTRAINED & 5.96 & 0.881 & \$74{,}416 & 0.937 & 0.978 & 22.80\% & 33.64\% & -9.27\% & 1.039 & -0.014 \\
% % LGBSurrogate & UTOPIA & 7.91 & 0.882 & \$74{,}939 & 0.935 & 0.974 & 23.02\% & 33.85\% & -8.36\% & 1.036 & -0.007 \\
% % LGBSurrogate & NASH & 5.96 & 0.881 & \$74{,}416 & 0.937 & 0.978 & 22.80\% & 33.64\% & -9.27\% & 1.039 & -0.014 \\
% % \midrule
% % \multicolumn{12}{l}{\textbf{Constraints: PRD, MEAN\_RATIO, and COD}} \\
% % \midrule
% % LGBSurrogate & INFEASIBLE* & 5.96 & 0.881 & \$74{,}416 & 0.937 & 0.978 & 22.80\% & 33.64\% & -9.27\% & 1.039 & -0.014 \\
% % LGBSurrogate & UTOPIA & 6.87 & 0.882 & \$74{,}449 & 0.937 & 0.977 & 22.84\% & 33.72\% & -9.38\% & 1.038 & -0.011 \\
% % LGBSurrogate & NASH & 2.95 & 0.883 & \$73{,}395 & 0.938 & 0.988 & 22.63\% & 33.63\% & -12.20\% & 1.048 & -0.034 \\
% % \midrule
% % \multicolumn{12}{l}{\textbf{Constraints: All listed metrics$^\dagger$}} \\
% % \midrule
% % LGBSurrogate & INFEASIBLE* & 5.96 & 0.881 & \$74{,}416 & 0.937 & 0.978 & 22.80\% & 33.64\% & -9.27\% & 1.039 & -0.014 \\
% % LGBSurrogate & UTOPIA & 5.96 & 0.881 & \$74{,}416 & 0.937 & 0.978 & 22.80\% & 33.64\% & -9.27\% & 1.039 & -0.014 \\
% % LGBSurrogate & NASH & 5.96 & 0.881 & \$74{,}416 & 0.937 & 0.978 & 22.80\% & 33.64\% & -9.27\% & 1.039 & -0.014 \\
% % \bottomrule
% % \multicolumn{12}{l}{\footnotesize *INFEASIBLE abbreviates CONSTRAINED\_CLOSEST\_INFEASIBLE.} \\
% % \multicolumn{12}{l}{\footnotesize $^\dagger$ PRD, PRB, VEI, COD, COV, MEAN\_RATIO, MEDIAN\_RATIO, and WEIGHTED\_MEAN\_RATIO.}
% % \end{tabular}}
% % \end{table}


% \paragraph{Constraining PRD}
% \begin{table}[!htbp]
% \centering
% \caption{Average cross-validation performance. The penalty parameter $\rho$ is specific to the penalty definition.}
% \label{tab:val_results}
% \resizebox{\textwidth}{!}{%
% \begin{tabular}{llc|cc|cc|ccc}
% \toprule
% Family & Solution & $\rho$ & $R^2$ & MAE & COD & COV & VEI & PRD & PRB \\
% \midrule
% Linear Regression & BASELINE & -- & 0.667 & \$92{,}441.063 & 32.765\% & 49.185\% & 8.633\% & 1.026 & 0.061 \\
% LGBMRegressor & BASELINE & -- & 0.890 & \$58{,}097.609 & 22.990\% & 37.623\% & -24.247\% & 1.069 & -0.077 \\
% \midrule
% LGBSurrogatePenalty & CONSTRAINED & 5.789 & 0.886 & \$61{,}122.635 & 23.124\% & 35.511\% & -3.446\% & 1.026 & -0.001 \\
% LGBSurrogatePenalty & UTOPIA & 12.390 & 0.885 & \$61{,}769.215 & 23.352\% & 35.024\% & -2.065\% & 1.019 & 0.011 \\
% LGBSurrogatePenalty & NASH & 50.000 & 0.882 & \$63{,}662.712 & 24.188\% & 35.539\% & -0.101\% & 1.009 & 0.028 \\
% \midrule
% LGBCovPenalty & CONSTRAINED\_CLOSEST\_INFEASIBLE & 44.044 & 0.853 & \$68{,}940.597 & 24.902\% & 39.202\% & 2.768\% & 1.023 & -0.008 \\
% LGBCovPenalty & UTOPIA & 15.968 & 0.873 & \$62{,}126.710 & 23.631\% & 37.174\% & -4.155\% & 1.031 & -0.024 \\
% LGBCovPenalty & NASH & 34.176 & 0.868 & \$62{,}871.817 & 23.766\% & 37.780\% & -0.488\% & 1.025 & -0.016 \\
% \bottomrule
% \end{tabular}}
% \end{table}

% \begin{table}[!htbp]
% \centering
% \caption{Held-out test set performance. The penalty parameter $\rho$ is specific to the penalty definition.}
% \label{tab:test_results}
% \resizebox{\textwidth}{!}{%
% \begin{tabular}{llc|cc|cc|ccc}
% \toprule
% Family & Solution & $\rho$ & $R^2$ & MAE & COD & COV & VEI & PRD & PRB \\
% \midrule
% Linear Regression & BASELINE & -- & 0.762 & \$86{,}256.879 & 25.013\% & 36.825\% & -13.474\% & 1.039 & -0.011 \\
% LGBMRegressor & BASELINE & -- & 0.883 & \$72{,}063.853 & 23.106\% & 35.398\% & -36.816\% & 1.085 & -0.118 \\
% \midrule
% LGBSurrogatePenalty & CONSTRAINED & 5.789 & 0.880 & \$74{,}575.537 & 22.764\% & 33.423\% & -9.092\% & 1.039 & -0.013 \\
% LGBSurrogatePenalty & UTOPIA & 12.390 & 0.879 & \$75{,}846.355 & 23.249\% & 33.684\% & -7.446\% & 1.032 & 0.003 \\
% LGBSurrogatePenalty & NASH & 50.000 & 0.878 & \$78{,}155.478 & 24.092\% & 34.336\% & -4.883\% & 1.021 & 0.027 \\
% \midrule
% LGBCovPenalty & CONSTRAINED\_CLOSEST\_INFEASIBLE & 44.044 & 0.834 & \$95{,}978.845 & 23.558\% & 34.817\% & -9.602\% & 1.036 & -0.037 \\
% LGBCovPenalty & UTOPIA & 15.968 & 0.876 & \$73{,}837.127 & 23.239\% & 34.563\% & -14.985\% & 1.045 & -0.053 \\
% LGBCovPenalty & NASH & 34.176 & 0.689 & \$141{,}725.628 & 23.684\% & 34.853\% & -10.155\% & 1.037 & -0.037 \\
% \bottomrule
% \end{tabular}}
% \end{table}
% \paragraph{Constraining PRD,PRB, and VEI}




% \begin{table}[!htbp]
% \centering
% \caption{Average cross-validation results. The penalty parameter $\rho$ is specific to the penalty definition.}
% \label{tab:val_results}
% \resizebox{\textwidth}{!}{%
% \begin{tabular}{llc|cc|cc|ccc}
% \toprule
% Family & Solution & $\rho$ & $R^2$ & MAE & COD & COV & VEI & PRD & PRB \\
% \midrule
% Linear Regression & baseline & -- & 0.667 & \$92{,}441.063 & 32.765\% & 49.185\% & 8.633\% & 1.026 & 0.061 \\
% LGBMRegressor & baseline & -- & 0.890 & \$58{,}097.609 & 22.990\% & 37.623\% & -24.247\% & 1.069 & -0.077 \\
% \midrule
% LGBSurrogatePenalty & CONSTRAINED & 5.789 & 0.886 & \$61{,}122.635 & 23.124\% & 35.511\% & -3.446\% & 1.026 & -0.001 \\
% LGBSurrogatePenalty & UTOPIA & 6.572 & 0.885 & \$61{,}259.246 & 23.161\% & 35.469\% & -3.180\% & 1.025 & 0.002 \\
% LGBSurrogatePenalty & NASH & 4.492 & 0.886 & \$60{,}947.328 & 23.035\% & 35.503\% & -3.794\% & 1.028 & -0.005 \\
% \midrule
% LGBCovPenalty & CONSTRAINED\_CLOSEST\_INFEASIBLE & 44.044 & 0.853 & \$68{,}940.597 & 24.902\% & 39.202\% & 2.768\% & 1.023 & -0.008 \\
% LGBCovPenalty & UTOPIA & 20.578 & 0.872 & \$62{,}550.405 & 23.741\% & 37.499\% & -2.633\% & 1.029 & -0.021 \\
% LGBCovPenalty & NASH & 18.127 & 0.873 & \$62{,}151.192 & 23.628\% & 37.199\% & -3.827\% & 1.030 & -0.023 \\
% \bottomrule
% \end{tabular}}
% \end{table}

% \begin{table}[!htbp]
% \centering
% \caption{Held-out test set results. The penalty parameter $\rho$ is specific to the penalty definition.}
% \label{tab:test_results}
% \resizebox{\textwidth}{!}{%
% \begin{tabular}{llc|cc|cc|ccc}
% \toprule
% Family & Solution & $\rho$ & $R^2$ & MAE & COD & COV & VEI & PRD & PRB \\
% \midrule
% Linear Regression & baseline & -- & 0.762 & \$86{,}256.879 & 25.013\% & 36.825\% & -13.474\% & 1.039 & -0.011 \\
% LGBMRegressor & baseline & -- & 0.883 & \$72{,}063.853 & 23.106\% & 35.398\% & -36.816\% & 1.085 & -0.118 \\
% \midrule
% LGBSurrogatePenalty & CONSTRAINED & 5.789 & 0.880 & \$74{,}575.537 & 22.764\% & 33.423\% & -9.092\% & 1.039 & -0.013 \\
% LGBSurrogatePenalty & UTOPIA & 6.572 & 0.882 & \$74{,}428.114 & 22.806\% & 33.386\% & -8.545\% & 1.038 & -0.011 \\
% LGBSurrogatePenalty & NASH & 4.492 & 0.881 & \$73{,}941.570 & 22.645\% & 33.216\% & -10.249\% & 1.043 & -0.020 \\
% \midrule
% LGBCovPenalty & CONSTRAINED\_CLOSEST\_INFEASIBLE & 44.044 & 0.834 & \$95{,}978.845 & 23.558\% & 34.817\% & -9.602\% & 1.036 & -0.037 \\
% LGBCovPenalty & UTOPIA & 20.578 & 0.874 & \$74{,}209.834 & 23.257\% & 34.494\% & -14.353\% & 1.044 & -0.051 \\
% LGBCovPenalty & NASH & 18.127 & 0.765 & \$119{,}396.571 & 23.476\% & 34.580\% & -13.525\% & 1.043 & -0.048 \\
% \bottomrule
% \end{tabular}}
% \end{table}

% % \begin{table}[!htbp]
% % \centering
% % \caption{Summary of model performance metrics grouped by fold split, model family, and solution type.}
% % \label{tab:summary_models}
% % \resizebox{\textwidth}{!}{%
% % \begin{tabular}{lll|c|cc|cc|ccc}
% % \toprule
% % Split & Family & Solution & $\rho$ & $R^2$ & MAE & COD & COV & VEI & PRD & PRB \\
% % \midrule
% % validation\_avg & baseline & LinearRegression & -- & 0.667 & \$92{,}441.063 & 32.765\% & 49.185\% & 8.633\% & 1.026 & 0.061 \\
% % test & baseline & LinearRegression & -- & 0.762 & \$86{,}256.879 & 25.013\% & 36.825\% & -13.474\% & 1.039 & -0.011 \\
% % validation\_avg & baseline & LGBMRegressor & -- & 0.890 & \$58{,}097.609 & 22.990\% & 37.623\% & -24.247\% & 1.069 & -0.077 \\
% % test & baseline & LGBMRegressor & -- & 0.883 & \$72{,}063.853 & 23.106\% & 35.398\% & -36.816\% & 1.085 & -0.118 \\
% % \midrule
% % validation\_avg & simple\_selection & LGBCovPenalty\_CONSTRAINED\_CLOSEST\_INFEASIBLE & 44.044 & 0.853 & \$68{,}940.597 & 24.902\% & 39.202\% & 2.768\% & 1.023 & -0.008 \\
% % test & simple\_selection & LGBCovPenalty\_CONSTRAINED\_CLOSEST\_INFEASIBLE & 44.044 & 0.834 & \$95{,}978.845 & 23.558\% & 34.817\% & -9.602\% & 1.036 & -0.037 \\
% % validation\_avg & simple\_selection & LGBSurrogatePenalty\_CONSTRAINED & 5.789 & 0.886 & \$61{,}122.635 & 23.124\% & 35.511\% & -3.446\% & 1.026 & -0.001 \\
% % test & simple\_selection & LGBSurrogatePenalty\_CONSTRAINED & 5.789 & 0.880 & \$74{,}575.537 & 22.764\% & 33.423\% & -9.092\% & 1.039 & -0.013 \\
% % validation\_avg & simple\_selection & LGBCovPenalty\_UTOPIA & 20.578 & 0.872 & \$62{,}550.405 & 23.741\% & 37.499\% & -2.633\% & 1.029 & -0.021 \\
% % test & simple\_selection & LGBCovPenalty\_UTOPIA & 20.578 & 0.874 & \$74{,}209.834 & 23.257\% & 34.494\% & -14.353\% & 1.044 & -0.051 \\
% % validation\_avg & simple\_selection & LGBSurrogatePenalty\_UTOPIA & 6.572 & 0.885 & \$61{,}259.246 & 23.161\% & 35.469\% & -3.180\% & 1.025 & 0.002 \\
% % test & simple\_selection & LGBSurrogatePenalty\_UTOPIA & 6.572 & 0.882 & \$74{,}428.114 & 22.806\% & 33.386\% & -8.545\% & 1.038 & -0.011 \\
% % validation\_avg & simple\_selection & LGBCovPenalty\_NASH & 18.127 & 0.873 & \$62{,}151.192 & 23.628\% & 37.199\% & -3.827\% & 1.030 & -0.023 \\
% % test & simple\_selection & LGBCovPenalty\_NASH & 18.127 & 0.765 & \$119{,}396.571 & 23.476\% & 34.580\% & -13.525\% & 1.043 & -0.048 \\
% % validation\_avg & simple\_selection & LGBSurrogatePenalty\_NASH & 4.492 & 0.886 & \$60{,}947.328 & 23.035\% & 35.503\% & -3.794\% & 1.028 & -0.005 \\
% % test & simple\_selection & LGBSurrogatePenalty\_NASH & 4.492 & 0.881 & \$73{,}941.570 & 22.645\% & 33.216\% & -10.249\% & 1.043 & -0.020 \\
% % \bottomrule
% % \end{tabular}}
% % \end{table}




\subsection{Practical Implications for Assessor Workflows}

The empirical message is operationally simple. The proposed correction can be inserted into an existing assessor workflow with minimal disruption. It uses the same preprocessing pipeline, the same base model family, and the same downstream diagnostics already familiar to assessors. The only new tuning dimension is the penalty parameter \(\rho\), whose qualitative effect is transparent: larger values tend to reduce value-dependent ratio trends. In that sense, the contribution is not a replacement for current workflows, but a targeted correction layer that can expand the attainable accuracy--equity tradeoff set while remaining easy to calibrate, audit, and communicate.



% \section{Empirical Study: Cook County Assessor's Office}
% \label{sec:empirical}

% \subsection{Scope, Pipeline, and Data}

% The empirical study is based on the Cook County Assessor's Office workflow (CCAO). The CCAO manages extensive documentation and a Github repository \todo{add Github link as a citation}. For the experiments, we follow the documentation, available in the \texttt{model-res-avm} codebase and, specifically, the version available for 2025 assessment process (\texttt{assessment-year-2025}). This includes data preparation, preprocessing, and feature engineering, mostly implemented in the R programming language. However, our implementation is in Python~3.10.14 \todo{cite?}, as a translation of the original R-based workflow. The temporal splits are chosen to mimic a realistic assessor setting: 2016--2022 form the train/validation development period, 2023 is used as the out-of-time test year, and 2024 is treated as the current assessment year reserved for later deployment analysis (see for specific details  \todo{Is this necessary? I was working with 2023 as test and 2024 as assessment, but not sure if makes sense for the paper}).

% \todo{
% At the current stage of the project, the empirical evidence reported in the manuscript comes from the validation and test components of this translated pipeline.

% The final paper should continue to state clearly whether any remaining differences persist between the translated implementation and the original production workflow, since that distinction matters for both reproducibility and external interpretation.}

% % \subsection{Preprocessing, Filtering, and Sample Construction}

% \todo{
% Because the empirical contribution is intended as a real workflow study rather than a purely conceptual example, the final paper should report the main inclusion rules, exclusion rules, and final sample sizes used in the translated CCAO pipeline. Table~\ref{tab:sample_construction_placeholder} is a placeholder for that information. At a minimum, the final version should document the property scope, any multicard or outlier exclusions, the construction of the sale-price target, and the final observation counts by temporal split.
% }

% \begin{table}[!ht]
% \centering
% \caption{Placeholder for sample construction summary. Replace with the final counts and main filtering rules from the experiment pipeline.}
% \label{tab:sample_construction_placeholder}
% \fbox{\parbox{0.9\textwidth}{
% \textbf{Insert final sample-construction table here.} Suggested columns: stage, main rule applied, number of remaining observations, and split-specific counts. Suggested rows: raw extract, sale verification filters, multicard exclusion, outlier exclusion, modeling-ready sample, and split counts for train/validation, test, and assessment.
% }}
% \end{table}

% \subsection{Exact Validation Windows and Fold Weights}

% \todo{
% The constrained calibration framework in Section~\ref{sec:model_selection} is defined over validation folds, so the final paper should report the exact rolling-origin or expanding-window construction used inside the 2016--2022 development period. Table~\ref{tab:fold_definition_placeholder} is a placeholder for that exact specification. This should include the training window, validation window, whether the training sample expands over time, and the fold weights \(\alpha_v\) used in metric aggregation.
% }

% \begin{table}[!ht]
% \centering
% \caption{Placeholder for exact validation-window construction.}
% \label{tab:fold_definition_placeholder}
% \fbox{\parbox{0.9\textwidth}{
% \textbf{Insert final fold-definition table here.} Suggested columns: fold index, training years, validation years, number of validation observations, and aggregation weight \(\alpha_v\). This table should make the calibration protocol fully reproducible.
% }}
% \end{table}

% \subsection{Baselines: Linear Versus Nonlinear Models}

% As mentioned in Section \todo{Check}, the CCAO used Linear Regression-based prediction models for the mass appraisal process until 2019 \todo{check}, and then transitioned to the nonlinear boosted tree model LightGBM. In terms of accuracy, the transition is consistent with the literature; boosted trees tend to perform better on tabular data \todo{cite}.

% Tables~\ref{tab:baseline_models} and~\ref{tab:baseline_models_2} report the current benchmark comparison for the CCAO's data on the test set and assessment sets. The baseline LightGBM model substantially improves predictive accuracy relative to linear regression. However, at the same time it worsens vertical-equity diagnostics. Relative to Linear Regression, the baseline LightGBM model improves \(R^2\), MAE, COD, and COV, but its VEI, PRD, and PRB values are worsening, until the point they are clearly signaling a regressive behavior, while Linear regression seems much more stable in this dimension.

% \medskip

% \noindent \textbf{Remark.} This is the main empirical motivation for the paper: stronger nonlinear predictive fit does not automatically imply better vertical equity. This not a surprising behavior; nonlinear models can capture much more stronger relationships, and thus are more susceptible to produce non-equitable patters like in (\ref{eq:covariance_population_identity}) \todo{Verify it?}\todo{Cite maybe?}. Therefore, there is a necessity for a middle-ground solution between a simple but equitable model and more complex one that can achieve higher prediction performance, which motivates our solution.

% \begin{table}[!ht]
% \centering
% \caption{Baseline comparison on the \textit{test} set ($\sim$2023). Lower is better for MAE, COD, and COV. For PRD, values closer to one are preferred; for VEI and PRB, values closer to zero are preferred.}
% \label{tab:baseline_models}
% \resizebox{\textwidth}{!}{%
% \begin{tabular}{lrrrrrrr}
% \toprule
% Model & \(R^2\) & MAE & COD & COV & PRD & PRB & VEI \\
% \midrule
% Linear Regression & 0.762 & \$86{,}257 & 25.013\% & 36.825\% & 1.039 & -0.011 & -13.474\% \\
% LightGBM & 0.887 & \$71{,}337 & 22.739\% & 34.805\% & 1.082 & -0.114 & -36.810\% \\
% \bottomrule
% \end{tabular}}
% \end{table}

% \begin{table}[!ht]
% \centering
% \caption{Baseline comparison on the \textit{assessment} set (2024). Lower is better for MAE, COD, and COV. For PRD, values closer to one are preferred; for VEI and PRB, values closer to zero are preferred.}
% \label{tab:baseline_models_2}
% \resizebox{\textwidth}{!}{%
% \begin{tabular}{lrrrrrrr}
% \toprule
% Model & \(R^2\) & MAE & COD & COV & PRD & PRB & VEI \\
% \midrule
% LinearRegression & 0.802 & \$90{,}749 & 24.199\% & 36.622\% & 1.042 & -0.022 & -13.352\% \\
% LGBMRegressor & 0.886 & \$79{,}816 & 21.310\% & 35.097\% & 1.071 & -0.096 & -28.208\% \\
% \bottomrule
% \end{tabular}}
% \end{table}

% \todo{Table~\ref{tab:baseline_models} should be paired with the main motivation plot in the final version. Figure~\ref{fig:baseline_regressivity_placeholder} is the placeholder for that visual.
% Why would we add the plot?}

% \begin{figure}[t]
% \centering
% \fbox{\parbox{0.9\textwidth}{
% \textbf{Insert baseline regressivity plot here.} Suggested figure: assessment ratio \(r=\hat P/P\) versus log market value under the baseline LightGBM model, including the horizontal \(r=1\) line, a LOWESS trend, and an optional linear trend line. This is the main visual motivation for the regressivity problem.
% }}
% \caption{Placeholder for the baseline regressivity motivation figure.}
% \label{fig:baseline_regressivity_placeholder}
% \end{figure}

% \subsection{Effect of Covariance Penalization}

% We aim to get a middle-ground solution that can achieve better tradeoffs between accuracy and vertical-equity, when compared to the baseline models (Linear Regression and LightGBM). We propose to add covariance-based penalties to the baseline LightGBM learner. The current empirical results show the basic mechanism clearly \todo{CHECK THE NUMBERS}. For the baseline model, the ratio--log-value relationship exhibits a strong negative trend, with a reported linear slope of approximately \(-0.2156\) and ratio--price correlation of approximately \(-0.2666\). Under the squared covariance penalty, the slope moves to about \(-0.1263\) and the correlation to about \(-0.1224\). Under the separable surrogate penalty, the slope moves even closer to flat, about \(-0.0222\), and the correlation to about \(-0.0239\). These are not deployment diagnostics by themselves, but they visually illustrate that the training-time penalties act in the intended direction.

% The final paper should show that before/after comparison directly. Figure~\ref{fig:penalty_mechanism_placeholder} is the placeholder for the main paired ratio-trend figure.

% \begin{figure}[!htbp]
%     \centering
%     \begin{subfigure}{0.48\textwidth}
%         \centering
%         \includegraphics[width=\linewidth, trim={-3cm 0mm 0cm 3mm}]{Fairness Price/GE Paper/img/Baseline LightGBM_0_motivation.pdf}
%         \caption{\scriptsize Baseline LightBGM}
%     \end{subfigure}\hfill
%     \begin{subfigure}{0.48\textwidth}
%         \centering
%         \includegraphics[width=\linewidth, trim={0cm 0mm -3cm 3mm}]{Fairness Price/CCAO meeting/img/Covariance-Penalized LightGBM_2_motivation.pdf}
%         \caption{\scriptsize Covariance-Penalized LightGBM}
%     \end{subfigure}

%     \caption{LightGBM example comparison. / Placeholder for the main mechanism figure showing how covariance penalization flattens the ratio--value trend. /  Suggested layout: two panels comparing baseline LightGBM against a squared-covariance or surrogate-penalized LightGBM, each showing ratio versus log value, the \(r=1\) line, and trend overlays. Add the reported slope and correlation summaries in the caption or callout box.}
% \end{figure}
% \vspace{-5mm}

% \subsection{Penalized LightGBM Performance on the Test Set}

% Table~\ref{tab:penalized_models} summarizes the current test-set results for the main penalized LightGBM variants. All four reported penalized models move the vertical-equity diagnostics sharply toward assessor-preferred ranges. For example, the baseline LightGBM has PRD \(=1.082\) and PRB \(=-0.108\), both strongly regressive, whereas the penalized variants move PRD into the 1.029--1.030 range and PRB into the \(-0.026\) to \(0.008\) range. At the same time, the best reported squared-covariance configuration improves \(R^2\) from \(0.860\) to \(0.880\) and improves MAE from \$73{,}398 to \$70{,}855. These results are exactly the kind of middle-ground behavior the method is designed to find.

% \begin{table}[htbp]
%     \centering
%     \caption{\textbf{Penalized Objective Comparison:} Pred. Error + $\rho\; \cdot$ Covariance(ratios, prices)$^2$. Test results for Accuracy (R2, MAE), horizontal equity (COD, COV), and vertical equity (VEI, PRD, PRB, MKI). / / Covariance-penalized LightGBM results on the test set. The reported \(\rho\) values are penalty-specific tuning parameters and are not directly comparable across penalty definitions.}
%     \label{tab:penalized_models}
%     \resizebox{\textwidth}{!}{%
%     \begin{tabular}{lc|cc|cc|cccc}
%     \toprule
%      Model & Penalty $\rho$ & R2 & MAE & COD & COV (IAAO) & VEI & PRD & PRB & MKI \\
%     \midrule
%     LinearRegression & -- & 0.798 & \$85,273 & 24.306 & 0.354 & 6.223\% & 1.016 & 0.033 & 1.009 \\
%     LGBMRegressor & -- & 0.860 & \$73,398 & 21.647 & 0.330 & \textcolor{red}{-31.431\%} & \textcolor{red}{1.082} & \textcolor{red}{-0.108} & \textcolor{red}{0.890} \\
%     \midrule
%     LGBMSurrPenalty [v1] & 2.043 &
%     \makecell{\textbf{0.872}\\{\scriptsize(\textcolor{green!60!black}{+1.40\%})}} &
%     \makecell{\$74,579\\{\scriptsize(\textcolor{red!70!black}{-1.61\%})}} &
%     \makecell{21.770\\{\scriptsize(\textcolor{red!70!black}{-0.57\%})}} &
%     \makecell{\textbf{0.316}\\{\scriptsize(\textcolor{green!60!black}{+4.24\%})}} &
%     \makecell{\textbf{-0.043\%}\\{\scriptsize(\textcolor{green!60!black}{+99.86\%})}} &
%     \makecell{\textbf{1.030}\\{\scriptsize(\textcolor{green!60!black}{+63.41\%})}} &
%     \makecell{\textbf{+0.003}\\{\scriptsize(\textcolor{green!60!black}{+97.22\%})}} &
%     \makecell{\textbf{0.970}\\{\scriptsize(\textcolor{green!60!black}{+72.73\%})}} \\
%     [1.5mm]

%     LGBMSurrPenalty [v2] & 280.7 &
%     \makecell{\textbf{0.867}\\{\scriptsize(\textcolor{green!60!black}{+0.81\%})}} &
%     \makecell{\$75,641\\{\scriptsize(\textcolor{red!70!black}{-3.06\%})}} &
%     \makecell{21.829\\{\scriptsize(\textcolor{red!70!black}{-0.84\%})}} &
%     \makecell{\textbf{0.313}\\{\scriptsize(\textcolor{green!60!black}{+5.15\%})}} &
%     \makecell{\textbf{+3.403\%}\\{\scriptsize(\textcolor{green!60!black}{+89.17\%})}} &
%     \makecell{\textbf{1.029}\\{\scriptsize(\textcolor{green!60!black}{+64.63\%})}} &
%     \makecell{\textbf{+0.008}\\{\scriptsize(\textcolor{green!60!black}{+92.59\%})}} &
%     \makecell{\textbf{0.970}\\{\scriptsize(\textcolor{green!60!black}{+72.73\%})}} \\
%     [1.5mm]

%     LGBMCovPenalty [v1] & 6.58 &
%     \makecell{\textbf{0.877}\\{\scriptsize(\textcolor{green!60!black}{+1.98\%})}} &
%     \makecell{\textbf{\$72,042}\\{\scriptsize(\textcolor{green!60!black}{+1.85\%})}} &
%     \makecell{21.693\\{\scriptsize(\textcolor{red!70!black}{-0.21\%})}} &
%     \makecell{\textbf{0.319}\\{\scriptsize(\textcolor{green!60!black}{+3.33\%})}} &
%     \makecell{\textbf{-6.579\%}\\{\scriptsize(\textcolor{green!60!black}{+79.07\%})}} &
%     \makecell{\textbf{1.030}\\{\scriptsize(\textcolor{green!60!black}{+63.41\%})}} &
%     \makecell{\textbf{-0.026}\\{\scriptsize(\textcolor{green!60!black}{+75.93\%})}} &
%     \makecell{\textbf{0.989}\\{\scriptsize(\textcolor{green!60!black}{+90.00\%})}} \\
%     [1.5mm]

%     LGBMCovPenalty [v2] & 1487.0 &
%     \makecell{\textbf{0.880}\\{\scriptsize(\textcolor{green!60!black}{+2.33\%})}} &
%     \makecell{\textbf{\$70,855}\\{\scriptsize(\textcolor{green!60!black}{+3.46\%})}} &
%     \makecell{\textbf{21.467}\\{\scriptsize(\textcolor{green!60!black}{+0.83\%})}} &
%     \makecell{\textbf{0.317}\\{\scriptsize(\textcolor{green!60!black}{+3.94\%})}} &
%     \makecell{\textbf{-4.043\%}\\{\scriptsize(\textcolor{green!60!black}{+87.14\%})}} &
%     \makecell{\textbf{1.029}\\{\scriptsize(\textcolor{green!60!black}{+64.63\%})}} &
%     \makecell{\textbf{-0.022}\\{\scriptsize(\textcolor{green!60!black}{+79.63\%})}} &
%     \makecell{\textbf{0.986}\\{\scriptsize(\textcolor{green!60!black}{+87.27\%})}} \\
%     \bottomrule
%     \end{tabular}
%     }
% \end{table}

% If desired, the final paper can also add a second table with percentage improvements relative to the baseline LightGBM model. That presentation is useful for quick interpretation, but Table~\ref{tab:penalized_models} should remain the primary source of the raw metric values.

% \subsection{Tradeoff Frontiers}

% The current results indicate that increasing \(\rho\) traces a meaningful accuracy--equity tradeoff curve rather than producing a degenerate monotone sacrifice of fit. In PRD--\(R^2\) space, the penalized models move from the unconstrained LightGBM point toward the assessor compliance band while remaining near the upper edge of the accuracy frontier. A similar pattern appears in VEI--\(R^2\) space. These frontier views justify the constrained calibration logic of Section~\ref{sec:model_selection}: in a genuinely multi-criteria setting, the practical question is not which model is globally best, but which feasible model offers the best accuracy among those that satisfy the chosen assessment-quality requirements.

% \begin{figure}[!htbp]
%     \centering
%     \begin{subfigure}[b]{0.48\textwidth}
%         \centering
%         \includegraphics[width=\textwidth, trim={2mm 0 4mm 2mm}]{Fairness Price/GE Paper/img/tradeoffs/PRD_vs_R2_band.pdf}
%         \caption{Insert PRD description here}
%         \label{fig:prd_tradeoff}
%     \end{subfigure}
%     \hfill
%     \begin{subfigure}[b]{0.48\textwidth}
%         \centering
%         \includegraphics[width=\textwidth, trim={4mm 0 0 2mm}]{Fairness Price/GE Paper/img/tradeoffs/VEI_vs_R2_band.pdf}
%         \caption{Insert VEI description here}
%         \label{fig:vei_tradeoff}
%     \end{subfigure}

%     \caption{\textbf{Pareto Frontier:} Test set behavior for different penalization terms ($\rho >0$).}
%     \label{fig:pareto_frontier}
% \end{figure}

% \begin{figure}[!ht]
% \centering
% \includegraphics[width=0.7\linewidth, trim={2mm 2mm 4mm 4mm}]{Fairness Price/GE Paper/img/model-selection/pareto_optimal.pdf}
% \caption{\parbox{0.9\textwidth}{ Suggested layout: one panel for PRD versus \(R^2\) with the compliance band indicated, and one panel for VEI versus \(R^2\). Label baseline linear regression, baseline LightGBM, and the penalized families across increasing \(\rho\).
% }}
% \label{fig:frontier_placeholder}
% \end{figure}

% \subsection{Validation-Fold Constrained Calibration}

% One of the paper's methodological contributions is the constrained calibration rule of Section~\ref{sec:model_selection}, so the empirical section should also instantiate that rule directly. The final version should report the exact active validation constraints, the number of candidate configurations evaluated, the number of feasible configurations, the selected configuration, and the corresponding performance on the untouched test year. Table~\ref{tab:calibration_constraints_placeholder} and Table~\ref{tab:selected_model_placeholder} are placeholders for that information.

% \begin{table}[t]
% \centering
% \caption{Placeholder for the exact empirical calibration rule used on validation folds.}
% \label{tab:calibration_constraints_placeholder}
% \fbox{\parbox{0.9\textwidth}{
% \textbf{Insert final calibration-constraints table here.} Suggested columns: diagnostic, constraint type, lower bound, upper bound, and whether the constraint is aggregated across folds or enforced foldwise. Typical rows may include PRD, PRB, VEI, COD, COV, and ratio-statistic guardrails.
% }}
% \end{table}

% \begin{table}[t]
% \centering
% \caption{Placeholder for the final selected configuration under constrained calibration.}
% \label{tab:selected_model_placeholder}
% \fbox{\parbox{0.9\textwidth}{
% \textbf{Insert selected-model summary here.} Suggested rows or columns: selected model family, penalty type, \(\rho\), number of feasible candidates, best infeasible accuracy-only model, final selected feasible model, and held-out test metrics for the selected configuration.
% }}
% \end{table}

% This subsection is also the natural place to explain any difference between (i) the full frontier of candidate models and (ii) the single final model selected for deployment under the chosen assessment-quality requirements.

% \subsection{Uncertainty, Stability, and Fold Heterogeneity}

% Because regressivity diagnostics can be noisy across time windows, the final paper should also report a stability summary rather than relying only on point estimates. This can be done through bootstrap confidence intervals on the held-out test metrics, foldwise distributions across validation windows, or both. Figure~\ref{fig:stability_placeholder} is the placeholder for that summary.

% \begin{figure}[t]
% \centering
% \fbox{\parbox{0.9\textwidth}{
% \textbf{Insert stability/uncertainty figure here.} Suggested options: foldwise boxplots for PRD, PRB, VEI, COD, and \(R^2\); bootstrap confidence intervals for the selected model versus the main baselines; or a table of foldwise means and standard deviations.
% }}
% \caption{Placeholder for stability and uncertainty reporting.}
% \label{fig:stability_placeholder}
% \end{figure}

% \subsection{Practical Implications for Assessor Workflows}

% The empirical message is operationally simple. The method can be inserted into an existing valuation workflow with minimal changes. It uses the same preprocessing, the same base model classes, and the same downstream diagnostics already familiar to assessors. It adds one main penalty parameter with a transparent qualitative effect: larger \(\rho\) values tend to move models toward less regressive behavior. This makes the method attractive not only because it improves the attainable tradeoffs, but because it does so in a form that can be calibrated, audited, and communicated within existing assessor practice.





% \section{(Optional) important points of the approach. Other have empirical examination or case of study can be an option.}

\section{Discussion and Conclusions}
\label{sec:discussion}

We studied how to reduce regressivity in ML-based mass appraisal with minimal modification to existing assessor workflows. The main message is empirical: predictive accuracy alone does not guarantee vertical equity, and the move from linear regression to a nonlinear learner inside the CCAO pipeline improves $R^2$, MAE, COD, and COV while simultaneously worsening PRD, PRB, and VEI. A soft penalty on the covariance between log-residuals and log-prices recovers most of this equity gap at little cost in predictive fit, and does so with a single hyperparameter. The two penalty forms play complementary roles: the direct squared-covariance penalty is the closest instantiation of the dependence-control target; the separable surrogate is sample-additive and fits inside standard boosting custom-objective interfaces. In our CCAO runs the surrogate delivers the strongest combined performance and is the variant we recommend for routine deployment.

A second contribution is a calibration perspective. Once equity enters the loss through $\rho$, a single accuracy score no longer summarizes candidate quality, and the constrained calibration problem we use is a direct analogue of the ratio-study compliance logic that assessors already apply to candidate models. The alternative Pareto criteria of Section~\ref{subsubsec:pareto_alternatives} and the convex-stacking extension described in Appendix~\ref{app:pareto_and_stacking} are available when explicit tolerances are unavailable or jointly infeasible; they did not dominate the constrained rule in our experiments, and we keep them as secondary tools rather than as coequal contributions.

\paragraph{Limitations.}
The method controls the covariance between log-residuals and log-prices, which is directionally aligned with the standard vertical-equity diagnostics but is not equal to them. In-sample or validation control of $\widehat{\Cov}(e,y)$ therefore does not imply out-of-sample PRD/PRB/VEI compliance, and we do not claim such a guarantee. The covariance is a first-order summary, so nonmonotone local distortions can cancel; the surrogate partially mitigates this by penalizing ratio deviations pointwise, but does not eliminate it. The penalties are global across the sample and do not currently use sub-market or geographic structure. The constrained calibration rule can become infeasible when many ratio-study tolerances are enforced simultaneously. Finally, the CCAO pipeline on which we evaluate the method is a translation rather than a direct execution of the production code, so the absolute levels we report should be interpreted within that translated pipeline.

\paragraph{Extensions.}
Two directions are especially natural. First, moving beyond first-order dependence via nonlinear measures such as HSIC or maximal correlation \cite{LiEtAl2022Fairness,lee2022maximal}, while keeping the separable implementation pattern used here. Second, stratifying the penalty by sub-market, neighborhood, or property class so that local value-dependent patterns are not averaged away by a global correction. Broader directions include applying the framework to additional assessor offices and base learners (random forests, neural networks, tabular foundation models), extending the constraint set to include central-tendency and dispersion diagnostics, and combining the penalty with interpretability tools that communicate the effect of $\rho$ to non-technical audiences.

% \section{Discussion}



% Our contribution is to connect these two literatures in an assessor-compatible way. We start from an assessor-defined problem---regressivity as a systematic pattern in the relationship between assessment ratios and value level---and recast it in a form that fair-regression tools can act on directly during training. This shifts the role of vertical-equity analysis from purely ex-post diagnosis toward in-processing control. Concretely, we use this interpretation to build a regularization approach for standard regression models used in mass appraisal and embed the resulting corrections in a multi-objective calibration framework based on accuracy--equity tradeoffs. In this sense, the paper does not introduce equity concerns into mass appraisal, nor does it introduce fairness regularization into machine learning; rather, it develops a practical bridge between assessor-defined vertical equity and fairness-aware regression methods.

% % \section{Discussion/Conclusions, Limitations, Extensions \& Future Work}

% % \subsection{Dicussion}

% \subsection{Limitations.}


% \begin{itemize}
%     \item  This does not capture all nonlinear dependencies, so we could be enforcing them. Specially, if penalizing too much.
%     \item Is not aware of sub-market behaviors. The nonlinear relationships enforce seem to have a neighborhood or sub-market structure.
%     \item 
    
% \end{itemize}





% \subsection{Extensions.}


% \subsection{Future Work}

% Same direction: control the relationship, but now aware of the sub-market clusters, which can create nonlinear relationships when asking for too much overall linear correction (see image of 0-correlation plot of surrogate penalty).

% Interpretable \& fair models. Apply this to multiple cases of assessors offices, neural networks as n alternative methodology because of their predictive power (?); in particular, their ability to handle Geospatial data (e.g., apply an embedding using the information or methodology provided by AlphaEarth of Google). Another alternative is to test the foundational models (e.g., tabular models). Try to incorpore it into Random Forest/Bagging, which is also widely used apparently.





% % \section{MY NOTES}

% % Paper structure references. INFORMS seems to use a lot of plots and tables in general (make sense since they care about the analytics part).
% % \begin{itemize}
% %     \item INFORMS data science: https://pubsonline.informs.org/doi/epdf/10.1287/ijds.2024.0029
% %     \item INFORMS managements science: https://pubsonline.informs.org/doi/epdf/10.1287/mnsc.2022.00345
% %     \item INFORM journal on optimization: https://arxiv.org/pdf/2412.17116
% % \end{itemize}

% % \noindent Haihao whiteboard:
% % \begin{itemize}
% %     \item Introduction: motivation, contribution, related work.
% %     \item Setting: Examples, assumptions, etc.
% %     \item Methodology
% %     \item Experiments
% %     \item Extension (2-pages at most)
% %     \item Conclusions
% % \end{itemize}

% % \noindent Swati Comments / Suggestions:

% % \begin{itemize}
% %     \item What is novel in this research (theory and applications-wise)? Has anything like this be done already?
% %     \item Good thing: it looked like I owned the project. Lets keep it that way for the GE.
% %     \item It has to be a journal-level paper. Good venues: Journal of Machine Learning Research (JMLR) or Management Sciences.
% %     \item Paper Critique: adjacent to what I've done. Maybe ML, fairness, etc. The idea is to be able to read, summarize, critique (Not just chatGPT-style). Even go on details, proofs, technical stuff. I can do extra research on the topic of the paper.
% %     \item Important: GE as a message of what I've been doing during the phd.
% %     \item Presentation (oral-exam): Be prepared for presenting + questions at the same time. Bad things about past presentations: too much notation, not be able to answer extra questions, etc.
% %     \item Benchmark v/s more complex correlation/covariance methods. They don't work/generalize as good. Show it.
% %     \item Extras: copy of the paper \& slides in physical for reading.
% % \end{itemize}

\newpage

\printbibliography %Prints bibliography

\newpage

\appendix



% \section{A Theoretical Analysis of Regressivity as a Dependence Problem}
% \label{anex:regressivity_dependence}

% The assessor literature describes \emph{regressivity} as a price-related pattern in which assessment ratios tend to decline as property value increases \cite{IAAO2013Ratio,IAAO2025ExposureRatio}. In the predictive setting of Section~\ref{subsec:prediction_setting}, this pattern can be expressed directly in terms of the model outputs.

% For a fitted model $f$ and an evaluation sample $S$, recall the assessment ratio
% \[
% r_i(f):=\frac{\widehat P_i(f)}{P_i},
% \qquad i\in S,
% \]
% and define the log-ratio residual
% \[
% u_i(f):=\log r_i(f)=\hat y_i(f)-y_i,
% \]
% where the identity follows from \eqref{eq:log_ratio_identity}. Thus, any price-related pattern in the assessment ratios can equivalently be studied as a dependence pattern between $u_i(f)$ and a value index.

% To keep the exposition slightly more general, let
% \[
% z_i:=\phi(P_i),
% \]
% where $\phi:\mathbb{R}_{++}\to\mathbb{R}$ is any strictly increasing transform of price. Throughout the paper the default choice is $\phi(P)=\log P$, but the arguments below only require monotonicity. For any pair of sample vectors $a=\{a_i\}_{i\in S}$ and $b=\{b_i\}_{i\in S}$, define
% \begin{equation}
% \label{eq:sample_cov_general}
% \widehat{\operatorname{Cov}}_S(a,b)
% :=
% \frac{1}{|S|}
% \sum_{i\in S}
% (a_i-\bar a_S)(b_i-\bar b_S),
% \end{equation}
% with the analogous definition for $\widehat{\operatorname{Var}}_S(b)$. A natural log-scale summary of directional regressivity is therefore
% \begin{equation}
% \label{eq:cov_log_resid}
% \widehat{\operatorname{Cov}}_S(u(f),z),
% \end{equation}
% and the corresponding linear trend is the slope
% \begin{equation}
% \label{eq:sample_slope_logz}
% \widehat\beta^{\,\log}_S(f)
% :=
% \frac{\widehat{\operatorname{Cov}}_S(u(f),z)}
% {\widehat{\operatorname{Var}}_S(z)}.
% \end{equation}
% Because $\widehat{\operatorname{Var}}_S(z)$ is fixed once the evaluation sample is fixed, controlling the sign and magnitude of the linear price-related trend is equivalent to controlling the covariance in \eqref{eq:cov_log_resid}.

% At the same time, implementation is often simpler on the original ratio scale around the target of perfect equity. For this reason, define the price-scale deviation
% \[
% e_i(f):=r_i(f)-1.
% \]
% Since subtracting a constant does not change covariance,
% \[
% \widehat{\operatorname{Cov}}_S(e(f),z)
% =
% \widehat{\operatorname{Cov}}_S(r(f),z).
% \]
% Moreover, because $\sum_{i\in S}(z_i-\bar z_S)=0$,
% \begin{equation}
% \label{eq:cov_linear_form}
% \widehat{\operatorname{Cov}}_S(e(f),z)
% =
% \frac{1}{|S|}
% \sum_{i\in S} e_i(f)\,(z_i-\bar z_S).
% \end{equation}
% This representation motivates a separable surrogate penalty.

% \begin{lemma}[Exact surrogate upper bound]
% \label{lem:surrogate_upper_bound}
% Define
% \begin{equation}
% \label{eq:surrogate_penalty}
% \Psi^{\mathrm{surr}}_S(f)
% :=
% \frac{1}{|S|}
% \sum_{i\in S} e_i(f)^2\,(z_i-\bar z_S)^2.
% \end{equation}
% Then, for every fitted model $f$,
% \begin{equation}
% \label{eq:cov_surrogate_bound}
% \widehat{\operatorname{Cov}}_S(e(f),z)^2
% \le
% \Psi^{\mathrm{surr}}_S(f).
% \end{equation}
% \end{lemma}

% \begin{proof}
% By \eqref{eq:cov_linear_form},
% \[
% \widehat{\operatorname{Cov}}_S(e(f),z)
% =
% \frac{1}{|S|}
% \sum_{i\in S} a_i(f),
% \qquad
% a_i(f):=e_i(f)\,(z_i-\bar z_S).
% \]
% Hence, by Jensen's inequality for $x\mapsto x^2$,
% \[
% \widehat{\operatorname{Cov}}_S(e(f),z)^2
% =
% \left(
% \frac{1}{|S|}
% \sum_{i\in S} a_i(f)
% \right)^2
% \le
% \frac{1}{|S|}
% \sum_{i\in S} a_i(f)^2
% =
% \frac{1}{|S|}
% \sum_{i\in S} e_i(f)^2\,(z_i-\bar z_S)^2,
% \]
% which is exactly \eqref{eq:cov_surrogate_bound}.
% \end{proof}

% Lemma~\ref{lem:surrogate_upper_bound} gives the optimization interpretation used later: $\Psi_S^{\mathrm{surr}}(f)$ is not merely heuristic, but an exact upper bound on the squared covariance between the price-scale deviation from perfect equity and the value index. In addition, unlike the covariance itself, the surrogate is separable across observations, which makes it more convenient for gradient-based learners and libraries that are built around per-observation derivatives.

% To connect this dependence view to standard ratio-study metrics, we now formalize the common monotone picture of regressivity.

% \begin{assumption}[Monotone price-related behavior]
% \label{ass:monotone_regressivity}
% For a fitted model $f$ and evaluation sample $S$, there exists a nonincreasing function
% \[
% g_f:\phi(\mathbb{R}_{++})\to\mathbb{R}_{++}
% \]
% such that
% \[
% r_i(f)=g_f(z_i),
% \qquad i\in S.
% \]
% Equivalently, whenever $z_i\le z_j$, one has $r_i(f)\ge r_j(f)$.
% \end{assumption}

% Assumption~\ref{ass:monotone_regressivity} is the finite-sample counterpart of the usual assessor description that ratios tend to fall with value. Because both $\phi$ and $\log$ are increasing, it immediately implies that $u_i(f)=\log r_i(f)$ is also nonincreasing in $z_i$.

% \begin{proposition}[Sign of covariances under monotone regressivity]
% \label{prop:cov_signs}
% Under Assumption~\ref{ass:monotone_regressivity},
% \begin{align}
% \widehat{\operatorname{Cov}}_S(r(f),z) &\le 0,
% \label{eq:cov_r_z_negative}
% \\
% \widehat{\operatorname{Cov}}_S(u(f),z) &\le 0,
% \label{eq:cov_u_z_negative}
% \\
% \widehat{\operatorname{Cov}}_S(r(f),P) &\le 0.
% \label{eq:cov_r_p_negative}
% \end{align}
% If $g_f$ is not constant on $S$, then the inequalities are strict.
% \end{proposition}

% \begin{proof}
% For any two sample vectors $a$ and $b$,
% \begin{equation}
% \label{eq:pairwise_cov_identity}
% \widehat{\operatorname{Cov}}_S(a,b)
% =
% \frac{1}{2|S|^2}
% \sum_{i,j\in S}
% (a_i-a_j)(b_i-b_j).
% \end{equation}
% Applying \eqref{eq:pairwise_cov_identity} with $a_i=r_i(f)$ and $b_i=z_i$, each summand is nonpositive because $r_i(f)=g_f(z_i)$ with $g_f$ nonincreasing. This proves \eqref{eq:cov_r_z_negative}.

% Next, since $u_i(f)=\log r_i(f)$ and $\log$ is increasing, the map $z\mapsto u(f)=\log(g_f(z))$ is also nonincreasing, so the same argument yields \eqref{eq:cov_u_z_negative}.

% Finally, because $\phi$ is strictly increasing, its inverse preserves order on the sample, so $P_i$ is increasing in $z_i$. Therefore $r_i(f)$ is nonincreasing in $P_i$ as well, and another application of \eqref{eq:pairwise_cov_identity} gives \eqref{eq:cov_r_p_negative}.
% \end{proof}

% Proposition~\ref{prop:cov_signs} shows that the sign of the covariance in either log space or price space is aligned with the usual notion of regressivity. The next result connects this to the assessor-facing vertical-equity metrics of Section~\ref{subsec:equity_metrics}.

% \begin{proposition}[Implications for ratio-study metrics]
% \label{prop:metrics_alignment}
% Assume Assumption~\ref{ass:monotone_regressivity}. Then the following hold.

% \smallskip
% \noindent
% \emph{(i) Log-slope.} The slope of log ratios on the value index is nonpositive:
% \[
% \widehat\beta^{\,\log}_S(f)\le 0.
% \]

% \smallskip
% \noindent
% \emph{(ii) PRD.} Let $\bar r_S:=|S|^{-1}\sum_{i\in S} r_i(f)$ and $\bar P_S:=|S|^{-1}\sum_{i\in S} P_i$. Then
% \begin{equation}
% \label{eq:prd_cov_identity}
% \operatorname{PRD}(S)
% =
% \frac{1}
% {1+\widehat{\operatorname{Cov}}_S(r(f),P)/(\bar r_S\bar P_S)}.
% \end{equation}
% Consequently,
% \[
% \widehat{\operatorname{Cov}}_S(r(f),P)\le 0
% \quad\Longrightarrow\quad
% \operatorname{PRD}(S)\ge 1,
% \]
% with strict inequality whenever $\widehat{\operatorname{Cov}}_S(r(f),P)<0$.

% \smallskip
% \noindent
% \emph{(iii) PRB.} Let $v_i$ denote the value proxy used in the PRB calculation, and assume that it preserves the ordering by value on $S$ (that is, $P_i\le P_j$ implies $v_i\le v_j$). Then
% \[
% \operatorname{PRB}(S)\le 0.
% \]

% \smallskip
% \noindent
% \emph{(iv) VEI.} Let the lowest-value and highest-value groups used in the VEI be formed by the same ordering proxy $v_i$. Then, under the sign convention of Section~\ref{subsec:equity_metrics},
% \[
% \operatorname{VEI}(S)\le 0.
% \]
% \end{proposition}

% \begin{proof}
% Part (i) follows immediately from \eqref{eq:sample_slope_logz}, since $\widehat{\operatorname{Var}}_S(z)>0$ on any nondegenerate evaluation sample and Proposition~\ref{prop:cov_signs} gives $\widehat{\operatorname{Cov}}_S(u(f),z)\le 0$.

% For part (ii), note that
% \[
% \overline{r_W}
% =
% \frac{\sum_{i\in S} r_i(f)P_i}{\sum_{i\in S} P_i}
% =
% \frac{\overline{rP}_S}{\bar P_S},
% \]
% where $\overline{rP}_S:=|S|^{-1}\sum_{i\in S} r_i(f)P_i$. Since
% \[
% \overline{rP}_S
% =
% \bar r_S\bar P_S
% +
% \widehat{\operatorname{Cov}}_S(r(f),P),
% \]
% we obtain
% \[
% \operatorname{PRD}(S)
% =
% \frac{\bar r_S}{\overline{r_W}}
% =
% \frac{\bar r_S\bar P_S}
% {\bar r_S\bar P_S+\widehat{\operatorname{Cov}}_S(r(f),P)},
% \]
% which is exactly \eqref{eq:prd_cov_identity}. The sign implication is immediate.

% For part (iii), the PRB response variable $(r_i(f)-m_S)/m_S$ is an affine transform of $r_i(f)$ and therefore remains nonincreasing in $v_i$ whenever $r_i(f)$ is nonincreasing in value and the proxy $v_i$ preserves that ordering. The PRB slope is therefore nonpositive by the same covariance-slope argument used in part (i).

% For part (iv), if the sample is sorted by an ordering proxy that preserves value ranking, then every observation in the highest-value group has ratio no larger than every observation in the lowest-value group under Assumption~\ref{ass:monotone_regressivity}. Hence the median ratio of the highest-value group is no larger than the median ratio of the lowest-value group, which implies $\operatorname{VEI}(S)\le 0$ under the definition in Section~\ref{subsec:equity_metrics}.
% \end{proof}

% Proposition~\ref{prop:metrics_alignment} should be read as a \emph{directional} result: under a monotone ratio-value relationship, the covariance target and the standard vertical-equity metrics all move in the same regressive or non-regressive direction. This is precisely the sense in which covariance acts as a root measure of regressivity in the present framework.

% The next corollary describes the neutral case.

% \begin{corollary}[Flattening to perfect equity]
% \label{cor:neutralization}
% Let $\{f_m\}_{m\ge 1}$ be a sequence of fitted models such that, for each $m$, Assumption~\ref{ass:monotone_regressivity} holds on $S$ and
% \[
% \max_{i\in S}|r_i(f_m)-1|\to 0.
% \]
% Then
% \[
% \Psi_S^{\mathrm{surr}}(f_m)\to 0,
% \qquad
% \widehat{\operatorname{Cov}}_S(r(f_m),z)\to 0,
% \qquad
% \widehat{\operatorname{Cov}}_S(u(f_m),z)\to 0,
% \]
% and the assessor-facing vertical-equity metrics converge to their neutral values:
% \[
% \operatorname{PRD}(S)\to 1,
% \qquad
% \operatorname{PRB}(S)\to 0,
% \qquad
% \operatorname{VEI}(S)\to 0.
% \]
% \end{corollary}

% \begin{proof}
% Uniform convergence $r_i(f_m)\to 1$ on the finite sample $S$ implies $e_i(f_m)=r_i(f_m)-1\to 0$ uniformly, and therefore $\Psi_S^{\mathrm{surr}}(f_m)\to 0$ by \eqref{eq:surrogate_penalty}. The covariances with $z$ also converge to zero because they are continuous finite-sample functionals of $r_i(f_m)$ and $u_i(f_m)=\log r_i(f_m)$.

% For the metrics, $\operatorname{PRD}(S)\to 1$ follows from \eqref{eq:prd_cov_identity}. Likewise, the PRB slope converges to zero because its response variable converges uniformly to zero, and $\operatorname{VEI}(S)\to 0$ because the medians in the top and bottom value groups both converge to $1$.
% \end{proof}

% \begin{remark}[Why the monotone assumption is needed]
% \label{rem:cov_not_sufficient}
% A small covariance by itself does not rule out more complex nonmonotone price effects, because positive and negative local trends can cancel in the aggregate. The role of Assumption~\ref{ass:monotone_regressivity} is precisely to align the covariance target with the assessor metrics. Accordingly, the covariance penalty proposed here should not be interpreted as a replacement for PRD, PRB, or VEI, but rather as a smooth training-time control for the first-order directional component of regressivity, whose reduction moves the standard vertical-equity diagnostics in the desired direction under the monotone root behavior described above.
% \end{remark}




\section{Alternative Pareto Criteria and Convex Stacking}
\label{app:pareto_and_stacking}

This appendix collects the material for the two Pareto-based selection rules of Section~\ref{subsubsec:pareto_alternatives} and for the convex-stacking extension of the constrained calibration problem. Both are deferred from the main text because they are secondary to the constrained rule in \eqref{opt:model_selection_single}, which is the deployment rule we recommend.

\subsection{Pareto-Based Selection Rules}
\label{app:pareto_rules}

Let $\mathcal C$ denote the candidate set, and let $\bar{\mathrm{Perf}}(c)$ and $\{\bar Q_m(c)\}_{m\in\mathcal M}$ denote the aggregated performance and diagnostic scores defined in Section~\ref{subsec:model_selection}. Combining the metrics into a single tuple
$
\bar u(c):=\bigl(\bar{\mathrm{Perf}}(c),\bar Q_1(c),\dots,\bar Q_{|\mathcal M|}(c)\bigr),
$
and applying a monotone rescaling into the unit box $[0,1]^{|\mathcal M|+1}$ in which larger is better on every axis yields a normalized score vector $\tilde u(c)$. The two rules below are defined on these normalized scores.

\paragraph{Closest-to-Utopian.}
Let the Utopian reference be
$
u^\star := \big(\max_{c\in\mathcal C}\tilde u_1(c),\dots,\max_{c\in\mathcal C}\tilde u_{|\mathcal M|+1}(c)\big).
$
The rule selects
$
c^\star_{\text{Uto}} \in \arg\min_{c\in\mathcal C} \lVert u^\star - \tilde u(c)\rVert_2.
$

\paragraph{Nash equilibrium / maximum-volume.}
The rule selects
$
c^\star_{\text{Nash}} \in \arg\max_{c\in\mathcal C} \prod_{k} \tilde u_k(c),
$
i.e., the configuration whose normalized score vector has the largest product across dimensions. Because the criterion is proportional, only the sign of each dimension matters.

Both rules operate on the same aggregated scores as \eqref{opt:model_selection_single} and therefore inherit the same fold structure. They differ only in how points on the empirical Pareto frontier are compared. In our CCAO study, the constrained rule, the closest-to-Utopian rule, and the Nash rule all select penalized surrogate configurations in the same region of the frontier (Table~\ref{tab:test_results}).

\subsection{Convex Stacking Extension}
\label{app:stacking}

A natural extension of \eqref{opt:model_selection_single} replaces discrete candidate selection by a convex combination of candidate predictions. Let $\mathcal C_{\mathrm{stack}}=\{c_1,\dots,c_M\}\subseteq\mathcal C$ denote the stacking pool. For each fold $v$ and observation $i\in V_v$, let $\hat P_{i,j}^{\,v}:=\hat P_i^{\,v}(c_j)$ denote the out-of-fold price-scale prediction produced by configuration $c_j$. For weights $w$ on the simplex $\Delta_M:=\{w\in\R_+^M:\sum_j w_j=1\}$, define the stacked prediction on fold $v$:
\begin{align}
    \label{eq:stacked_prediction}
    \hat P_i^{\,v}(w)
    :=
    \sum_{j=1}^M w_j\,\hat P_{i,j}^{\,v},
    \qquad
    r_i^{\,v}(w)
    :=
    \frac{\hat P_i^{\,v}(w)}{P_i^{\,v}}.
\end{align}
Stacking is defined on the price scale because the assessor diagnostics are ratio-study quantities on that scale. Fold-level performance and diagnostics $\mathrm{Perf}_v(w)$ and $Q_{m,v}(w)$ are computed exactly as in the single-model case, now as functions of $w$. The stacked calibration problem is
\begin{align}
    \label{opt:model_selection_stacking}
    w^\star
    \in
    \arg\max_{w\in\Delta_M}
    \quad
    & \sum_{v\in\mathcal V}\alpha_v\,\mathrm{Perf}_v(w)
    \\
    \text{s.t.}\quad
    & \underline{\kappa}_m
      \le
      \sum_{v\in\mathcal V}\alpha_v\,Q_{m,v}(w)
      \le
      \overline{\kappa}_m,
      \qquad
      m\in\mathcal M_{\mathrm{int}}.
      \nonumber
\end{align}
Once $w^\star$ is identified, the deployed stacked predictor is obtained by refitting each $c_j\in\mathcal C_{\mathrm{stack}}$ on the full pre-assessment development sample and combining its price-scale predictions with weights $w^\star$.

\paragraph{A convex quadratic special case.}
Take the average validation MSE as the objective and worst-fold box constraints on selected ratio-study diagnostics. Let $n_v:=|V_v|$ and define the stacked MSE
$
\mathrm{MSE}_v(w)
:=
n_v^{-1}\sum_{i\in V_v}
\bigl(\hat P_i^{\,v}(w)-P_i^{\,v}\bigr)^2.
$
Consider the problem
\begin{align}
    \label{opt:stacking_convex_example}
    \min_{w\in\Delta_M}
    \quad
    & \frac{1}{V}\sum_{v\in\mathcal V}\mathrm{MSE}_v(w)
    \\
    \text{s.t.}\quad
    & \underline{\tau}_{\mathrm{mean}}
      \le
      \frac{1}{n_v}\sum_{i\in V_v}\frac{\hat P_i^{\,v}(w)}{P_i^{\,v}}
      \le
      \overline{\tau}_{\mathrm{mean}},
      \qquad \forall v\in\mathcal V,
      \nonumber\\
    & \underline{\tau}_{\mathrm{wmean}}
      \le
      \frac{\sum_{i\in V_v}\hat P_i^{\,v}(w)}{\sum_{i\in V_v}P_i^{\,v}}
      \le
      \overline{\tau}_{\mathrm{wmean}},
      \qquad \forall v\in\mathcal V,
      \nonumber\\
    & \underline{\tau}_{\mathrm{PRD}}
      \le
      \frac{
        \frac{1}{n_v}\sum_{i\in V_v}\hat P_i^{\,v}(w)/P_i^{\,v}
      }{
        \sum_{i\in V_v}\hat P_i^{\,v}(w)\big/\sum_{i\in V_v}P_i^{\,v}
      }
      \le
      \overline{\tau}_{\mathrm{PRD}},
      \qquad \forall v\in\mathcal V.
      \nonumber
\end{align}
The objective is convex quadratic in $w$ because each $\hat P_i^{\,v}(w)$ is affine in $w$. The mean-ratio and weighted-mean-ratio constraints are linear in $w$. The PRD box constraints can be rewritten as linear inequalities on each fold:
\begin{align}
    \label{eq:prd_linearized_bounds}
    \frac{1}{n_v}\sum_{i\in V_v}\frac{\hat P_i^{\,v}(w)}{P_i^{\,v}}
    -
    \underline{\tau}_{\mathrm{PRD}}
    \frac{\sum_{i\in V_v}\hat P_i^{\,v}(w)}{\sum_{i\in V_v}P_i^{\,v}}
    &\ge 0,
    \qquad \forall v\in\mathcal V,
    \\
    \frac{1}{n_v}\sum_{i\in V_v}\frac{\hat P_i^{\,v}(w)}{P_i^{\,v}}
    -
    \overline{\tau}_{\mathrm{PRD}}
    \frac{\sum_{i\in V_v}\hat P_i^{\,v}(w)}{\sum_{i\in V_v}P_i^{\,v}}
    &\le 0,
    \qquad \forall v\in\mathcal V,
    \nonumber
\end{align}
so \eqref{opt:stacking_convex_example} is a convex quadratic program over the simplex.

\section{Algorithmic Details for Covariance-Penalized Boosting}
\label{app:boosting_algorithms}

This appendix collects pseudocode for the boosted-tree implementation of the squared covariance penalty introduced in Section~\ref{sec:covariance_penalty_learning}. Throughout, we use the manuscript's sign convention
\[
e_i = f(x_i)-y_i,
\qquad
c_i = y_i-\bar y,
\qquad
z = \widehat{\Cov}_n(e,y)=\frac{1}{n}\sum_{i=1}^n e_i c_i.
\]
Under this convention, the squared covariance penalty is
$
\Psi^{\mathrm{cov}^2}(f)=\frac{1}{2}z^2.
$

\subsection{Exact Newton-style tree boosting with squared covariance penalty}
\label{app:exact_newton_cov_boosting}

The following algorithm implements the exact leaf-wise Newton update discussed in Section~\ref{subsec:tree_boosting}.

\begin{algorithm}[!ht]
\caption{Exact Newton-style tree boosting with squared covariance penalty}
\label{alg:exact_newton_cov_boosting}
\begin{algorithmic}[1]
\Require Training data $\D=\{(x_i,y_i)\}_{i=1}^n$, penalty weight $\rho \ge 0$, leaf regularization $\lambda \ge 0$, tree complexity penalty $\gamma \ge 0$, learning rate $\nu \in (0,1]$, number of rounds $M$
\State Compute $\bar y \gets \frac{1}{n}\sum_{i=1}^n y_i$
\State Compute $c_i \gets y_i-\bar y$ for all $i=1,\dots,n$
\State Initialize $f_0(x) \gets \bar y$
\For{$m=1,\dots,M$}
    \State Compute current predictions $f_i \gets f_{m-1}(x_i)$ for all $i=1,\dots,n$
    \State Compute residuals $e_i \gets f_i-y_i$ for all $i=1,\dots,n$
    \State Compute covariance statistic
    \[
    z \gets \frac{1}{n}\sum_{i=1}^n e_i c_i
    \]
    \State Compute exact gradients
    \[
    g_i \gets \frac{e_i}{n} + \frac{\rho}{n} z c_i
    \qquad \forall i=1,\dots,n
    \]
    \State Search over candidate tree structures $T$
    \For{each candidate tree $T$ with leaves $I_1,\dots,I_J$}
        \For{$j=1,\dots,J$}
            \State $G_j \gets \sum_{i\in I_j} g_i$
            \State $n_j \gets |I_j|$
            \State $s_j \gets \sum_{i\in I_j} c_i$
        \EndFor
        \State Form
        \[
        D \gets \diag\!\left(\frac{n_1}{n}+\lambda,\dots,\frac{n_J}{n}+\lambda\right),
        \qquad
        \beta \gets \frac{\rho}{n^2},
        \qquad
        A_T \gets D + \beta ss^\top
        \]
        \State Compute exact tree gain
        \[
        \operatorname{Gain}(T)
        \gets
        \frac{1}{2}G^\top A_T^{-1}G - \gamma J
        \]
    \EndFor
    \State Select the best tree structure $T_m$
    \State For the selected tree, form $A_m \gets D + \beta ss^\top$
    \State Compute optimal leaf weights
    \[
    w^\star \gets -A_m^{-1}G
    \]
    \State Update
    \[
    f_m(\cdot) \gets f_{m-1}(\cdot) + \nu\,T_m(\cdot;w^\star)
    \]
\EndFor
\State \Return final predictor $f_M$
\end{algorithmic}
\end{algorithm}


% \subsection{Optional first-order boosting variant}
% \label{app:first_order_cov_boosting}

% For completeness, we also record a first-order boosting variant based on fitting base learners to the negative gradient of the penalized objective. This algorithm is mainly pedagogical, and can be omitted if a shorter appendix is preferred.

% \begin{algorithm}[!ht]
% \caption{First-order gradient boosting with squared covariance penalty}
% \label{alg:first_order_cov_boosting}
% \begin{algorithmic}[1]
% \Require Training data $\D=\{(x_i,y_i)\}_{i=1}^n$, penalty weight $\rho \ge 0$, learning rate $\nu \in (0,1]$, number of rounds $M$
% \State Compute $\bar y \gets \frac{1}{n}\sum_{i=1}^n y_i$
% \State Compute $c_i \gets y_i-\bar y$ for all $i=1,\dots,n$
% \State Initialize $f_0(x) \gets \bar y$
% \For{$m=1,\dots,M$}
%     \State Compute current predictions $f_i \gets f_{m-1}(x_i)$ for all $i=1,\dots,n$
%     \State Compute residuals $e_i \gets f_i-y_i$ for all $i=1,\dots,n$
%     \State Compute covariance statistic
%     \[
%     z \gets \frac{1}{n}\sum_{i=1}^n e_i c_i
%     \]
%     \State Form negative-gradient targets
%     \[
%     u_i \gets -\bigl(e_i + \rho z c_i\bigr)
%     \qquad \forall i=1,\dots,n
%     \]
%     \State Fit a base learner structure $T_m$ to $\{(x_i,u_i)\}_{i=1}^n$
%     \If{$T_m$ is a generic weak learner with output vector $h \in \R^n$}
%         \State Compute
%         \[
%         q \gets \frac{1}{n}c^\top h
%         \]
%         \State Set exact line-search step
%         \[
%         \gamma_m
%         \gets
%         -
%         \frac{\frac{1}{n}e^\top h + \rho z q}
%         {\frac{1}{n}\|h\|_2^2 + \rho q^2}
%         \]
%         \State Update
%         \[
%         f_m(\cdot) \gets f_{m-1}(\cdot) + \nu\,\gamma_m h(\cdot)
%         \]
%     \Else
%         \State Let $I_1,\dots,I_J$ be the leaves of $T_m$
%         \For{$j=1,\dots,J$}
%             \State $n_j \gets |I_j|$
%             \State $b_j \gets \sum_{i\in I_j} e_i$
%             \State $s_j \gets \sum_{i\in I_j} c_i$
%             \State $q_j \gets s_j/n$
%         \EndFor
%         \State Form
%         \[
%         D \gets \diag\!\left(\frac{n_1}{n},\dots,\frac{n_J}{n}\right),
%         \qquad
%         q \gets (q_1,\dots,q_J)^\top,
%         \qquad
%         \eta \gets \frac{1}{n}b + \rho z q
%         \]
%         \State Solve the coupled leaf system
%         \[
%         (D+\rho qq^\top)w = -\eta
%         \]
%         \State Update
%         \[
%         f_m(\cdot) \gets f_{m-1}(\cdot) + \nu\,T_m(\cdot;w)
%         \]
%     \EndIf
% \EndFor
% \State \Return final predictor $f_M$
% \end{algorithmic}
% \end{algorithm}

% \begin{table}[!htbp]
% \centering
% \renewcommand{\arraystretch}{1.5} % Extra stretch for stacked cell contents
% \caption{Validation Results: Metrics in bold show improvement over LGBMRegressor. Colored percentages in parentheses under other models indicate percentual change (improvement / worsening) compared to LGBMRegressor.}
% \label{tab:validation_metrics}
% \resizebox{\textwidth}{!}{%
% \begin{tabular}{ll|rr|rr|rrr}
% \toprule
% Model & Config ID & R2 & MAE & COD & COV (IAAO) & VEI & PRD & PRB \\
% \midrule
% LinearRegression & fd63507d & 0.668 & \$92,314 & 32.716 & 0.490 & 8.425\% & 1.026 & 0.059 \\
% LGBMRegressor & 0c59889c & 0.889 & \$58,430 & 23.283 & 0.377 & -19.961\% & 1.066 & -0.069 \\
% \midrule
% SINGLE\_OPT\_REQ\_ROBUST & - &
% \begin{tabular}{@{}c@{}} 0.881 \\ \small\color{worsRed}{(-0.90\%)} \end{tabular} &
% \begin{tabular}{@{}c@{}} \$63,918 \\ \small\color{worsRed}{(-9.39\%)} \end{tabular} &
% \begin{tabular}{@{}c@{}} 24.465 \\ \small\color{worsRed}{(-5.08\%)} \end{tabular} &
% \begin{tabular}{@{}c@{}} \textbf{0.361} \\ \small\color{impGreen}{(+4.24\%)} \end{tabular} &
% \begin{tabular}{@{}c@{}} \textbf{2.538\%} \\ \small\color{impGreen}{(+87.29\%)} \end{tabular} &
% \begin{tabular}{@{}c@{}} \textbf{1.009} \\ \small\color{impGreen}{(+86.36\%)} \end{tabular} &
% \begin{tabular}{@{}c@{}} \textbf{0.026} \\ \small\color{impGreen}{(+62.32\%)} \end{tabular} \\
% SINGLE\_OPT\_UTOPIA\_ROBUST & - &
% \begin{tabular}{@{}c@{}} 0.881 \\ \small\color{worsRed}{(-0.90\%)} \end{tabular} &
% \begin{tabular}{@{}c@{}} \$63,918 \\ \small\color{worsRed}{(-9.39\%)} \end{tabular} &
% \begin{tabular}{@{}c@{}} 24.465 \\ \small\color{worsRed}{(-5.08\%)} \end{tabular} &
% \begin{tabular}{@{}c@{}} \textbf{0.361} \\ \small\color{impGreen}{(+4.24\%)} \end{tabular} &
% \begin{tabular}{@{}c@{}} \textbf{2.538\%} \\ \small\color{impGreen}{(+87.29\%)} \end{tabular} &
% \begin{tabular}{@{}c@{}} \textbf{1.009} \\ \small\color{impGreen}{(+86.36\%)} \end{tabular} &
% \begin{tabular}{@{}c@{}} \textbf{0.026} \\ \small\color{impGreen}{(+62.32\%)} \end{tabular} \\
% \bottomrule
% \end{tabular}%
% }
% \end{table}

% \vspace{3em}

% \begin{table}[!htbp]
% \centering
% \renewcommand{\arraystretch}{1.5} % Extra stretch for stacked cell contents
% \caption{Test Results: Metrics in bold show improvement over LGBMRegressor. Colored percentages in parentheses under other models indicate percentual change (improvement / worsening) compared to LGBMRegressor.}
% \label{tab:test_metrics}
% \resizebox{\textwidth}{!}{%
% \begin{tabular}{ll|rr|rr|rrr}
% \toprule
% Model & Config ID & R2 & MAE & COD & COV (IAAO) & VEI & PRD & PRB \\
% \midrule
% LinearRegression & fd63507d & 0.762 & \$86,244 & 25.013 & 0.369 & -13.527\% & 1.039 & -0.011 \\
% LGBMRegressor & 0c59889c & 0.882 & \$72,672 & 23.437 & 0.359 & -37.257\% & 1.088 & -0.120 \\
% \midrule
% STACKING\_OPT\_AVG\_ROBUST & - &
% \begin{tabular}{@{}c@{}} \textbf{0.883} \\ \small\color{impGreen}{(+0.11\%)} \end{tabular} &
% \begin{tabular}{@{}c@{}} \$73,080 \\ \small\color{worsRed}{(-0.56\%)} \end{tabular} &
% \begin{tabular}{@{}c@{}} \textbf{22.823} \\ \small\color{impGreen}{(+2.62\%)} \end{tabular} &
% \begin{tabular}{@{}c@{}} \textbf{0.335} \\ \small\color{impGreen}{(+6.69\%)} \end{tabular} &
% \begin{tabular}{@{}c@{}} \textbf{-13.100\%} \\ \small\color{impGreen}{(+64.84\%)} \end{tabular} &
% \begin{tabular}{@{}c@{}} \textbf{1.040} \\ \small\color{impGreen}{(+54.55\%)} \end{tabular} &
% \begin{tabular}{@{}c@{}} \textbf{-0.020} \\ \small\color{impGreen}{(+83.33\%)} \end{tabular} \\
% STACKING\_OPT\_UTOPIA\_ROBUST & - &
% \begin{tabular}{@{}c@{}} 0.879 \\ \small\color{worsRed}{(-0.34\%)} \end{tabular} &
% \begin{tabular}{@{}c@{}} \$76,258 \\ \small\color{worsRed}{(-4.94\%)} \end{tabular} &
% \begin{tabular}{@{}c@{}} 23.636 \\ \small\color{worsRed}{(-0.85\%)} \end{tabular} &
% \begin{tabular}{@{}c@{}} \textbf{0.345} \\ \small\color{impGreen}{(+3.90\%)} \end{tabular} &
% \begin{tabular}{@{}c@{}} \textbf{-3.157\%} \\ \small\color{impGreen}{(+91.53\%)} \end{tabular} &
% \begin{tabular}{@{}c@{}} \textbf{1.027} \\ \small\color{impGreen}{(+69.32\%)} \end{tabular} &
% \begin{tabular}{@{}c@{}} \textbf{0.014} \\ \small\color{impGreen}{(+88.33\%)} \end{tabular} \\
% SINGLE\_OPT\_REQ\_ROBUST & - &
% \begin{tabular}{@{}c@{}} 0.878 \\ \small\color{worsRed}{(-0.45\%)} \end{tabular} &
% \begin{tabular}{@{}c@{}} \$77,163 \\ \small\color{worsRed}{(-6.18\%)} \end{tabular} &
% \begin{tabular}{@{}c@{}} 23.929 \\ \small\color{worsRed}{(-2.10\%)} \end{tabular} &
% \begin{tabular}{@{}c@{}} \textbf{0.348} \\ \small\color{impGreen}{(+3.06\%)} \end{tabular} &
% \begin{tabular}{@{}c@{}} \textbf{-2.237\%} \\ \small\color{impGreen}{(+94.00\%)} \end{tabular} &
% \begin{tabular}{@{}c@{}} \textbf{1.025} \\ \small\color{impGreen}{(+71.59\%)} \end{tabular} &
% \begin{tabular}{@{}c@{}} \textbf{0.019} \\ \small\color{impGreen}{(+84.17\%)} \end{tabular} \\
% SINGLE\_OPT\_UTOPIA\_ROBUST & - &
% \begin{tabular}{@{}c@{}} 0.878 \\ \small\color{worsRed}{(-0.45\%)} \end{tabular} &
% \begin{tabular}{@{}c@{}} \$77,163 \\ \small\color{worsRed}{(-6.18\%)} \end{tabular} &
% \begin{tabular}{@{}c@{}} 23.929 \\ \small\color{worsRed}{(-2.10\%)} \end{tabular} &
% \begin{tabular}{@{}c@{}} \textbf{0.348} \\ \small\color{impGreen}{(+3.06\%)} \end{tabular} &
% \begin{tabular}{@{}c@{}} \textbf{-2.237\%} \\ \small\color{impGreen}{(+94.00\%)} \end{tabular} &
% \begin{tabular}{@{}c@{}} \textbf{1.025} \\ \small\color{impGreen}{(+71.59\%)} \end{tabular} &
% \begin{tabular}{@{}c@{}} \textbf{0.019} \\ \small\color{impGreen}{(+84.17\%)} \end{tabular} \\
% \bottomrule
% \end{tabular}%
% }
% \end{table}


% \section{More results}


% \begin{table}[!htbp]
% \centering
% \caption{Held-out test-set performance expressed as percentage change relative to the LGBM baseline. For Med.\ $r$, Mean $r$, PRD, PRB, and VEI, percentages are computed as improvement in distance to the target value rather than raw relative change.}
% \label{tab:test_results_pct_change}
% \resizebox{\textwidth}{!}{%
% \begin{tabular}{llc|cc|cc|cc|ccc}
% \toprule
% Solution & Selected model & $\rho$ & $R^2$ & MAE & Med.\ $r$ & Mean $r$ & COD & COV & PRD & PRB & VEI \\
% \midrule
% \multicolumn{12}{l}{\textbf{Unconstrained baselines}} \\
% \midrule
% BASELINE & Linear Reg. & -- & -13.9\% & +19.9\% & +98.7\% & -142.7\% & +8.4\% & +3.9\% & +54.7\% & +91.1\% & +64.1\% \\
% BASELINE & LGBM & -- & +0.0\% & +0.0\% & +0.0\% & +0.0\% & +0.0\% & +0.0\% & +0.0\% & +0.0\% & +0.0\% \\
% \midrule
% \multicolumn{12}{l}{\textbf{Constraints: PRD}} \\
% \midrule
% CONSTRAINED & LGBSurrogate & 5.96 & -0.4\% & +3.5\% & -13.3\% & -4.9\% & -1.2\% & -5.2\% & +54.0\% & +88.2\% & +75.4\% \\
% UTOPIA & LGBSurrogate & 100.00 & -0.9\% & +10.2\% & -26.0\% & -128.4\% & +6.1\% & -1.7\% & +77.6\% & +73.1\% & +88.3\% \\
% NASH & LGBSurrogate & 2.95 & -0.3\% & +2.1\% & -10.7\% & +41.5\% & -1.9\% & -5.2\% & +43.2\% & +71.0\% & +67.6\% \\
% \midrule
% \multicolumn{12}{l}{\textbf{Constraints: PRD, PRB, and VEI}} \\
% \midrule
% CONSTRAINED & LGBSurrogate & 5.96 & -0.4\% & +3.5\% & -13.3\% & -4.9\% & -1.2\% & -5.2\% & +54.0\% & +88.2\% & +75.4\% \\
% UTOPIA & LGBSurrogate & 7.91 & -0.4\% & +4.2\% & -17.4\% & -25.7\% & -0.2\% & -4.6\% & +57.5\% & +93.9\% & +77.8\% \\
% NASH & LGBSurrogate & 5.96 & -0.4\% & +3.5\% & -13.3\% & -4.9\% & -1.2\% & -5.2\% & +54.0\% & +88.2\% & +75.4\% \\
% \midrule
% \multicolumn{12}{l}{\textbf{Constraints: PRD, Mean $r$, and COD}} \\
% \midrule
% INFEASIBLE* & --- & --- & --- & --- & --- & --- & --- & --- & --- & --- & --- \\
% UTOPIA & LGBSurrogate & 6.87 & -0.3\% & +3.5\% & -13.9\% & -12.3\% & -1.0\% & -5.0\% & +55.1\% & +90.5\% & +75.1\% \\
% NASH & LGBSurrogate & 2.95 & -0.3\% & +2.1\% & -10.7\% & +41.5\% & -1.9\% & -5.2\% & +43.2\% & +71.0\% & +67.6\% \\
% \midrule
% \multicolumn{12}{l}{\textbf{Constraints: Full metric set$^\dagger$}} \\
% \midrule
% INFEASIBLE* & --- & --- & --- & --- & --- & --- & --- & --- & --- & --- & --- \\
% UTOPIA & LGBSurrogate & 5.96 & -0.4\% & +3.5\% & -13.3\% & -4.9\% & -1.2\% & -5.2\% & +54.0\% & +88.2\% & +75.4\% \\
% NASH & LGBSurrogate & 5.96 & -0.4\% & +3.5\% & -13.3\% & -4.9\% & -1.2\% & -5.2\% & +54.0\% & +88.2\% & +75.4\% \\
% \bottomrule
% \multicolumn{12}{l}{\footnotesize *INFEASIBLE abbreviates CONSTRAINED INFEASIBLE; no feasible selected solution.} \\
% \multicolumn{12}{l}{\footnotesize $^\dagger$ PRD, PRB, VEI, COD, COV, Mean $r$, Median $r$, and Weighted Mean $r$.} \\
% \multicolumn{12}{l}{\footnotesize For Med.\ $r$, Mean $r$, and PRD, the target is 1. For PRB and VEI, the target is 0.} \\
% \multicolumn{12}{l}{\footnotesize Positive values indicate movement toward the target (or, for $R^2$, MAE, COD, and COV, improvement relative to LGBM).}
% \end{tabular}}
% \end{table}



\end{document}
